\documentclass[11pt,a4paper]{article}

%================= Packages ===================
\usepackage[utf8]{inputenc}
\usepackage[T1]{fontenc}
\usepackage[french]{babel}
\usepackage{lmodern}
\usepackage{amsmath, amssymb, amsfonts}
\usepackage{mathtools}
\usepackage{amsthm}
\usepackage{geometry}
\usepackage{setspace}
\usepackage{hyperref}
\usepackage{enumitem}
\usepackage{microtype}
\usepackage{tikz}
\usetikzlibrary{arrows.meta,positioning,calc}
\usepackage{url}
\usepackage{graphicx}
\usepackage{booktabs}
\usepackage{tabularx}
\usepackage{longtable}

\geometry{margin=2.3cm}
\setstretch{1.15}
\setlist[itemize]{itemsep=3pt}
\setcounter{secnumdepth}{3}
\setcounter{tocdepth}{3}

\hypersetup{
    colorlinks=true,
    linkcolor=blue,
    citecolor=blue,
    urlcolor=blue
}

%================= Macros =====================
\newcommand{\dd}{\mathrm{d}}
\newcommand{\ee}{\mathrm{e}}
\newcommand{\ii}{\mathrm{i}}
\newcommand{\R}{\mathbb{R}}
\newcommand{\N}{\mathbb{N}}
\newcommand{\E}{\mathbb{E}}
\newcommand{\PP}{\mathbb{P}}
\newcommand{\QQ}{\mathbb{Q}}
\newcommand{\1}{\mathbf{1}}
\DeclareMathOperator*{\argmin}{arg\,min}
\DeclareMathOperator*{\argmax}{arg\,max}

%================= Figures (largeur A4) =======
% Images : ne jamais d\'epasser la largeur du texte (A4 avec marges via geometry)
\makeatletter
\newcommand{\maxwidth}{%
  \ifdim\Gin@nat@width>\linewidth
    \linewidth
  \else
    \Gin@nat@width
  \fi
}
\makeatother
\setkeys{Gin}{width=\maxwidth,keepaspectratio}

% TikZ : encadrer les figures pour tenir dans la largeur du texte.
% On ne redimensionne que si la figure d\'epasse \linewidth (marges A4 via geometry).
\newsavebox{\AFourTikZBox}
\newenvironment{A4TikZ}[1][]
{%
  \begin{center}%
  \begin{lrbox}{\AFourTikZBox}%
  \begin{tikzpicture}[#1]%
}
{%
  \end{tikzpicture}%
  \end{lrbox}%
  \ifdim\wd\AFourTikZBox>\linewidth
    \resizebox{\linewidth}{!}{\usebox{\AFourTikZBox}}%
  \else
    \usebox{\AFourTikZBox}%
  \fi
  \end{center}%
}

%================= Titre ======================
\title{\vspace{-2cm}%
    \textbf{PaperTradingApp : Pricing d'options, mod\'elisation de la volatilit\'e et gestion du risque}\\[0.3cm]
    \large M\'emoire de Master --- synth\`ese th\'eorique et mise en perspective
}
\author{Pr\'enom Nom}
\date{\today}

\begin{document}

%================= Page de garde ==============
\begin{titlepage}
    \centering
    \vspace*{2cm}
    {\Large Programme / \'Ecole}\\[0.7cm]
    {\huge \bfseries PaperTradingApp\\[0.3cm]
    Plateforme de paper trading d'options : synth\`ese th\'eorique}\\[2cm]
    {\large Pr\'enom Nom}\\[0.5cm]
    {\large Ann\'ee universitaire 20XX--20XX}\\[3cm]

    \vfill
    \begin{flushleft}
    \textbf{R\'esum\'e :}\\[0.2cm]
    Ce m\'emoire pr\'esente une synth\`ese structur\'ee des notions th\'eoriques mobilis\'ees
    par une plateforme de \emph{paper trading} d'options. Il couvre :
    \begin{itemize}
      \item les produits d\'eriv\'es (options vanilles et exotiques) et leurs profils de risque ;
      \item le pricing (Black--Scholes, arbres, Monte Carlo, m\'ethodes de Fourier) et les Greeks ;
      \item la volatilit\'e implicite, les surfaces de volatilit\'e et la calibration de mod\`eles
            (Heston, SABR, mod\`eles \`a m\'emoire de type Volterra/rough, diffusions \`a sauts) ;
      \item la construction de courbes de taux (facteurs d'actualisation, forwards, bootstrapping,
            param\'etrisation Nelson--Siegel) ;
      \item la gestion du risque et du portefeuille (PnL, drawdowns, VaR/ES, allocation) ;
      \item l'analyse de r\'egimes de volatilit\'e (vol r\'ealis\'ee, mean reversion, cha\^ines de Markov),
            et une ouverture sur des approches d'apprentissage (surrogates, couverture sous forme MDP).
    \end{itemize}
    L'objectif est que le lecteur puisse constater une ma\^{i}trise explicative de chaque notion,
    ainsi qu'une capacit\'e \`a les articuler dans un cadre coh\'erent d'ing\'enierie quantitative.

    \medskip
    \textbf{Important :} le document est volontairement \emph{strictement th\'eorique} : aucun code,
    aucun pseudo-code, et aucun d\'etail d'impl\'ementation. Lorsque des \'el\'ements empiriques
    d\'ependraient de donn\'ees non fournies, on explicite une \emph{m\'ethode de validation} sans
    avancer de r\'esultat chiffr\'e.
    \end{flushleft}

    \vfill
\end{titlepage}

%================= Table des matières =========
\tableofcontents
\newpage

%================= Notations ==================
\section*{Notations et conventions}
\addcontentsline{toc}{section}{Notations et conventions}

\begin{center}
\begin{tabularx}{\textwidth}{@{}lX@{}}
\toprule
\textbf{Notation} & \textbf{Description (conventions)} \\
\midrule
$\bigl(\Omega,\mathcal{F},(\mathcal{F}_t)_{t\in[0,T]},\PP\bigr)$ & Espace probabilis\'e filtr\'e (information disponible \`a l'instant $t$). \\
$\QQ$ & Mesure risque--neutre (\'equivalente \`a $\PP$) lorsque l'absence d'arbitrage est satisfaite. \\
$B_t$ & Num\'eraire (compte mon\'etaire) ; typiquement $B_t=\exp\!\bigl(\int_0^t r_u\,\dd u\bigr)$. \\
$t\in[0,T]$ & Temps ; $T$ d\'esigne une maturit\'e ; $\Delta t$ un horizon. \\
$S_t$ & Prix spot du sous-jacent ; $X_t=\ln S_t$ est le log-prix. \\
$D_t$ / $q_t$ & Dividendes (processus) / taux de dividende (portage) dans des conventions simplifi\'ees. \\
$r_t$ & Taux court (instantan\'e) ; $r$ d\'esigne un taux constant lorsque suppos\'e. \\
$P(0,T)$ & Prix d'un z\'ero-coupon ; $D(0,T)=P(0,T)$ facteur d'actualisation. \\
$f(0,T)$ & Taux forward instantan\'e : $f(0,T)=-\partial_T \ln P(0,T)$. \\
$K$ & Strike ; $F(t,T)$ un prix forward \`a maturit\'e $T$ (convention selon dividendes/num\'eraire). \\
$V(t,S)$ & Prix d'un d\'eriv\'e ; $C(t,S)$ (call) et $P(t,S)$ (put). Dans la partie taux, $P(0,T)$ d\'esigne un z\'ero-coupon. \\
$\sigma$ & Volatilit\'e (ex. Black--Scholes) ; $\sigma_{\mathrm{imp}}$ vol implicite ; $\widehat\sigma$ vol r\'ealis\'ee. \\
$(W_t)$ & Mouvement brownien standard (sous $\PP$ ou $\QQ$ selon le contexte). \\
$(\Delta,\Gamma,\nu,\Theta,\rho_{\mathrm{IR}})$ & Greeks : delta, gamma, vega, theta et rho (sensibilit\'e au taux). \\
$v_t$ & Variance instantan\'ee ; param\`etres Heston $(\kappa,\theta,\xi,\rho,v_0)$. \\
$(\alpha,\beta,\rho,\nu)$ & Param\`etres SABR (niveau de vol, \'elasticit\'e, corr\'elation, vol de vol). \\
$H$ & Exposant de Hurst (rough volatility : $H<\tfrac12$). \\
$\lambda$ & Intensit\'e de sauts (Merton/Bates) ; loi des amplitudes selon le mod\`ele. \\
$\mu,\Sigma,w$ & Rendements moyens, covariance, et poids d'allocation d'un portefeuille. \\
$L$ & Variable de perte ; $\mathrm{VaR}_\alpha$ et $\mathrm{ES}_\alpha$ au niveau $\alpha$. \\
$(\mathcal{S},\mathcal{A},\mathcal{P},\mathcal{R},\gamma_{\mathrm{RL}})$ & MDP (couverture) : \'etats, actions, transitions $\mathcal{P}(s'|s,a)$, r\'ecompense $\mathcal{R}$ et facteur d'actualisation $\gamma_{\mathrm{RL}}$. \\
$e,\,\mathcal{I}_e$ & Era (bloc temporel/r\'egime) et ensemble d'indices associ\'es ; $\rho_e$ d\'esigne une mesure calcul\'ee sur l'era $e$ (ex. une corr\'elation). \\
\bottomrule
\end{tabularx}
\end{center}

%================= Glossaire ==================
\section*{Glossaire (termes cl\'es)}
\addcontentsline{toc}{section}{Glossaire (termes cl\'es)}
\begin{description}[style=nextline,leftmargin=2.8cm]
  \item[Absence d'arbitrage] Propri\'et\'e d'un march\'e o\`u aucune strat\'egie auto-financ\'ee \`a co\^ut initial nul ne permet un gain certain avec une probabilit\'e strictement positive de gain strict.
  \item[Mesure risque--neutre] Mesure $\QQ$ sous laquelle les prix actualis\'es des actifs (au num\'eraire choisi) sont des martingales ; elle permet la valorisation par esp\'erance conditionnelle.
  \item[Smile / skew] D\'eformation de la volatilit\'e implicite en fonction du strike (et/ou de la maturit\'e), r\'ev\'elatrice d'asym\'etries, de queues et de dynamique de volatilité.
  \item[Surface de vol] Fonction $(K,T)\mapsto \sigma_{\mathrm{imp}}(K,T)$ observ\'ee sur le march\'e ; elle doit satisfaire des contraintes d'absence d'arbitrage statique.
  \item[Calibration] Probl\`eme inverse consistant \`a choisir des param\`etres de mod\`ele de sorte \`a reproduire des prix ou des volatilit\'es implicites de march\'e.
  \item[Risque de mod\`ele] Risque li\'e \`a l'\'ecart entre le mod\`ele choisi et la dynamique r\'eelle ; distinct du risque de march\'e (mouvements des facteurs) et du risque de donn\'ees (qualit\'e/couverture des observations).
  \item[Greeks] Sensibilit\'es locales du prix d'un d\'eriv\'e aux facteurs (spot, vol, taux, temps), outils centraux de couverture et d'attribution de risque.
  \item[MDP (couverture)] Formalisation de la couverture comme d\'ecision s\'equentielle : l'\'etat r\'esume l'information disponible, l'action d\'efinit la r\'eallocation, la r\'ecompense p\'enalise co\^uts et risque.
  \item[Era] Bloc temporel (ou r\'egime) utilis\'e pour structurer les donn\'ees et \'{e}valuer la stabilit\'e d'un signal dans le temps.
  \item[Leakage] Fuite d'information (souvent temporelle) entre entra\^{i}nement et validation, menant \`a une surestimation de performance hors \'echantillon.
  \item[OOF / stacking / ranking] OOF : pr\'edictions hors-fold (sur donn\'ees non vues) ; stacking : combinaison par meta-mod\`ele ; ranking : transformation monotone par rang (CDF empirique) pour stabiliser la distribution des scores.
\end{description}

%================= Figures & Tables ============
\section*{Liste des figures et tables propos\'ees}
\addcontentsline{toc}{section}{Liste des figures et tables propos\'ees}

\subsection*{Figures propos\'ees}
\begin{enumerate}[leftmargin=*,label=F\arabic*.]
  \item Profils de payoff : call/put, straddle, strangle, butterfly, vertical spread, calendar spread.
  \item Bornes et monotonicit\'e : prix d'option en fonction de $S$ et de $K$ (convexit\'e en strike).
  \item Illustration du th\'eor\`eme de Breeden--Litzenberger (densit\'e risque--neutre \`a partir de $\partial_{KK}C$).
  \item Smile/skew : exemple de coupes $(K\mapsto \sigma_{\mathrm{imp}}(K,T))$ et surface $(K,T)$.
  \item Comparaison qualitative : Black--Scholes vs volatilité stochastique vs sauts (impact sur les queues).
  \item Effet des param\`etres Heston $(\kappa,\theta,\xi,\rho,v_0)$ sur le smile (intuition).
  \item SABR : r\^ole de $(\alpha,\beta,\rho,\nu)$ et domaine de validit\'e de l'approximation de Hagan.
  \item Rough volatility : noyau de Volterra et impact de $H$ sur le skew court terme (qualitatif).
  \item Calibration : sch\'ema conceptuel ``donn\'ees $\rightarrow$ surface $\rightarrow$ param\`etres $\rightarrow$ diagnostic''.
  \item Courbe de taux : $P(0,T)$, $f(0,T)$ et d\'ecomposition Nelson--Siegel (niveau/pente/courbure).
  \item Mesures de risque : sch\'ema VaR vs ES sur une distribution de pertes (conceptuel).
  \item Allocation : fronti\`ere efficiente, contributions au risque (risk parity) et eigen-portfolios (PCA).
  \item Couverture : erreur de delta-hedging (discr\'etisation, co\^uts, risque de mod\`ele) (conceptuel).
  \item MDP de couverture : graphe \'etat--action--r\'ecompense et boucle de r\'etroaction.
  \item R\'egimes : matrice de transition d'une cha\^ine de Markov et distribution stationnaire (conceptuel).
  \item Validation temporelle : d\'ecoupage en fen\^etres (rolling/expanding) et distribution d'une m\'etrique par eras (conceptuel).
\end{enumerate}

\subsection*{Tables propos\'ees}
\begin{enumerate}[leftmargin=*,label=T\arabic*.]
  \item Synth\`ese des hypoth\`eses et param\`etres des mod\`eles : BS, Heston, SABR, Merton, Bates, rBergomi, Rough Heston, Kalman (\'etat-espace).
  \item Comparaison des m\'ethodes de pricing : PDE, arbres, Monte Carlo, Fourier/FFT (avantages, limites, complexit\'e).
  \item Greeks : d\'efinitions, interpr\'etations \'economiques et sensibilit\'es dominantes par produit.
  \item Contraintes d'absence d'arbitrage statique sur une surface : butterfly/vertical/calendar et bornes asymptotiques.
  \item Fonctions de perte de calibration : prix vs vol implicite, pond\'erations (vega), r\'egularisation et contraintes.
  \item Courbe de taux : instruments (z\'ero-coupons, swaps), bootstrapping, interpolation (pi\`eges) et diagnostics.
  \item Mesures de risque : VaR/ES, hypoth\`eses, m\'ethodes d'estimation et limites.
  \item Allocation : Markowitz vs risk parity vs eigen-portfolios (hypoth\`eses, robustesse, sensibilit\'e \`a l'estimation).
  \item RL pour la couverture : Q-learning vs DQN vs PPO (id\'ees, stabilit\'e, risques de sur-apprentissage, validation).
  \item Mod\`eles de r\'egimes : d\'efinitions, estimation, interpr\'etation et limites (non-stationnarit\'e).
  \item Pr\'ediction sous non-stationnarit\'e : risques de leakage, validation temporelle, m\'etriques par eras, et stabilisation (ensembles, post-traitements monotones).
\end{enumerate}
\newpage

%================= Introduction générale ======
\section*{Introduction g\'en\'erale}
\addcontentsline{toc}{section}{Introduction g\'en\'erale}

La finance de march\'e moderne repose sur trois piliers principaux : les \emph{produits}
(options vanilles et exotiques, instruments de taux), les \emph{mod\`eles} (processus stochastiques,
mesure risque--neutre, volatilit\'e locale/stochastique/rough), et les \emph{m\'ethodes num\'eriques}
(arbres, PDE, Monte Carlo, transform\'ees de Fourier).

Parall\`element, certaines approches d'\emph{approximation} et de \emph{d\'ecision s\'equentielle}
compl\`etent l'arsenal classique : (i) des \textbf{surrogates} (r\'eseaux) pour acc\'el\'erer des cartes de prix/Greeks
ou des it\'erations de calibration, et (ii) une formulation \textbf{MDP/RL} pour mod\'eliser la couverture
sous contraintes (co\^uts de transaction, risque de mod\`ele, non-stationnarit\'e).

Ce dossier a un double objectif :
\begin{itemize}
  \item exposer de mani\`ere structur\'ee les notions fondamentales mobilis\'ees par le projet ;
  \item relier ces notions \`a un cas d'usage coh\'erent : une plateforme de paper trading d'options.
\end{itemize}

\bigskip

Le fil directeur suit une progression naturelle : absence d'arbitrage $\rightarrow$ produits et payoffs
$\rightarrow$ calcul stochastique $\rightarrow$ m\'ethodes de pricing $\rightarrow$ volatilit\'e implicite et surfaces
$\rightarrow$ mod\`eles de volatilité et calibration $\rightarrow$ courbe de taux $\rightarrow$ gestion du risque et allocation
$\rightarrow$ couverture et d\'ecision s\'equentielle $\rightarrow$ r\'egimes de volatilité $\rightarrow$ m\'ethodologie de validation.

\newpage

%================= Cartographie du projet ======
\section*{Cartographie du projet PaperTradingApp}
\addcontentsline{toc}{section}{Cartographie du projet PaperTradingApp}

\subsection*{Positionnement}
Ce m\'emoire est structur\'e comme une \emph{preuve de compr\'ehension} : pour chaque fonctionnalit\'e majeure de la
plateforme, on identifie les notions th\'eoriques sous-jacentes et on explique comment elles s'articulent.
L'objectif n'est pas de d\'etailler une impl\'ementation, mais de pr\'esenter un cadre conceptuel rigoureux et
coh\'erent.

\subsection*{Correspondance ``modules \texorpdfstring{$\leftrightarrow$}{<->} notions''}
\begin{center}
\begin{tabularx}{\textwidth}{@{}lX@{}}
\toprule
\textbf{Module (onglet)} & \textbf{Notions th\'eoriques mobilis\'ees} \\
\midrule
Options & Payoffs, parit\'es, smile/skew, vol implicite, Greeks, pricing (BS, arbres, MC, Fourier), exotiques. \\
Calibration avanc\'ee & Probl\`eme inverse, fonctions de perte, contraintes, identifiabilit\'e, r\'egularisation,
surfaces de vol, mod\`eles affines (Heston), SABR, mod\`eles \`a m\'emoire (Volterra/rough), diffusions \`a sauts,
surrogates par apprentissage. \\
Yield Curve & Z\'ero-coupons, facteurs d'actualisation, taux forwards, bootstrapping, interpolation/extrapolation,
param\'etrisations (Nelson--Siegel), coh\'erence absence d'arbitrage. \\
Hedging Systems & Couverture delta/gamma/vega, erreurs de couverture, co\^uts de transaction, contr\^ole stochastique,
formulation MDP et apprentissage par renforcement (DQN/PPO) pour la d\'ecision s\'equentielle. \\
Portefeuille \& risque & Valorisation, expositions, PnL r\'ealis\'e/non r\'ealis\'e, attribution, drawdowns,
VaR/ES, stress tests, allocation (Markowitz, risk parity, eigen-portfolios). \\
Trading & Logique d'ex\'ecution (ordres, slippage), strat\'egies simples (DCA/grid), backtesting et mesures de performance. \\
Bots \& Assistant & Strat\'egies de volatilité (straddle, IV crush), vol r\'ealis\'ee, mean reversion, r\'egimes
et transitions Markov, diagnostic textuel et limites de l'assistance IA. \\
\bottomrule
\end{tabularx}
\end{center}

\subsection*{Plan de r\'edaction recommand\'e (niveau ``m\'emoire'')}
\begin{itemize}
  \item Partie 1 : th\'eorie des produits, du pricing et de la volatilit\'e ; puis courbes de taux et gestion du risque.
  \item Partie 2 : apprentissage utile (surrogates pour calibration, r\'egimes, MDP/RL pour couverture).
  \item Partie 3 : synth\`ese par modules (objectif, hypoth\`eses, limites, liens vers les chapitres th\'eoriques).
  \item Conclusion : limites (donn\'ees, hypoth\`eses de mod\`ele, co\^uts, robustesse) et pistes d'am\'elioration.
\end{itemize}
\newpage

%================= Figures TikZ globales ======
\section*{Figures globales}
\addcontentsline{toc}{section}{Figures globales}

\subsection*{Sch\'ema global du m\'emoire}

\begin{A4TikZ}[
    node distance=2.2cm,
    box/.style={rectangle, draw, rounded corners, minimum width=3.2cm, minimum height=1cm, align=center},
    arrow/.style={-Latex, thick}
]
\node[box] (finance) {Partie 1\\Finance de march\'e};
\node[box, right=of finance] (ml) {Partie 2\\D\'ecision\\(MDP/RL)};
\node[box, below=of finance] (projects) {Partie 3\\Synth\`ese\\projet};
\node[box, right=of projects] (concl) {Partie 4\\Conclusion\\\& ouverture};

\draw[arrow] (finance) -- (projects);
\draw[arrow] (ml) -- (projects);
\draw[arrow] (projects) -- (concl);

\end{A4TikZ}

\subsection*{Carte conceptuelle : mod\'elisation, risque et d\'ecision}

\begin{A4TikZ}[
    node distance=1.8cm,
    box/.style={rectangle, draw, rounded corners, minimum width=3.0cm, minimum height=0.8cm, align=center, font=\small},
    arrow/.style={-Latex, thin}
]
\node[box] (sl) {Surrogates\\(approx.)};
\node[box, below=of sl] (ul) {R\'egimes\\(Markov)};
\node[box, right=3.2cm of sl] (dl) {Mod\`eles\\de vol};
\node[box, below=of dl] (rl) {Couverture\\MDP/RL};

\node[box, above left=2.0cm and 0.5cm of sl] (pricing) {Pricing\\(PDE/MC/FFT)};
\node[box, below left=2.0cm and 0.5cm of ul] (risk) {Gestion du\\risque};
\node[box, right=3.5cm of dl] (vol) {Surfaces de\\volatilité};
\node[box, below=2.0cm of rl] (trading) {D\'ecision\\s\'equentielle};

\draw[arrow] (sl) -- (pricing);
\draw[arrow] (ul) -- (risk);
\draw[arrow] (dl) -- (vol);
\draw[arrow] (rl) -- (trading);

\draw[arrow] (sl) -- (dl);
\draw[arrow] (ul) -- (dl);
\draw[arrow] (dl) -- (rl);

\end{A4TikZ}

\newpage

%==========================================================
%                 PARTIE 1 : FINANCE
%==========================================================
\section{Finance de march\'e}

\subsection{Cadre g\'en\'eral et absence d'arbitrage}

Cette sous-partie fixe le cadre math\'ematique de r\'ef\'erence : mod\'elisation probabiliste,
actualisation, absence d'arbitrage et valorisation sous mesure risque--neutre. Il s'agit du socle
sur lequel s'appuient le pricing, la calibration, la mesure des risques et la couverture.

%%%%%%%%%%%%%%%%%%%%%%%%%%%%%%%%%%%%%%%%%%%%%%%%%%%%%%%%%%%%%%
\subsubsection{D\'efinitions : march\'e, information et strat\'egies}

On consid\`ere un espace probabilis\'e filtr\'e $\bigl(\Omega,\mathcal{F},(\mathcal{F}_t)_{t\in[0,T]},\PP\bigr)$,
o\`u $\mathcal{F}_t$ repr\'esente l'information disponible \`a la date $t$.
Un \emph{march\'e} est mod\'elis\'e par un num\'eraire $B_t>0$ (compte mon\'etaire) et un vecteur de prix
d'actifs risqu\'es $(S_t^1,\dots,S_t^d)$, adapt\'es \`a $(\mathcal{F}_t)$.

\paragraph{Actualisation.}
On dit qu'un processus $X_t$ est \emph{actualis\'e} par $B_t$ via $\widetilde X_t:=X_t/B_t$.
L'intuition est que $B_t$ sert d'unit\'e de compte : tout prix est mesur\'e en ``unit\'es de num\'eraire''.

\paragraph{Strat\'egies auto-financ\'ees.}
Une strat\'egie de trading (simplifi\'ee) est un processus pr\'evisible $\phi_t=(\phi_t^1,\dots,\phi_t^d)$
repr\'esentant les quantit\'es d'actifs d\'etenues, \'eventuellement compl\'et\'e par une position en num\'eraire.
La valeur du portefeuille est $V_t=\sum_{i=1}^d \phi_t^i S_t^i + \psi_t B_t$.
La condition \emph{auto-financ\'ee} formalise l'absence d'apports/retraits : toute variation de $V_t$ provient
des variations de prix des actifs d\'etenus.

\paragraph{Arbitrage.}
Intuitivement, un arbitrage est une strat\'egie auto-financ\'ee qui ne co\^ute rien au d\'epart et qui ne peut
pas perdre, tout en ayant une probabilit\'e strictement positive de gain strict. Selon les cadres, on utilise
des notions plus fines (NA, NFLVR) pour traiter les limites de march\'es continus.

%%%%%%%%%%%%%%%%%%%%%%%%%%%%%%%%%%%%%%%%%%%%%%%%%%%%%%%%%%%%%%
\subsubsection{Intuition : num\'eraire, martingales et prix ``justes''}

En l'absence d'arbitrage et sous hypoth\`eses techniques usuelles, il existe une mesure $\QQ$ telle que les
prix \emph{actualis\'es} des actifs risqu\'es soient des martingales :
\[
\widetilde S_t^i=\frac{S_t^i}{B_t}
\quad \text{est une martingale sous } \QQ,\qquad i=1,\dots,d.
\]
Cette propri\'et\'e encode l'id\'ee que, une fois mesur\'es dans la bonne unit\'e (le num\'eraire), les prix n'ont
pas de tendance exploitable ``gratuitement''. Elle n'implique pas que les actifs soient sans risque, mais que
le risque est compens\'e par le choix de la mesure $\QQ$ (changement de mesure).

%%%%%%%%%%%%%%%%%%%%%%%%%%%%%%%%%%%%%%%%%%%%%%%%%%%%%%%%%%%%%%
\subsubsection{R\'esultat cl\'e : th\'eor\`eme fondamental et valorisation risque--neutre}

\paragraph{Th\'eor\`eme fondamental (version informelle).}
Sous des hypoth\`eses d'absence d'arbitrage (de type NFLVR) et de r\'egularit\'e, il existe une mesure $\QQ$
\'equivalente \`a $\PP$ telle que les prix actualis\'es soient des martingales. R\'eciproquement, l'existence d'une
telle mesure exclut l'arbitrage.

\paragraph{Prix d'un payoff.}
Pour un payoff $\Phi$ \`a maturit\'e $T$ (mesurable par rapport \`a $\mathcal{F}_T$), le \emph{prix} \`a la date $t$
admet la repr\'esentation :
\[
V_t = B_t\, \E^{\QQ}\!\left[\frac{\Phi}{B_T}\,\middle|\,\mathcal{F}_t\right].
\]
Dans un march\'e \emph{complet}, la mesure $\QQ$ est unique et le prix est unique. Dans un march\'e \emph{incomplet}
(typiquement d\`es que l'on introduit de la volatilité stochastique non r\'epliquable, des sauts, ou des co\^uts),
il existe plusieurs mesures martingales : la formule reste un principe, mais le choix de $\QQ$ (ou, \'equivalemment,
la sp\'ecification du mod\`ele) devient une hypoth\`ese de mod\'elisation.

%%%%%%%%%%%%%%%%%%%%%%%%%%%%%%%%%%%%%%%%%%%%%%%%%%%%%%%%%%%%%%
\subsubsection{Actualisation, dividendes et forwards : r\^ole des taux}

\paragraph{Dividendes et portage.}
Si le sous-jacent verse un dividende continu au taux $q$ (cas simplifi\'e), l'absence d'arbitrage conduit \`a la relation
de portage pour le forward europ\'een :
\[
F(0,T)=S_0\,\ee^{(r-q)T},
\]
et plus g\'en\'eralement $F(0,T)=S_0\,\exp\!\bigl(\int_0^T (r_u-q_u)\,\dd u\bigr)$ si $r_u,q_u$ sont d\'eterministes.
Lorsque les taux sont stochastiques, le prix d\'epend du num\'eraire choisi (mesure \`a terme, mesure associ\'ee au z\'ero-coupon),
ce qui souligne l'importance d'une courbe de taux coh\'erente et d'une convention d'actualisation explicite.

%%%%%%%%%%%%%%%%%%%%%%%%%%%%%%%%%%%%%%%%%%%%%%%%%%%%%%%%%%%%%%
\subsubsection{Trois familles de risques : march\'e, mod\`ele, donn\'ees}

\paragraph{Risque de march\'e.}
Risque li\'e \`a l'\'evolution des facteurs (spot, volatilité, taux, corr\'elations) : m\^eme avec un mod\`ele parfait,
les PnL fluctuent.

\paragraph{Risque de mod\`ele.}
Risque li\'e aux hypoth\`eses (log-normalit\'e, absence de sauts, markovianit\'e, forme de smile dynamics, etc.).
Il se manifeste notamment par des erreurs de couverture (hedging error) et des instabilit\'es de calibration.

\paragraph{Risque de donn\'ees.}
Risque li\'e \`a la qualit\'e des observations : microstructure, quotes non synchrones, bid--ask, erreurs de
nettoyage, maturit\'es manquantes. Il affecte directement les surfaces, les courbes et la calibration.

%%%%%%%%%%%%%%%%%%%%%%%%%%%%%%%%%%%%%%%%%%%%%%%%%%%%%%%%%%%%%%
\subsubsection{Limites et pr\'ecautions}

Le cadre ``arbitrage-free'' suppose des march\'es frictionless et des strat\'egies admissibles. En pratique,
la couverture est discr\`ete, co\^uteuse, soumise \`a des contraintes (liquidit\'e, marges), et les march\'es peuvent
\^etre incomplets. Une cons\'equence centrale est que le \emph{pricing} est ins\'eparable d'une gouvernance des
hypoth\`eses : choix du num\'eraire, conventions de dividendes, s\'election de mod\`ele et protocole de validation.

%%%%%%%%%%%%%%%%%%%%%%%%%%%%%%%%%%%%%%%%%%%%%%%%%%%%%%%%%%%%%%
\subsubsection{Lien avec le projet \emph{PaperTradingApp}}

Sans d\'ecrire d'impl\'ementation, le projet mobilise ce cadre pour : (i) d\'efinir une valorisation coh\'erente
par actualisation et absence d'arbitrage, (ii) relier prix, surface de volatilité et densit\'e risque--neutre,
(iii) interpr\'eter les \'ecarts entre mod\`eles comme du risque de mod\`ele, (iv) distinguer la variabilit\'e
des PnL due aux facteurs (risque de march\'e) de celle due aux erreurs de donn\'ees et de calibration.

\subsection{Produits d\'eriv\'es - Options}

\subsubsection{Vanilla / Early Exercise}

Les options vanilles constituent la classe fondamentale des produits dérivés d’options.
Leur payoff dépend uniquement du prix du sous-jacent $S_T$ à l’échéance $T$, ce qui permet
une étude rigoureuse via les outils classiques de calcul stochastique, de théorie de l’arbitrage
et d’analyse fonctionnelle.

%------------------------------------------------
\paragraph{Options européennes.}

Une option européenne ne peut être exercée qu’à la date de maturité $T$.

Les payoffs sont donnés par :
\[
\text{Call} : (S_T - K)^+,
\qquad
\text{Put} : (K - S_T)^+,
\]
où $(x)^+ = \max(x,0)$.

Sous la mesure risque--neutre $\mathbb{Q}$ et en l’absence d’arbitrage, les prix sont :
\[
C_0 = e^{-rT}\,\mathbb{E}^{\mathbb{Q}}\big[(S_T-K)^+\big],
\qquad
P_0 = e^{-rT}\,\mathbb{E}^{\mathbb{Q}}\big[(K-S_T)^+\big].
\]

Dans le modèle de mouvement brownien géométrique
\[
dS_t = r S_t\,dt + \sigma S_t\, dW_t,
\]
l’application du lemme d’Itô permet de dériver la PDE de Black--Scholes ainsi que la formule
fermée du call européen.

La parité call--put s’écrit :
\[
C_0 - P_0 = S_0 - K e^{-rT},
\]
relation qui doit être satisfaite dans tout modèle arbitrage--free.

%------------------------------------------------
\paragraph{Options américaines.}

Une option américaine peut être exercée à tout instant $\tau \in [0,T]$.
La valorisation s’écrit comme un problème de temps d’arrêt optimal :
\[
V_0 = \sup_{\tau \in \mathcal{T}}
\mathbb{E}^{\mathbb{Q}}\!\left[e^{-r\tau}\,\Phi(S_\tau)\right],
\]
où $\mathcal{T}$ désigne l’ensemble des temps d’arrêt adaptés.

La frontière d’exercice est caractérisée par l’inégalité :
\[
\Phi(S_t) \;\ge\; V_t.
\]

Cas classiques :
\begin{itemize}
  \item Call américain \emph{sans dividende} : l’exercice anticipé n’est jamais optimal et l’on a
  $V_0 = C_0^{EU}$.
  \item Put américain : un exercice anticipé peut être optimal lorsque $S_t$ franchit une
  frontière critique $S^\*(t)$, solution d’un problème de frontière libre.
\end{itemize}

En pratique, ces produits nécessitent des méthodes numériques (arbres binomiaux ou
trinômiaux, PDE variationnelles, méthodes de type Longstaff--Schwartz).

%------------------------------------------------
\paragraph{Options bermudéennes.}

Les options bermudéennes constituent un intermédiaire entre européennes et américaines :
l’exercice n’est possible qu’à un ensemble discret de dates
$\{t_1,\dots,t_m\} \subset [0,T]$.

La valorisation se fait par récurrence arrière :
\[
V(t_i,S)
=
\max\Big(
\Phi(S),
e^{-r (t_{i+1}-t_i)}
\mathbb{E}\big[ V(t_{i+1},S_{t_{i+1}})\mid S_{t_i}=S \big]
\Big).
\]

On rencontre ce type de produits en taux (swaptions bermudéennes), dans certains produits
structurés ou produits de crédit exotiques.

%------------------------------------------------
\paragraph{Options vanilles sous volatilité stochastique : modèle de Heston.}

Dans le modèle de Heston, sous la mesure risque--neutre, le sous--jacent $S_t$ et la variance
instantanée $v_t$ satisfont :
\[
\begin{cases}
dS_t = r S_t \, dt + \sqrt{v_t}\, S_t\, dW_t^{(1)}, \\[3pt]
dv_t = \kappa(\theta - v_t)\, dt + \xi \sqrt{v_t}\, dW_t^{(2)},
\end{cases}
\qquad
dW_t^{(1)} dW_t^{(2)} = \rho\, dt.
\]

Le prix du call européen admet une formule semi--fermée de la forme
\[
C_0 = S_0 P_1 - K e^{-rT} P_2,
\]
où $P_1$ et $P_2$ s’expriment comme des intégrales de Fourier faisant intervenir
la fonction caractéristique du log--prix.

La dérivation complète de cette formule (changement de mesure, résolution de l’ODE
pour la fonction caractéristique, inversion de Fourier) est fournie en annexe.

%------------------------------------------------
\paragraph{Volatilité rugueuse et modèle rBergomi.}

Le modèle rBergomi repose sur un mouvement brownien fractionnaire de paramètre de Hurst
$H < \tfrac{1}{2}$, conduisant à une dynamique de variance à mémoire longue de type Volterra.
Il est motivé par la g\'eom\'etrie du skew \`a court terme observ\'ee sur les march\'es actions.

Les éléments clés sont :
\begin{itemize}
  \item définition du fractional Brownian motion $B^H$,
  \item représentation de Mandelbrot--Van Ness,
  \item lien avec les SDE de type Volterra et la variance rough.
\end{itemize}

Les d\'etails (construction du fBM, noyaux, choix de sch\'emas) peuvent \^etre ajout\'es en annexe ; on renvoie sinon \`a la litt\'erature
sur la rough volatility pour les d\'erivations compl\`etes.

%------------------------------------------------
\paragraph{Rôle structurel des options vanilles.}

Les options vanilles sont utilisées pour :
\begin{itemize}
  \item \textbf{calibrer} les modèles de volatilité, en ajustant un vecteur de paramètres $\theta$
  pour minimiser l’écart entre volatilités implicites modèle et marché ;
  \item \textbf{extraire les sensibilités} (Greeks) : 
  $\Delta = \frac{\partial V}{\partial S}$,
  $\Gamma = \frac{\partial^2 V}{\partial S^2}$,
  $\text{Vega} = \frac{\partial V}{\partial \sigma}$, etc. ;
  \item \textbf{construire des payoffs exotiques synthétiques} par superposition linéaire
  d’options vanilles (spreads, butterflies, arbitrages statiques).
\end{itemize}


\subsubsection{Path-dependent options}

Les options path-dependent sont des produits dérivés dont le payoff dépend de
l’ensemble de la trajectoire $(S_t)_{0 \le t \le T}$, et non uniquement de la valeur 
du sous-jacent à maturité. Leur valorisation requiert l'utilisation de probabilités 
fonctionnelles, de techniques avancées de Monte Carlo et, dans certains cas, de 
résultats de théorie des extrêmes du mouvement brownien.

%====================================================
\paragraph{Asian options.}

\textbf{Structure.}  
On distingue deux types principaux de moyennes discrètes :
\[
A_T^{\mathrm{arith}} = \frac{1}{n}\sum_{i=1}^n S_{t_i}, 
\qquad
A_T^{\mathrm{geom}} = \left(\prod_{i=1}^n S_{t_i}\right)^{1/n}.
\]

\textbf{Payoff.}
\[
\text{Call asiatique}: (A_T - K)^+,
\qquad
\text{Put asiatique}: (K - A_T)^+.
\]

\textbf{Pricing.}
\begin{itemize}
  \item Cas géométrique : moyenne lognormale $\Rightarrow$ formule fermée.
  \item Cas arithmétique : pas de forme fermée.
\end{itemize}
Méthodes principales :
\begin{itemize}
  \item Monte Carlo (méthode standard),
  \item variance reduction (antithetic, control variate utilisant la moyenne géométrique),
  \item approximation lognormale de Turnbull--Wakeman.
\end{itemize}

%====================================================
\paragraph{Lookback options.}

\textbf{Structure.}
\[
M_T = \max_{0 \le t \le T} S_t,
\qquad
m_T = \min_{0 \le t \le T} S_t.
\]

\textbf{Payoff.}
\begin{itemize}
  \item \textbf{Floating strike} :
  \[
  \text{Call}: (S_T - m_T)^+,
  \qquad
  \text{Put}: (M_T - S_T)^+.
  \]
  \item \textbf{Fixed strike} :
  \[
  \text{Call}: (M_T - K)^+,
  \qquad
  \text{Put}: (K - m_T)^+.
  \]
\end{itemize}

\textbf{Pricing.}
\begin{itemize}
  \item Formules fermées en Black--Scholes via le principe de réflexion.
  \item En général : simulation Monte Carlo indispensable.
  \item Gestion des franchissements : Brownian bridge pour limiter les erreurs.
\end{itemize}

%====================================================
\paragraph{Forward-start options.}

\textbf{Structure.}  
Le strike est fixé à une date future $t_0 < T$ :
\[
K = \alpha S_{t_0}.
\]

\textbf{Payoff.}
\[
\Phi = (S_T - \alpha S_{t_0})^+.
\]

\textbf{Pricing.}
\begin{itemize}
  \item Interprétation comme un call européen conditionnel à $S_{t_0}$.
  \item En Black--Scholes : équivalent à un call sur un forward.
  \item En général : simulation jointe $(S_{t_0}, S_T)$.
\end{itemize}

%====================================================
\paragraph{Cliquet / Ratchet options.}

\textbf{Structure.}  
Options à réinitialisation périodique :
\[
\Phi = \sum_{i=1}^n (S_{t_i} - S_{t_{i-1}})^+.
\]

\textbf{Payoff.}  
Basé sur les increments positifs ou une accumulation de gains ``lock-in''.

\textbf{Pricing.}
\begin{itemize}
  \item Simulation Monte Carlo multidimensionnelle,
  \item corrélations fortes entre increments $\Rightarrow$ variance élevée,
  \item PDE multidimensionnelle impraticable en général.
\end{itemize}


\subsubsection{Barrier options}

Les options barrières sont des produits dérivés dont le payoff dépend du 
franchissement ou non d’un niveau prédéfini appelé \emph{barrière}. Elles sont 
fortement sensibles à la trajectoire du sous--jacent et requièrent la maîtrise 
des temps de passage du mouvement brownien.

%====================================================
\paragraph{Structure.}

On définit le temps d'atteinte de la barrière $B$ par :
\[
\tau_B = \inf \{ t \ge 0 : S_t \in \mathcal{B} \}.
\]

Deux familles principales :

\begin{itemize}
  \item \textbf{Knock-out} : l’option est annulée si $\tau_B \le T$.  
        \\Exemples : up-and-out, down-and-out.

  \item \textbf{Knock-in} : l’option n’existe que si la barrière est atteinte avant $T$.  
        \\Exemples : up-and-in, down-and-in.
\end{itemize}

%====================================================
\paragraph{Payoff.}

Pour un payoff vanille $\Phi(S_T)$ :

\begin{itemize}
  \item \textbf{Knock-out} :
  \[
  \Phi_{\mathrm{KO}} = \Phi(S_T)\, \mathbf{1}_{\{\tau_B > T\}}.
  \]

  \item \textbf{Knock-in} :
  \[
  \Phi_{\mathrm{KI}} = \Phi(S_T)\, \mathbf{1}_{\{\tau_B \le T\}}.
  \]
\end{itemize}

Relation fondamentale de parité :
\[
\Phi_{\mathrm{KI}} + \Phi_{\mathrm{KO}} = \Phi_{\text{vanilla}}.
\]

%====================================================
\paragraph{Pricing.}

\subparagraph{Cas Black--Scholes : formules fermées.}
Sous un GBM, les hitting times sont explicitables par le principe de 
réflexion. Cela conduit à des formules fermées pour l’ensemble des barrières 
up/down, in/out.

Exemple (call down--and--out) :  
\[
C_{\mathrm{DO}} 
= C_{\mathrm{vanilla}}
  - \left( \frac{B}{S_0} \right)^{2\lambda - 2}
    C_{\mathrm{vanilla}}(S_0', K),
\]
avec
\[
\lambda = \frac{r}{\sigma^2} + \frac12,
\qquad
S_0' = \frac{B^2}{S_0}.
\]

\subparagraph{Observation continue vs discrète.}
\begin{itemize}
  \item Observation continue : barrière active en tout temps → hitting times exacts.
  \item Observation discrète : barrière testée à des dates $\{t_i\}$ seulement  
        $\Rightarrow$ 
        \(
        \mathrm{Price}_{\mathrm{disc}} 
        >
        \mathrm{Price}_{\mathrm{cont}}
        \).
\end{itemize}

Correction : utilisation d’un \textbf{Brownian bridge} entre deux pas de temps.

%----------------------------------------------------
\subparagraph{Pricing sous volatilité locale / stochastique / rough.}

\begin{itemize}
  \item \textbf{Local vol} : PDE 2D (spot + temps), solution numérique stable.
  \item \textbf{Heston} : simulation Monte Carlo nécessaire, schéma QE pour $v_t$.
  \item \textbf{Rough volatility} : trajectoires plus irrégulières, hitting times plus fréquents.
\end{itemize}

%----------------------------------------------------
\subparagraph{Monte Carlo.}

Le Monte Carlo naïf ne détecte pas correctement les franchissements.  
Solution : Brownian bridge.

Probabilité conditionnelle de franchissement :
\[
\mathbb{P}(\text{hit})
=
\exp\!\left(
-\frac{2}{\sigma^2 h}
\,
\ln\!\frac{S_t}{B}
\cdot
\ln\!\frac{S_{t+h}}{B}
\right).
\]

%====================================================
\paragraph{Variantes.}

\begin{itemize}
  \item \textbf{Rebate options} :
  \[
  R \cdot \mathbf{1}_{\{\tau_B < T\}}.
  \]
  \item \textbf{Parisian barriers} :
  hitting cumulé durant $\delta$ :
  \[
  \tau^\star = \inf\left\{t : \int_0^t 
     \mathbf{1}_{\{S_u < B\}} \, du \ge \delta \right\}.
  \]
  \item \textbf{Partial barriers} :
  barrière active uniquement sur $[t_a, t_b]$.
\end{itemize}


\subsubsection{Spreads \& Wings}

Les stratégies « spreads » et « wings » sont des combinaisons linéaires d’options
vanilles visant à modifier le profil de payoff et les sensibilités (gamma,
vega, vomma, convexité). Elles reposent sur le principe :
\[
\Phi = \sum_{i} w_i \, \Phi_i,
\]
où chaque $\Phi_i$ est un payoff vanille.

%====================================================
\paragraph{Bull Call Spread.}

\textbf{Structure.}
\[
\text{Long Call}(K_1) + \text{Short Call}(K_2), \qquad K_1 < K_2.
\]

\textbf{Payoff.}
\[
\Phi(S_T) = (S_T-K_1)^+ - (S_T-K_2)^+.
\]

\textbf{Graphique du payoff.}

\begin{A4TikZ}[scale=0.12]
% Axes
\draw[->,thick] (-5,0) -- (220,0) node[right] {$S_T$};
\draw[->,thick] (0,-5) -- (0,40) node[above] {Payoff};

% Parameters
\def\Kone{50}
\def\Ktwo{120}

% Payoff line
\draw[thick,blue]
  (0,0) --
  (\Kone,0) --
  (\Ktwo,0) --
  (200,{ \Ktwo - \Kone });

% Vertical markers
\draw[dashed] (\Kone,0) -- (\Kone,10) node[above] {$K_1$};
\draw[dashed] (\Ktwo,0) -- (\Ktwo,10) node[above] {$K_2$};
\end{A4TikZ}

\textbf{Prix.}
\[
V_0 = C(K_1) - C(K_2).
\]


%====================================================
\paragraph{Bear Put Spread.}

\textbf{Structure.}
\[
\text{Long Put}(K_2) + \text{Short Put}(K_1),
\qquad K_1 < K_2.
\]

\textbf{Payoff.}
\[
\Phi(S_T) = (K_2-S_T)^+ - (K_1-S_T)^+.
\]

\textbf{Graphique du payoff.}

\begin{A4TikZ}[scale=0.12]
% Axes
\draw[->,thick] (-5,0) -- (220,0) node[right] {$S_T$};
\draw[->,thick] (0,-5) -- (0,40) node[above] {Payoff};

% Parameters
\def\Kone{60}
\def\Ktwo{120}

% Payoff line
\draw[thick,blue]
  (0,{ \Ktwo }) --
  (\Kone,{ \Ktwo - \Kone }) --
  (\Ktwo,0) --
  (200,0);

% Vertical markers
\draw[dashed] (\Kone,0) -- (\Kone,10) node[above] {$K_1$};
\draw[dashed] (\Ktwo,0) -- (\Ktwo,10) node[above] {$K_2$};
\end{A4TikZ}

\textbf{Prix.}
\[
V_0 = P(K_2) - P(K_1).
\]


%====================================================
\paragraph{Straddle.}

\textbf{Structure.}
\[
\text{Long Call}(K) + \text{Long Put}(K).
\]

\textbf{Payoff.}
\[
\Phi(S_T) = |S_T - K|.
\]

\textbf{Graphique du payoff.}

\begin{A4TikZ}[scale=0.12]
% Axes
\draw[->,thick] (-5,0) -- (220,0) node[right] {$S_T$};
\draw[->,thick] (0,-5) -- (0,60) node[above] {Payoff};

% Parameters
\def\K{100}

% Payoff line
\draw[thick,blue]
  (0,{ \K }) --
  (\K,0) --
  (200,{ 200 - \K });

% Symmetric line on left
\draw[thick,blue]
  (0,{ \K }) --
  (\K,0);

% Vertical marker
\draw[dashed] (\K,0) -- (\K,10) node[above] {$K$};
\end{A4TikZ}

\textbf{Prix.}
\[
V_0 = C(K) + P(K).
\]


%====================================================
\paragraph{Strangle.}

\textbf{Structure.}
\[
\text{Long Put}(K_1) + \text{Long Call}(K_2),
\qquad K_1 < K_2.
\]

\textbf{Payoff.}
\[
\Phi(S_T) = (K_1-S_T)^+ + (S_T-K_2)^+.
\]

\textbf{Graphique du payoff.}

\begin{A4TikZ}[scale=0.12]
\draw[->,thick] (-5,0) -- (220,0) node[right] {$S_T$};
\draw[->,thick] (0,-5) -- (0,50) node[above] {Payoff};

\def\Kone{70}
\def\Ktwo{130}

% Put wing
\draw[thick,blue]
  (0,{ \Kone }) --
  (\Kone,0);

% Flat part
\draw[thick,blue]
  (\Kone,0) -- (\Ktwo,0);

% Call wing
\draw[thick,blue]
  (\Ktwo,0) --
  (200,{ 200 - \Ktwo });

\draw[dashed] (\Kone,0) -- (\Kone,10) node[above] {$K_1$};
\draw[dashed] (\Ktwo,0) -- (\Ktwo,10) node[above] {$K_2$};
\end{A4TikZ}

\textbf{Prix.}
\[
V_0 = P(K_1) + C(K_2).
\]


%====================================================
\paragraph{Butterfly spread.}

\textbf{Structure.}
\[
\text{Long }C(K_1)
- 2C(K_2)
+ \text{Long }C(K_3),
\qquad K_1<K_2<K_3.
\]

\textbf{Payoff.}
\[
\Phi(S_T)
=(S_T-K_1)^+ - 2(S_T-K_2)^+ + (S_T-K_3)^+.
\]

\textbf{Graphique du payoff.}

\begin{A4TikZ}[scale=0.12]
\draw[->,thick] (-5,0) -- (240,0) node[right] {$S_T$};
\draw[->,thick] (0,-5) -- (0,40) node[above] {Payoff};

\def\Kone{60}
\def\Ktwo{100}
\def\Kthree{140}

\draw[thick,blue]
  (0,0) --
  (\Kone,0) --
  (\Ktwo,{ \Ktwo - \Kone }) --
  (\Kthree,0) --
  (200,0);

\draw[dashed] (\Kone,0) -- (\Kone,10) node[above] {$K_1$};
\draw[dashed] (\Ktwo,0) -- (\Ktwo,10) node[above] {$K_2$};
\draw[dashed] (\Kthree,0) -- (\Kthree,10) node[above] {$K_3$};
\end{A4TikZ}

\textbf{Prix.}
\[
V_0 = C(K_1) - 2C(K_2) + C(K_3).
\]

%====================================================
\paragraph{Condor spread.}

\textbf{Structure.}
\[
C(K_1) - C(K_2) - C(K_3) + C(K_4),
\qquad K_1<K_2<K_3<K_4.
\]

\textbf{Graphique du payoff.}

\begin{A4TikZ}[scale=0.12]
\draw[->,thick] (-5,0) -- (240,0) node[right] {$S_T$};
\draw[->,thick] (0,-5) -- (0,40) node[above] {Payoff};

\def\Kone{50}
\def\Ktwo{80}
\def\Kthree{120}
\def\Kfour{150}

\draw[thick,blue]
  (0,0) --
  (\Kone,0) --
  (\Ktwo,20) --
  (\Kthree,20) --
  (\Kfour,0) --
  (200,0);

\draw[dashed] (\Kone,0) -- (\Kone,10) node[above] {$K_1$};
\draw[dashed] (\Ktwo,0) -- (\Ktwo,10) node[above] {$K_2$};
\draw[dashed] (\Kthree,0) -- (\Kthree,10) node[above] {$K_3$};
\draw[dashed] (\Kfour,0) -- (\Kfour,10) node[above] {$K_4$};
\end{A4TikZ}


\subsubsection{Calendar \& Diagonal Spreads}

Les calendar spreads et diagonal spreads sont des stratégies multi--maturités
qui exploitent la structure par terme de la volatilité implicite ainsi que la
décroissance temporelle de la valeur temps (Theta). Elles reposent sur la
combinaison d’options de même sous--jacent mais de maturités, voire de strikes,
différents.

%====================================================
\paragraph{Calendar Spread.}

\textbf{Structure.}
\[
\text{Long Call}(K, T_2)
\;+\;
\text{Short Call}(K, T_1),
\qquad T_1 < T_2.
\]

\textbf{Payoff.}
À maturité $T_2$, le payoff du calendar spread est celui de l’option longue,
mais le coût initial est réduit par la jambe courte :
\[
\Pi_0 = C(K,T_2) - C(K,T_1).
\]

\textbf{Pricing.}
Sous Black--Scholes :
\[
C(K,T_i) = BS(S_0,K,r,\sigma,T_i), \quad i=1,2.
\]
Le spread dépend de :
\begin{itemize}
  \item la pente de la structure de volatilité,
  \item la position du strike par rapport au forward,
  \item la différence de theta entre $T_1$ et $T_2$.
\end{itemize}

\textbf{Graphique du payoff (profil à maturité).}

\begin{A4TikZ}[scale=0.13]
% Axes
\draw[->,thick] (-5,0) -- (250,0) node[right] {$S_{T_2}$};
\draw[->,thick] (0,-5) -- (0,50) node[above] {Payoff};

% Parameters
\def\K{120}

% Calendar payoff approximé (bosse autour de K)
\draw[thick,blue]
  (0,0) --
  (\K-40, 5) --
  (\K, 30) --
  (\K+40, 5) --
  (240,0);

% Strike marker
\draw[dashed] (\K,0) -- (\K,10) node[above] {$K$};
\end{A4TikZ}

%====================================================
\paragraph{Diagonal Spread.}

\textbf{Structure.}
\[
\text{Long Call}(K_2,T_2)
\;+\;
\text{Short Call}(K_1,T_1),
\qquad T_2 > T_1,\quad K_1 \neq K_2.
\]

\textbf{Payoff.}
\[
\Pi_0 = C(K_2,T_2) - C(K_1,T_1).
\]
Profil final : payoff de call $(S_{T_2}-K_2)^+$, mais avec dynamique
intermédiaire asymétrique en fonction de $K_1$.

\textbf{Pricing.}
\begin{itemize}
  \item Long vega long terme, short vega court terme.
  \item Sensibilité directionnelle via $(K_1,K_2)$.
  \item Interprétation : combinaison directionnelle + term structure.
\end{itemize}

\textbf{Graphique du payoff (profil typique).}

\begin{A4TikZ}[scale=0.13]
% Axes
\draw[->,thick] (-5,0) -- (250,0) node[right] {$S_{T_2}$};
\draw[->,thick] (0,-5) -- (0,50) node[above] {Payoff};

% Parameters
\def\Kone{100}
\def\Ktwo{140}

% Diagonal payoff approximation
\draw[thick,blue]
  (0,0) --
  (\Kone,0) --
  (\Ktwo,20) --
  (240,{240-\Ktwo});

% Strike markers
\draw[dashed] (\Kone,0) -- (\Kone,10) node[above] {$K_1$};
\draw[dashed] (\Ktwo,0) -- (\Ktwo,10) node[above] {$K_2$};
\end{A4TikZ}


\subsubsection{Exotic options}

Les options exotiques sont des dérivés dont la structure de payoff dépasse celle des
options vanilles. Elles permettent de modéliser des expositions complexes : vues
directionnelles asymétriques, paris sur la corrélation, paris sur la volatilité locale ou
stochastique, ou encore élimination du risque de change (produits quanto).

%%%%%%%%%%%%%%%%%%%%%%%%%%%%%%%%%%%%%%%%%%%%%%%%%%%%%%
% 1. Digital Options
%%%%%%%%%%%%%%%%%%%%%%%%%%%%%%%%%%%%%%%%%%%%%%%%%%%%%%
\paragraph{Digital options.}

\textbf{Structure.} 
Une option digitale verse un montant fixe $Q$ si la condition de moneyness est satisfaite.

\textbf{Payoff.}
\[
\Phi_{\mathrm{DC}} = Q \cdot \mathbf{1}_{\{ S_T > K \}}, 
\qquad
\Phi_{\mathrm{DP}} = Q \cdot \mathbf{1}_{\{ S_T < K \}}.
\]

\textbf{Pricing (Black--Scholes).}
\[
V_0^{\mathrm{DC}} = Q e^{-rT} N(d_2),
\qquad
d_2 = \frac{\ln(S_0/K) + (r - \tfrac12 \sigma^2)T}{\sigma \sqrt{T}}.
\]

\textbf{Graphique du payoff (digital call).}

\begin{A4TikZ}[scale=0.13]
\draw[->,thick] (-5,0) -- (220,0) node[right] {$S_T$};
\draw[->,thick] (0,-5) -- (0,40) node[above] {Payoff};

\def\K{100}

\draw[thick,blue]
  (0,0) --
  (\K,0) --
  (\K,30) --
  (200,30);

\draw[dashed] (\K,0) -- (\K,10) node[above] {$K$};
\end{A4TikZ}


%%%%%%%%%%%%%%%%%%%%%%%%%%%%%%%%%%%%%%%%%%%%%%%%%%%%%%
% 2. Asset-or-Nothing Options
%%%%%%%%%%%%%%%%%%%%%%%%%%%%%%%%%%%%%%%%%%%%%%%%%%%%%%
\paragraph{Asset-or-nothing options.}

\textbf{Structure.}  
Le payoff dépend du niveau du sous-jacent et d’une condition binaire.

\textbf{Payoff.}
\[
\Phi_{\mathrm{AON}} = S_T \cdot \mathbf{1}_{\{ S_T > K \}}.
\]

\textbf{Pricing (Black--Scholes).}
\[
V_0 = S_0 N(d_1).
\]

\textbf{Graphique du payoff (asset-or-nothing call).}

\begin{A4TikZ}[scale=0.13]
\draw[->,thick] (-5,0) -- (240,0) node[right] {$S_T$};
\draw[->,thick] (0,-5) -- (0,60) node[above] {Payoff};

\def\K{120}

\draw[thick,blue]
  (0,0) --
  (\K,0) --
  (200,200);

\draw[dashed] (\K,0) -- (\K,10) node[above] {$K$};
\end{A4TikZ}


%%%%%%%%%%%%%%%%%%%%%%%%%%%%%%%%%%%%%%%%%%%%%%%%%%%%%%
% 3. Chooser Options
%%%%%%%%%%%%%%%%%%%%%%%%%%%%%%%%%%%%%%%%%%%%%%%%%%%%%%
\paragraph{Chooser options.}

\textbf{Structure.}
À une date $t_c < T$, l’acheteur choisit entre un call et un put.

\textbf{Payoff.}
\[
\Phi = \max\big( C(t_c,T), \, P(t_c,T) \big).
\]

\textbf{Relation clé.}
\[
\Phi = C(t_c,T) + P(t_c,T) - |S_{t_c} - K|.
\]

\textbf{Pricing (Black--Scholes).}
\[
V_0 = C(S_0,K,T) + P(S_0,K,T) - P(S_0,K,t_c).
\]

\textbf{Graphique du payoff à $t_c$.}

\begin{A4TikZ}[scale=0.14]
\draw[->,thick] (-5,0) -- (220,0) node[right] {$S_{t_c}$};
\draw[->,thick] (0,-5) -- (0,60) node[above] {Payoff};

\def\K{100}

\draw[thick,blue]
  (0,\K) --
  (\K,0) --
  (200,{200 - \K});

\draw[dashed] (\K,0) -- (\K,10) node[above] {$K$};
\end{A4TikZ}


%%%%%%%%%%%%%%%%%%%%%%%%%%%%%%%%%%%%%%%%%%%%%%%%%%%%%%
% 4. Quanto Options
%%%%%%%%%%%%%%%%%%%%%%%%%%%%%%%%%%%%%%%%%%%%%%%%%%%%%%
\paragraph{Quanto options.}

\textbf{Structure.}
Option sur un actif étranger, payoff payé dans une autre devise.

\[
\Phi_{\mathrm{quanto}} = \Phi_{\mathrm{local}}(S_T) \times FX_0.
\]

\textbf{Drift risque--neutre ajusté.}
\[
dS_t = (r_d - q - \rho_{S,X}\sigma_S\sigma_X) S_t \, dt
      + \sigma_S S_t\, dW_t^{(S)}.
\]

\textbf{Graphique du payoff (quanto call).}

\begin{A4TikZ}[scale=0.13]
\draw[->,thick] (-5,0) -- (240,0) node[right] {$S_T$};
\draw[->,thick] (0,-5) -- (0,60) node[above] {Payoff};

\def\K{100}

\draw[thick,blue]
  (0,0) --
  (\K,0) --
  (200,{200 - \K});

\draw[dashed] (\K,0) -- (\K,10) node[above] {$K$};
\end{A4TikZ}


%%%%%%%%%%%%%%%%%%%%%%%%%%%%%%%%%%%%%%%%%%%%%%%%%%%%%%
% 5. Rainbow Options
%%%%%%%%%%%%%%%%%%%%%%%%%%%%%%%%%%%%%%%%%%%%%%%%%%%%%%
\paragraph{Rainbow options.}

\textbf{Structure.}
\[
\Phi_{\max} = (\max_i S_T^{(i)} - K)^+,
\qquad
\Phi_{\min} = (K - \min_i S_T^{(i)})^+.
\]

\textbf{Payoff.}
Best-of / worst-of.

\textbf{Pricing.}
Dépend de la corrélation :
\[
\sigma_B^2 = 
\sum_{i,j} w_i w_j \sigma_i \sigma_j \rho_{ij}.
\]
Méthodes :
\begin{itemize}
  \item simulation Monte Carlo multivariée,
  \item approximations lognormales,
  \item pricing via copules.
\end{itemize}

\textbf{Graphique du payoff (best-of, 2 actifs).}

\begin{A4TikZ}[scale=0.13]
\draw[->,thick] (-5,0) -- (240,0) node[right] {$\max(S_T^1,S_T^2)$};
\draw[->,thick] (0,-5) -- (0,60) node[above] {Payoff};

\def\K{120}

\draw[thick,blue]
  (0,0) --
  (\K,0) --
  (200,{200 - \K});

\draw[dashed] (\K,0) -- (\K,10) node[above] {$K$};
\end{A4TikZ}


\subsubsection{Basket options}

Les basket options sont des dérivés dont le payoff dépend d’un panier pondéré 
de plusieurs sous--jacents. Leur valorisation est profondément influencée par 
la corrélation entre les actifs, ainsi que par la structure de leurs 
volatilités respectives.

Le panier s’écrit :
\[
B_T = \sum_{i=1}^{n} w_i S_T^{(i)},
\qquad 
\sum_{i=1}^n w_i = 1, \; w_i \ge 0.
\]

%%%%%%%%%%%%%%%%%%%%%%%%%%%%%%%%%%%%%%%%%%%%%%%%%%%%%%
% 1. Basket Call / Put
%%%%%%%%%%%%%%%%%%%%%%%%%%%%%%%%%%%%%%%%%%%%%%%%%%%%%%
\paragraph{Basket call / put.}

\textbf{Payoff.}
\[
\Phi_{\mathrm{call}} = (B_T - K)^+,
\qquad
\Phi_{\mathrm{put}}  = (K - B_T)^+.
\]

\textbf{Pricing.}
Même lorsque chaque $S_t^{(i)}$ suit un GBM, le panier $B_T$ n’est pas 
lognormal. Il n’existe donc pas de formule fermée générale.

Approximations usuelles :
\begin{itemize}
  \item \textbf{Approximation lognormale} : $B_T \approx \text{LN}(\mu_B,\sigma_B^2)$ où
  \[
  \sigma_B^2 = \sum_{i,j} w_i w_j \sigma_i \sigma_j \rho_{ij}.
  \]
  \item \textbf{Moment matching} :
  ajustement d’une lognormale équivalente via 
  $\mathbb{E}[B_T]$ et $\mathrm{Var}(B_T)$.
  \item \textbf{Monte Carlo multivarié} :
  simulation corrélée via factorisation de Cholesky ou SVD.
\end{itemize}

\textbf{Graphique (Basket Call).}

\begin{A4TikZ}[scale=0.13]
\draw[->,thick] (-5,0) -- (240,0) node[right] {$B_T$};
\draw[->,thick] (0,-5) -- (0,60) node[above] {Payoff};

\def\K{120}

\draw[thick,blue]
  (0,0) --
  (\K,0) --
  (200,{200-\K});

\draw[dashed] (\K,0) -- (\K,10) node[above] {$K$};
\end{A4TikZ}


%%%%%%%%%%%%%%%%%%%%%%%%%%%%%%%%%%%%%%%%%%%%%%%%%%%%%%
% 2. Basket Put
%%%%%%%%%%%%%%%%%%%%%%%%%%%%%%%%%%%%%%%%%%%%%%%%%%%%%%
\paragraph{Basket put.}

\textbf{Graphique (Basket Put).}

\begin{A4TikZ}[scale=0.13]
\draw[->,thick] (-5,0) -- (240,0) node[right] {$B_T$};
\draw[->,thick] (0,-5) -- (0,60) node[above] {Payoff};

\def\K{120}

\draw[thick,blue]
  (0,\K) --
  (\K,0) --
  (200,0);

\draw[dashed] (\K,0) -- (\K,10) node[above] {$K$};
\end{A4TikZ}


%%%%%%%%%%%%%%%%%%%%%%%%%%%%%%%%%%%%%%%%%%%%%%%%%%%%%%
% 3. Best-of / Worst-of
%%%%%%%%%%%%%%%%%%%%%%%%%%%%%%%%%%%%%%%%%%%%%%%%%%%%%%
\paragraph{Best-of / Worst-of options.}

\textbf{Structure.}
\[
\Phi_{\max} = (\max_i S_T^{(i)} - K)^+,
\qquad
\Phi_{\min} = (K - \min_i S_T^{(i)})^+.
\]

\textbf{Sensibilités.}
La dépendance à la corrélation est immédiate :
\begin{itemize}
  \item Best-of : augmente lorsque la corrélation baisse.
  \item Worst-of : augmente lorsque la corrélation grimpe.
\end{itemize}

\textbf{Pricing.}
\begin{itemize}
  \item Monte Carlo multivarié (méthode standard).
  \item Approximations lognormales / copules gaussiennes.
  \item Pricing Fourier multivarié (méthodes COS).
\end{itemize}

\textbf{Graphique (Best-of, 2 actifs).}

\begin{A4TikZ}[scale=0.13]
\draw[->,thick] (-5,0) -- (240,0) node[right] {$\max(S_T^1,S_T^2)$};
\draw[->,thick] (0,-5) -- (0,60) node[above] {Payoff};

\def\K{120}

\draw[thick,blue]
  (0,0) --
  (\K,0) --
  (200,{200-\K});

\draw[dashed] (\K,0) -- (\K,10) node[above] {$K$};
\end{A4TikZ}



\subsubsection{Lien avec le projet \emph{PaperTradingApp}}

Cette partie correspond au socle ``produits'' : elle sert \`a (i) d\'efinir les payoffs (vanilles et exotiques) manipul\'es
par la plateforme, (ii) lire des profils de risque (lin\'earit\'e, convexit\'e, exposition aux queues), et (iii) analyser des
strat\'egies simples (spreads, wings, straddles) en termes de sensibilit\'es et de sc\'enarios. Le lien avec le projet n'est pas
une impl\'ementation, mais une exigence de coh\'erence : une strat\'egie n'a de sens que rapport\'ee \`a un payoff et \`a des bornes
absence d'arbitrage (parit\'es, monotonicit\'es, convexit\'e), puis \`a un mod\`ele de dynamique (volatilité, sauts) discut\'e plus loin.

\subsection{Pricing des options}

\subsubsection{Variation quadratique des processus}

\bigskip

\paragraph{Définition.}

Soit une partition d’un intervalle $[0,T]$ :
\[
\Pi_n = \{0 = t_0^{(n)} < t_1^{(n)} < \dots < t_{k_n}^{(n)} = T\},
\qquad
|\Pi_n| := \max_i (t^{(n)}_{i+1} - t^{(n)}_i).
\]

On définit la \emph{variation quadratique} d’un processus continu $(X_t)$ par la limite,
lorsqu’elle existe :
\[
[X]_T
:= \lim_{n\to\infty}
\sum_{i=0}^{k_n - 1}
\big( X_{t^{(n)}_{i+1}} - X_{t^{(n)}_{i}} \big)^2,
\qquad
|\Pi_n| \to 0.
\]

Lorsque la limite existe pour tout $t\le T$, on obtient un processus croissant 
$t\mapsto [X]_t$.

%%%%%%%%%%%%%%%%%%%%%%%%%%%%%%%%%%%%%%%%%%%%%%%%%%%%%%%%%%%%%%
\paragraph{Variation quadratique du mouvement brownien.}

Soit $(W_t)$ un mouvement brownien standard.  
Sur un intervalle $[t_i,t_{i+1}]$ de longueur $\Delta t_i$ :
\[
W_{t_{i+1}} - W_{t_i} \sim \mathcal{N}(0,\Delta t_i)
\quad\Rightarrow\quad
\mathbb{E}[(W_{t_{i+1}} - W_{t_i})^2] = \Delta t_i.
\]

Considérons :
\[
S_n := \sum_i (W_{t^{(n)}_{i+1}} - W_{t^{(n)}_{i}})^2.
\]

Alors :
\[
\mathbb{E}[S_n] = \sum_i \Delta t_i = T.
\]

Pour la variance :
\[
\mathrm{Var}(S_n)
= \sum_i \mathrm{Var}\big( (W_{t_{i+1}} - W_{t_i})^2 \big).
\]

Si $Z \sim \mathcal{N}(0,\Delta t)$, alors
\[
\mathbb{E}[Z^4] = 3 (\Delta t)^2, 
\qquad
\mathrm{Var}(Z^2) = 3(\Delta t)^2 - (\Delta t)^2 = 2(\Delta t)^2.
\]

Donc :
\[
\mathrm{Var}(S_n) 
= 2 \sum_i (\Delta t_i)^2.
\]

Or, lorsque $|\Pi_n|\to 0$ :
\[
\sum_i (\Delta t_i)^2 
\le |\Pi_n|\sum_i \Delta t_i 
= |\Pi_n|\,T 
\xrightarrow[n\to\infty]{} 0.
\]

Ainsi :
\[
\mathrm{Var}(S_n) \longrightarrow 0.
\]

Par l’inégalité de Chebyshev :
\[
\forall \varepsilon>0,\qquad
\mathbb{P}\big(|S_n - T| > \varepsilon\big)
\le \frac{\mathrm{Var}(S_n)}{\varepsilon^2}
\xrightarrow[n\to\infty]{} 0.
\]

Donc :
\[
S_n \xrightarrow[n\to\infty]{\mathbb{P}} T,
\]
et, en renforçant l’argument (partitions dyadiques + Borel--Cantelli),
\[
[W]_t = t
\qquad \text{presque sûrement}.
\]

%%%%%%%%%%%%%%%%%%%%%%%%%%%%%%%%%%%%%%%%%%%%%%%%%%%%%%%%%%%%%%
\paragraph{Variation quadratique d’un processus d’It\^{o}.}

Considérons une diffusion :
\[
X_t = X_0 + \int_0^t a_s\, ds + \int_0^t b_s\, dW_s.
\]

Sur un petit intervalle $[t_i,t_{i+1}]$ :
\[
\Delta X_i
=
X_{t_{i+1}} - X_{t_i}
=
\underbrace{\int_{t_i}^{t_{i+1}} a_s\, ds}\_{\mathcal{O}(\Delta t_i)}
+
\underbrace{\int_{t_i}^{t_{i+1}} b_s\, dW_s}\_{\mathcal{O}(\sqrt{\Delta t_i})}.
\]

Le terme stochastique domine dans le carré :
\[
(\Delta X_i)^2
\approx
\left(\int_{t_i}^{t_{i+1}} b_s\, dW_s\right)^2.
\]

Par isométrie d’It\^{o} :
\[
\mathbb{E}\!\left[
\left(\int_{t_i}^{t_{i+1}} b_s\, dW_s\right)^2
\right]
=
\mathbb{E}\!\left[
\int_{t_i}^{t_{i+1}} b_s^2\, ds
\right]
\approx b_{t_i}^2 \Delta t_i.
\]

Ainsi :
\[
\sum_i (\Delta X_i)^2 
\to \int_0^T b_s^2\, ds.
\]

On obtient rigoureusement :
\[
[X]_t = \int_0^t b_s^2\, ds.
\]

%%%%%%%%%%%%%%%%%%%%%%%%%%%%%%%%%%%%%%%%%%%%%%%%%%%%%%%%%%%%%%
\paragraph{Résumé.}
\begin{itemize}
  \item Pour le mouvement brownien : $[W]_t = t$.
  \item Pour une diffusion d’It\^{o} :
  \[
  dX_t = a_t\,dt + b_t\,dW_t
  \quad\Rightarrow\quad
  [X]_t = \int_0^t b_s^2\, ds.
  \]
  \item La variation quadratique justifie les règles :
  $(dW_t)^2 = dt$, $dt\,dW_t = 0$, $dt^2 = 0$.
\end{itemize}


\subsubsection{Lemme d'It\^{o} (d\'erivation ``for dummies'' rigoureuse)}

\bigskip

Soit un mouvement brownien $(W_t)_{t\ge 0}$ sur un espace probabilis\'e 
filtr\'e $(\Omega,\mathcal{F},(\mathcal{F}_t)_{t\ge 0},\mathbb{P})$, et soit
un processus d'It\^{o} scalaire $(X_t)_{t\in[0,T]}$ de la forme :
\[
dX_t = a(t,X_t)\,dt + b(t,X_t)\,dW_t,
\]
o\`u $a$ et $b$ sont des fonctions mesurables, suffisamment r\'eguli\`eres, et
adapt\'ees. On se donne une fonction
\[
f : [0,T]\times \mathbb{R} \to \mathbb{R},
\qquad f \in C^{1,2},
\]
c'est-\`a-dire $f$ est $C^1$ en $t$ et $C^2$ en $x$ avec toutes ces d\'eriv\'ees
continues.

On d\'efinit le processus
\[
Y_t := f(t,X_t).
\]
Le lemme d'It\^{o} donne l'expression diff\'erentielle de $(Y_t)$.

%%%%%%%%%%%%%%%%%%%%%%%%%%%%%%%%%%%%%%%%%%%%%%%%%%%%%%%%%%%%%%
\paragraph{Heuristique (id\'ee de base).}

Consid\'erons un petit incr\'ement de temps $h>0$ et notons $X_t = x$ pour fixer les id\'ees.
Sur l'intervalle $[t,t+h]$, on a formellement :
\[
X_{t+h} - X_t 
\simeq a(t,x)\,h + b(t,x)\,\Delta W_t,
\quad
\Delta W_t := W_{t+h} - W_t,
\]
avec $\Delta W_t \sim \mathcal{N}(0,h)$. En ordre de grandeur :
\[
a(t,x)\,h = \mathcal{O}(h), 
\qquad
b(t,x)\,\Delta W_t = \mathcal{O}(\sqrt{h}).
\]

On applique un d\'eveloppement de Taylor \`a deux variables \`a $f$ autour de $(t,x)$ :
\[
\begin{aligned}
f(t+h,X_{t+h}) 
&\approx f(t,x) 
+ f_t(t,x)\,h 
+ f_x(t,x)\,(X_{t+h}-x) \\
&\quad + \frac12 f_{xx}(t,x)\,(X_{t+h}-x)^2.
\end{aligned}
\]
En posant $\Delta X := X_{t+h}-X_t$, on obtient donc
\[
\Delta f := f(t+h,X_{t+h}) - f(t,X_t)
\approx
f_t h + f_x\,\Delta X + \tfrac12 f_{xx} (\Delta X)^2.
\]

%%%%%%%%%%%%%%%%%%%%%%%%%%%%%%%%%%%%%%%%%%%%%%%%%%%%%%%%%%%%%%
\paragraph{Comportement du terme $(\Delta X)^2$ et r\^ole de la variation quadratique.}

Avec $\Delta X \simeq a h + b \Delta W_t$, on obtient
\[
(\Delta X)^2 \simeq (a h + b \Delta W_t)^2
= a^2 h^2 + 2ab h \Delta W_t + b^2 (\Delta W_t)^2.
\]

En termes d'ordre :
\[
a^2 h^2 = \mathcal{O}(h^2), 
\qquad
2ab h\Delta W_t = \mathcal{O}(h^{3/2}),
\qquad
b^2 (\Delta W_t)^2 = \mathcal{O}(h).
\]

Les deux premiers termes sont n\'egligeables face \`a $h$ lorsque $h\to 0$,
tandis que le dernier terme est d'ordre $h$ et donc \emph{non n\'egligeable}.
Or la variation quadratique du brownien dit précisément que
\[
(\Delta W_t)^2 \approx h
\]
(en un sens probabiliste ou presque s\^ur via la d\'efinition de $[W]_t$).
On retient donc heuristiquement :
\[
(\Delta X)^2 \approx b(t,x)^2\,h.
\]

%%%%%%%%%%%%%%%%%%%%%%%%%%%%%%%%%%%%%%%%%%%%%%%%%%%%%%%%%%%%%%
\paragraph{Substitution dans le d\'eveloppement de Taylor.}

En rempla\c{c}ant $\Delta X$ et $(\Delta X)^2$ dans l'expression de $\Delta f$ :
\[
\Delta f 
\approx f_t h + f_x (a h + b\Delta W_t)
+ \tfrac12 f_{xx} b^2 h.
\]

On peut r\'earranger les termes :
\[
\Delta f
\approx 
\bigl( f_t + a f_x + \tfrac12 b^2 f_{xx} \bigr)\, h
+ b f_x \,\Delta W_t.
\]

Lorsque $h \to 0$, on interpr\`ete :
\[
\Delta f \sim df,
\qquad h \sim dt,
\qquad \Delta W_t \sim dW_t,
\]
et on obtient formellement la forme diff\'erentielle :

\[
df(t,X_t)
= \bigl( f_t + a f_x + \tfrac12 b^2 f_{xx} \bigr)(t,X_t)\, dt
+ b(t,X_t) f_x(t,X_t)\, dW_t.
\]

%%%%%%%%%%%%%%%%%%%%%%%%%%%%%%%%%%%%%%%%%%%%%%%%%%%%%%%%%%%%%%
\paragraph{Forme rigoureuse via la variation quadratique.}

De mani\`ere plus rigoureuse, on \'ecrit d'abord la version g\'en\'erale :
\[
df(t,X_t)
= f_t(t,X_t)\, dt
+ f_x(t,X_t)\, dX_t
+ \frac12 f_{xx}(t,X_t)\, d[X]_t,
\]
o\`u $[X]_t$ est la variation quadratique de $X$.

Avec 
\[
dX_t = a(t,X_t)\,dt + b(t,X_t)\,dW_t,
\]
on sait que
\[
[X]_t = \int_0^t b(s,X_s)^2\, ds
\quad\Rightarrow\quad
d[X]_t = b(t,X_t)^2\, dt.
\]

En substituant dans la formule ci-dessus, on obtient :
\[
df(t,X_t)
= f_t\,dt + f_x(a\,dt + b\,dW_t) + \tfrac12 f_{xx} b^2 dt,
\]
c'est-\`a-dire
\[
df(t,X_t)
= \bigl( f_t + a f_x + \tfrac12 b^2 f_{xx} \bigr)(t,X_t)\, dt
+ b(t,X_t) f_x(t,X_t)\, dW_t.
\]

%%%%%%%%%%%%%%%%%%%%%%%%%%%%%%%%%%%%%%%%%%%%%%%%%%%%%%%%%%%%%%
\paragraph{Forme int\'egrale.}

En int\'egrant de $0$ \`a $t$ :
\[
\begin{aligned}
f(t,X_t)
&= f(0,X_0)
+ \int_0^t f_t(s,X_s)\,ds
+ \int_0^t f_x(s,X_s)\,dX_s
+ \frac12 \int_0^t f_{xx}(s,X_s)\, d[X]_s \\
&= f(0,X_0)
+ \int_0^t \bigl(f_t + a f_x + \tfrac12 b^2 f_{xx}\bigr)(s,X_s)\,ds
+ \int_0^t b(s,X_s) f_x(s,X_s)\, dW_s.
\end{aligned}
\]

%%%%%%%%%%%%%%%%%%%%%%%%%%%%%%%%%%%%%%%%%%%%%%%%%%%%%%%%%%%%%%
\paragraph{R\'esum\'e.}

Pour un processus d’It\^{o}
\[
dX_t = a(t,X_t)\,dt + b(t,X_t)\,dW_t
\]
et une fonction $f \in C^{1,2}$, le lemme d’It\^{o} affirme que :
\[
df(t,X_t)
= \bigl( f_t + a f_x + \tfrac12 b^2 f_{xx} \bigr)(t,X_t)\, dt
+ b(t,X_t) f_x(t,X_t)\, dW_t.
\]

La pr\'esence du terme $\tfrac12 b^2 f_{xx}$ provient directement de la 
variation quadratique du brownien, via le fait que $(dW_t)^2$ se comporte 
comme $dt$ dans le calcul stochastique.


\subsubsection{Application du lemme d'It\^{o} \`a $\log S_t$}

On consid\`ere le processus de prix suivant sous une mesure $\mathbb{P}$ quelconque :
\[
dS_t = \mu S_t\, dt + \sigma S_t\, dW_t,
\qquad S_0 > 0.
\]

On souhaite d\'eterminer la dynamique de $X_t := \log S_t$ par le lemme d'It\^{o}.

On pose
\[
f(x) = \log x,
\qquad f'(x) = \frac{1}{x},
\qquad f''(x) = -\frac{1}{x^2}.
\]

En appliquant le lemme d'It\^{o} \`a $f(S_t)$, avec $a(t,S_t)=\mu S_t$ et $b(t,S_t)=\sigma S_t$ :
\[
df(S_t)
=
f'(S_t)\, dS_t 
+ \frac12 f''(S_t)\, d[S]_t.
\]

Or la variation quadratique de $S_t$ est
\[
d[S]_t = (\sigma S_t)^2 dt = \sigma^2 S_t^2 dt.
\]

Donc
\[
\begin{aligned}
d(\log S_t)
&= \frac{1}{S_t} (\mu S_t\, dt + \sigma S_t\, dW_t)
  + \frac12 \left(-\frac{1}{S_t^2}\right) \sigma^2 S_t^2 dt \\
&= \left(\mu - \frac12 \sigma^2\right) dt + \sigma\, dW_t.
\end{aligned}
\]

Sous la mesure risque--neutre $\mathbb{Q}$ (o\`u le drift devient $r$), on obtient :
\[
d(\log S_t) = \left(r - \frac12 \sigma^2\right) dt + \sigma\, dW_t^{\mathbb{Q}}.
\]

Cette forme est cruciale pour :
\begin{itemize}
  \item expliciter la loi de $S_T$ (lognormale),
  \item calculer les esp\'erances risque--neutres dans la formule de Black--Scholes,
  \item effectuer des changements de variables (par exemple dans les mod\`eles de vol locale).
\end{itemize}

%%%%%%%%%%%%%%%%%%%%%%%%%%%%%%%%%%%%%%%%%%%%%%%%%%%%%%
\subsubsection{D\'erivation de la PDE de Black--Scholes}

On se place sous la mesure risque--neutre $\mathbb{Q}$, o\`u le sous--jacent suit :
\[
dS_t = r S_t\, dt + \sigma S_t\, dW_t^{\mathbb{Q}},
\]
avec $r$ le taux sans risque constant.

Soit $V(t,S)$ le prix \`a la date $t$ d'un payoff europ\'een $\Phi(S_T)$.
On suppose $V \in C^{1,2}$ et on applique le lemme d'It\^{o} \`a $V(t,S_t)$.

On a :
\[
dV(t,S_t) = V_t\, dt + V_S\, dS_t 
+ \frac12 V_{SS}\, d[S]_t.
\]

En substituant la dynamique de $S_t$ :
\[
dS_t = r S_t\,dt + \sigma S_t\, dW_t,
\qquad
d[S]_t = (\sigma S_t)^2 dt,
\]
on obtient :
\[
\begin{aligned}
dV(t,S_t)
&= V_t\, dt
 + V_S (r S_t\,dt + \sigma S_t\, dW_t)
 + \frac12 V_{SS} \sigma^2 S_t^2 dt \\
&= \bigl( V_t + r S_t V_S + \tfrac12 \sigma^2 S_t^2 V_{SS} \bigr)\, dt
 + \sigma S_t V_S\, dW_t.
\end{aligned}
\]

On construit maintenant un portefeuille auto--finan\c{c}ant 
\[
\Pi_t = V(t,S_t) - \Delta_t S_t,
\]
o\`u $\Delta_t$ est une quantit\'e d'unit\'es du sous--jacent.  
On choisit
\[
\Delta_t := V_S(t,S_t),
\]
de sorte que le terme en $dW_t$ disparaisse dans $d\Pi_t$.

En effet,
\[
d\Pi_t = dV_t - \Delta_t\, dS_t.
\]

On a :
\[
dS_t = r S_t\,dt + \sigma S_t\, dW_t,
\]
et
\[
dV_t = 
\bigl( V_t + r S_t V_S + \tfrac12 \sigma^2 S_t^2 V_{SS} \bigr)\, dt
+ \sigma S_t V_S\, dW_t.
\]

Donc :
\[
\begin{aligned}
d\Pi_t
&= \left[
V_t + r S_t V_S + \tfrac12 \sigma^2 S_t^2 V_{SS}
\right] dt
+ \sigma S_t V_S\, dW_t 
- V_S (r S_t\,dt + \sigma S_t\, dW_t) \\
&= \left[
V_t + r S_t V_S + \tfrac12 \sigma^2 S_t^2 V_{SS} 
- r S_t V_S
\right] dt \\
&= \left[
V_t + \tfrac12 \sigma^2 S_t^2 V_{SS}
\right] dt.
\end{aligned}
\]

Ainsi, le portefeuille $\Pi_t$ est \emph{sûr} (plus de terme stochastique).  
Par absence d’arbitrage, un portefeuille sans risque doit cro\^{i}tre au taux $r$ :

\[
d\Pi_t = r \Pi_t\,dt = r\bigl(V(t,S_t) - V_S(t,S_t) S_t\bigr)\,dt.
\]

En identifiant les coefficients de $dt$ :

\[
V_t + \tfrac12 \sigma^2 S_t^2 V_{SS}
= r\bigl(V - S_t V_S\bigr).
\]

Autrement dit :
\[
V_t + \frac12 \sigma^2 S^2 V_{SS} + r S V_S - r V = 0.
\]

C’est la \textbf{PDE de Black--Scholes} pour tout $(t,S)$ dans $[0,T)\times(0,+\infty)$,
avec condition terminale :
\[
V(T,S) = \Phi(S).
\]

Cette PDE, coupl\'ee \`a la condition de terminale et aux conditions de croissance
appropri\'ees, caract\'erise de fa\c{c}on unique le prix d’option europ\'een
dans le cadre Black--Scholes.


%%%%%%%%%%%%%%%%%%%%%%%%%%%%%%%%%%%%%%%%%%%%%%%%%%%%%%%%%%%%%%
\subsubsection{Application : r\'esolution explicite de la PDE de Black--Scholes}

On consid\`ere la PDE de Black--Scholes pour le prix $V(t,S)$ d’un call europ\'een
de strike $K$ et maturit\'e $T$ :
\[
V_t + \frac12 \sigma^2 S^2 V_{SS} + r S V_S - r V = 0,
\qquad 0 \le t < T,\ S>0,
\]
avec condition terminale :
\[
V(T,S) = (S-K)^+.
\]

L’objectif est de retrouver la formule ferm\'ee de Black--Scholes \`a partir de cette PDE.

\paragraph{Changement de variables.}

On introduit les variables :
\[
\tau := T - t \quad\text{(temps \`a maturit\'e)}, 
\qquad
x := \ln S.
\]

On note $u(\tau,x) := V(t,S) = V(T-\tau,e^x)$.
On calcule les d\'eriv\'ees :

\[
V_t = -u_\tau,
\qquad
V_S = \frac{1}{S} u_x,
\qquad
V_{SS} = \frac{1}{S^2}(u_{xx} - u_x).
\]

En substituant dans la PDE de Black--Scholes :

\[
-u_\tau + \frac12 \sigma^2 S^2 \frac{1}{S^2}(u_{xx} - u_x)
+ r S \frac{1}{S}u_x - r u = 0.
\]

Donc :
\[
-u_\tau + \frac12 \sigma^2 (u_{xx} - u_x) + r u_x - r u = 0,
\]
i.e.
\[
u_\tau = \frac12 \sigma^2 u_{xx}
+ \left(r - \frac12 \sigma^2\right) u_x - r u.
\]

\paragraph{\'Elimination du terme en $u_x$ et en $u$.}

On pose
\[
u(\tau,x) = e^{\alpha x + \beta \tau} w(\tau,x),
\]
et on choisit $\alpha,\beta$ pour transformer l’\'equation en \emph{\'equation de la chaleur}
pour $w$.

On calcule :
\[
u_x = e^{\alpha x + \beta \tau} (\alpha w + w_x), \quad
u_{xx} = e^{\alpha x + \beta \tau} (\alpha^2 w + 2\alpha w_x + w_{xx}),
\]
\[
u_\tau = e^{\alpha x + \beta \tau} (\beta w + w_\tau).
\]

En substituant dans l’\'equation pour $u$ :

\[
\beta w + w_\tau
=
\frac12 \sigma^2(\alpha^2 w + 2\alpha w_x + w_{xx})
+ \left(r - \frac12 \sigma^2\right)(\alpha w + w_x) - r w.
\]

On regroupe les termes en $w$, $w_x$, $w_{xx}$ :

\[
w_\tau
=
\frac12 \sigma^2 w_{xx}
+
\Bigl( \sigma^2 \alpha + r - \tfrac12\sigma^2 \Bigr) w_x
+
\Bigl( \tfrac12 \sigma^2\alpha^2 + (r - \tfrac12\sigma^2)\alpha - r - \beta \Bigr) w.
\]

Pour obtenir une \'equation de la chaleur pure :
\[
w_\tau = \frac12 \sigma^2 w_{xx},
\]
il faut que le coefficient de $w_x$ soit nul et que celui de $w$ soit nul, d’o\`u :

\[
\sigma^2 \alpha + r - \tfrac12 \sigma^2 = 0,
\qquad
\tfrac12 \sigma^2\alpha^2 + (r - \tfrac12\sigma^2)\alpha - r - \beta = 0.
\]

On choisit :
\[
\alpha = \frac12 - \frac{r}{\sigma^2},
\]
puis
\[
\beta = -\frac12 \sigma^2 \alpha^2 - (r - \tfrac12\sigma^2)\alpha + r.
\]

Avec ce choix, $w$ satisfait :
\[
w_\tau = \frac12 \sigma^2 w_{xx},
\]
c’est-à-dire l’\'equation de la chaleur classique.

\paragraph{Condition initiale.}

La condition terminale $V(T,S) = (S-K)^+$ se traduit en condition initiale pour $w$ :
\[
u(0,x) = V(T,e^x) = (e^x - K)^+,
\]
\[
w(0,x) = e^{-\alpha x} u(0,x)
= e^{-\alpha x} (e^x - K)^+.
\]

L’\'equation de la chaleur se r\'esout par convolution avec le noyau gaussien.
On obtient au final, apr\`es transformation inverse, la formule de Black--Scholes :
\[
C_0 = S_0 N(d_1) - K e^{-rT} N(d_2),
\]
avec
\[
d_{1,2} =
\frac{\ln(S_0/K) + (r \pm \tfrac12 \sigma^2)T}{\sigma \sqrt{T}}.
\]

%%%%%%%%%%%%%%%%%%%%%%%%%%%%%%%%%%%%%%%%%%%%%%%%%%%%%%%%%%%%%%
\subsubsection{Arbres binomiaux (Cox--Ross--Rubinstein) : exercice et convergence}

\paragraph{Principe.}
Un arbre binomial approxime l'\'evolution de $S_t$ par une marche \`a pas discrets : sur un pas $\Delta t$,
le sous-jacent passe de $S$ \`a $uS$ (``up'') ou $dS$ (``down''). Le sch\'ema recombinent CRR choisit classiquement
\[
u=\ee^{\sigma\sqrt{\Delta t}},\qquad d=\ee^{-\sigma\sqrt{\Delta t}}=\frac{1}{u}.
\]
Sous l'absence d'arbitrage et en pr\'esence d'un portage $(r-q)$, la probabilit\'e risque--neutre est
\[
p=\frac{\ee^{(r-q)\Delta t}-d}{u-d},
\]
et la valeur se calcule par r\'ecurrence :
\[
V(t,S)=\ee^{-r\Delta t}\Bigl(p\,V(t+\Delta t,uS)+(1-p)\,V(t+\Delta t,dS)\Bigr),
\]
avec condition terminale $V(T,S)=\Phi(S)$.

\paragraph{Exercice anticip\'e (options am\'ericaines).}
Pour une option am\'ericaine, le principe est une comparaison \emph{valeur d'exercice} vs \emph{valeur de continuation} :
\[
V(t,S)=\max\bigl(\Phi_{\text{ex}}(S),\;V_{\text{cont}}(t,S)\bigr),
\]
ce qui met en \'evidence la fronti\`ere d'exercice (free boundary) et justifie les algorithmes d'induction arri\`ere.

\paragraph{Convergence et limites.}
Sous des hypoth\`eses usuelles, le prix CRR converge vers Black--Scholes quand $\Delta t\to 0$ (th\'eor\`eme de Donsker,
convergence faible vers le brownien g\'eom\'etrique). Les limites sont bien connues : complexit\'e qui cro\^it avec le nombre
de facteurs, sensibilit\'e aux hypoth\`eses (dividendes, taux), et traitement non trivial des barri\`eres/path-dependence
(m\^eme si des extensions existent).

%%%%%%%%%%%%%%%%%%%%%%%%%%%%%%%%%%%%%%%%%%%%%%%%%%%%%%%%%%%%%%
\subsubsection{Monte Carlo : estimateurs, erreur et r\'eduction de variance}

\paragraph{Estimateur.}
Dans un cadre risque--neutre, le prix d'un payoff europ\'een $\Phi(S_T)$ s'\'ecrit
$V_0=\ee^{-rT}\E^{\QQ}[\Phi(S_T)]$. Monte Carlo approxime cette esp\'erance par
\[
\widehat V_0=\ee^{-rT}\frac{1}{N}\sum_{n=1}^N \Phi\!\left(S_T^{(n)}\right),
\]
o\`u $S_T^{(n)}$ sont des trajectoires simul\'ees sous $\QQ$.

\paragraph{Erreur statistique.}
Par le th\'eor\`eme central limite, $\sqrt{N}(\widehat V_0-V_0)$ converge en loi vers une gaussienne de variance
$\mathrm{Var}^{\QQ}(\Phi(S_T))$, d'o\`u une erreur typiquement en $\mathcal{O}(N^{-1/2})$. Une erreur suppl\'ementaire
provient de la discr\'etisation en temps lorsqu'on simule une EDS (biais en $\Delta t$).

\paragraph{R\'eduction de variance.}
On peut r\'eduire la variance \`a $N$ fix\'e via : variables antith\'etiques, \emph{control variates}
(par exemple utiliser une option \`a prix ferm\'e comme contr\^ole), importance sampling, stratification ou quasi-MC.
Ces techniques modifient l'estimateur sans changer la quantit\'e cibl\'ee.

\paragraph{Greeks.}
Les sensibilit\'es se calculent par diff\'erenciation (pathwise), par \emph{likelihood ratio} (score function), ou par
diff\'erences finies (\emph{bump and revalue}). Les payoffs non lisses (digital, barri\`eres) demandent des pr\'ecautions :
approximation liss\'ee ou estimateurs adapt\'es.

\paragraph{Limites.}
Monte Carlo est naturellement adapt\'e aux payoffs path-dependent et \`a la grande dimension, mais il est lent sur des
quantit\'es rares (queues) et sur des options \`a exercice anticip\'e : pour ces derni\`eres, des m\'ethodes comme le
Least-Squares Monte Carlo (LSMC) introduisent une r\'egression et un risque d'approximation.

%%%%%%%%%%%%%%%%%%%%%%%%%%%%%%%%%%%%%%%%%%%%%%%%%%%%%%%%%%%%%%
\subsubsection{Fonctions caract\'eristiques, inversion de Fourier et FFT (Carr--Madan)}

\paragraph{Id\'ee.}
Lorsque la fonction caract\'eristique du log-prix est connue,
$\phi_T(u)=\E[\ee^{\ii u \ln S_T}]$, on peut repr\'esenter des prix europ\'eens par inversion de Fourier.
L'avantage est de remplacer des int\'egrales difficiles par des int\'egrales oscillantes contr\^ol\'ees, voire par une FFT.

\paragraph{Cadre Carr--Madan.}
Carr et Madan introduisent un \emph{damping} pour rendre la transform\'ee int\'egrable : en notant $k=\ln K$ et
$c_\alpha(k)=\ee^{\alpha k}C(k)$ avec $\alpha>0$, la transform\'ee de $c_\alpha$ s'exprime explicitement en fonction de $\phi_T$.
Une discr\'etisation sur une grille de fr\'equences permet ensuite de calculer simultan\'ement de nombreux prix de strikes par FFT.

\paragraph{Choix du damping et stabilit\'e num\'erique.}
Le choix de $\alpha$ contr\^ole l'int\'egrabilit\'e et la stabilit\'e : trop faible, l'int\'egrale diverge ; trop fort,
on amplifie des erreurs num\'eriques. D'autres sources d'erreur proviennent de la troncature, de l'aliasing et des pas de grille.
Une validation th\'eorique consiste \`a contr\^oler la convergence en augmentant les bornes d'int\'egration et en raffinant les grilles.

\paragraph{Port\'ee.}
Les m\'ethodes de Fourier/FFT sont particuli\`erement adapt\'ees aux mod\`eles \`a structure affine (Heston, Bates) et aux mod\`eles \`a sauts.
Elles ciblent principalement des payoffs europ\'eens ; les extensions \`a l'am\'ericain ou \`a des exotiques requi\`erent des techniques suppl\'ementaires.

\paragraph{Lien avec le projet \emph{PaperTradingApp}.}
Dans le projet, l'int\'er\^et th\'eorique de ces m\'ethodes est double : (i) disposer d'une famille de pricers coh\'erents entre mod\`eles (diffusion,
volatilité stochastique, sauts) via la m\^eme grammaire ``fonction caract\'eristique + inversion'', et (ii) comparer les compromis pr\'ecision/co\^ut
entre PDE, arbres, Monte Carlo et Fourier, sans d\'ependre d'une impl\'ementation particuli\`ere.

\noindent
\subsubsection{Mod\`ele de Heston et pricing semi--analytique}

On consid\`ere d\'esormais le mod\`ele de Heston, dans lequel le sous--jacent $S_t$ et
la variance instantan\'ee $v_t$ suivent la dynamique (sous la mesure risque--neutre $\mathbb{Q}$) :
\[
\begin{cases}
dS_t = r S_t \, dt + \sqrt{v_t}\, S_t \, dW_t^{(1)}, \\[4pt]
dv_t = \kappa(\theta - v_t)\, dt + \xi \sqrt{v_t} \, dW_t^{(2)},
\end{cases}
\qquad
dW_t^{(1)} dW_t^{(2)} = \rho\, dt.
\]

O\`u :
\begin{itemize}
  \item $\kappa > 0$ : vitesse de retour vers la moyenne,
  \item $\theta > 0$ : variance de long terme,
  \item $\xi > 0$ : vol-of-vol,
  \item $\rho \in [-1,1]$ : corr\'elation entre prix et variance.
\end{itemize}

\paragraph{Objectif.}

Pour un call europ\'een de strike $K$ et maturit\'e $T$, on cherche :
\[
C_0 = e^{-rT}\,\mathbb{E}^{\mathbb{Q}}[(S_T - K)^+].
\]

Dans le mod\`ele de Heston, le prix admet une formule semi--ferm\'ee en termes de 
transform\'ee de Fourier, gr\^ace au caract\`ere affine du mod\`ele.

\paragraph{Changement de variable.}

On introduit le log--prix :
\[
X_t = \ln S_t.
\]

En appliquant le lemme d’It\^{o} \`a $X_t$ (comme dans la section pr\'ec\'edente), on obtient :
\[
dX_t = \left(r - \frac12 v_t\right)\, dt + \sqrt{v_t} \, dW_t^{(1)}.
\]

Le couple $(X_t,v_t)$ forme un processus affine de dimension 2.

\paragraph{Fonction caract\'eristique.}

On consid\`ere la fonction caract\'eristique du log--prix sous $\mathbb{Q}$ :
\[
\phi(u; t) := \mathbb{E}^{\mathbb{Q}}\big[ e^{iu X_t} \,\big|\, X_0, v_0 \big].
\]

Le caract\`ere affine du mod\`ele implique que la fonction caract\'eristique s’\'ecrit sous la forme :
\[
\phi(u; t) = 
\exp\Big( C(u,t) + D(u,t) v_0 + iu X_0 \Big),
\]
o\`u $C(u,t)$ et $D(u,t)$ satisfont des \'{e}quations diff\'erentielles ordinaires de type Riccati.

Plus pr\'ecis\'ement, on montre que :
\[
\begin{cases}
\displaystyle \frac{\partial D}{\partial t}(u,t)
= \frac12 \xi^2 D(u,t)^2 + (\rho \xi iu - \kappa) D(u,t) - \frac12 (u^2 + iu), \\[8pt]
\displaystyle \frac{\partial C}{\partial t}(u,t)
= \kappa \theta D(u,t) + i u r,
\end{cases}
\]
avec conditions initiales $C(u,0)=0$, $D(u,0)=0$.

La r\'esolution explicite de ces EDO (voir annexe consacr\'ee \`a Heston) donne des
expressions ferm\'ees pour $C$ et $D$, d'o\`u une expression ferm\'ee de $\phi(u;T)$.

\paragraph{Formule de pricing (int\'egrale de Fourier).}

Le prix du call peut alors s'\'ecrire comme :
\[
C_0 = S_0 P_1 - K e^{-rT} P_2,
\]
o\`u
\[
P_j = \frac12 + \frac{1}{\pi} \int_0^{+\infty} 
\Re\left(
\frac{e^{-iu \ln K} \phi_j(u)}{iu}
\right)\, du,
\qquad j=1,2,
\]
et $\phi_j(u)$ sont des variantes de la fonction caract\'eristique, obtenues en choisissant
des arguments complexes diff\'erents (voir la d\'erivation classique de Heston).

\medskip

\noindent
\textit{Remarque.} La d\'erivation compl\`ete des expressions ferm\'ees de $C(u,t)$, $D(u,t)$
et des int\'egrales de pricing est fournie en annexe. 
Elle repose sur la r\'esolution explicite des EDO de Riccati et sur l'utilisation 
de la transform\'ee de Fourier invers\'ee.


\subsubsection{Surface de volatilit\'e implicite et calibration de mod\`eles}

La volatilité implicite est définie comme la valeur de $\sigma$ qui, insérée dans 
la formule de Black--Scholes, reproduit le prix de marché d’un call europ\'een :
\[
C^{BS}(S_0,K,r,\sigma,T) = C^{\mathrm{mkt}}(K,T).
\]
En inversant cette relation pour chaque strike $K$ et chaque maturité $T$, on
obtient la \emph{surface de volatilité implicite} :
\[
(K,T) \;\longmapsto\; \sigma_{\mathrm{impl}}(K,T).
\]

%%%%%%%%%%%%%%%%%%%%%%%%%%%%%%%%%%%%%%%%%%%%%%%%%%%%%%%%%%%%%%
\paragraph{Contraintes d’arbitrage sur la surface.}

Pour qu’une surface de volatilité soit compatible avec l’absence d’arbitrage
statique, certaines conditions doivent être satisfaites.

\begin{itemize}
  \item \textbf{Monotonie en maturité (pas d’arbitrage de calendrier)} :
  \[
  T_2 > T_1 \quad\Rightarrow\quad C(K,T_2) \ge C(K,T_1).
  \]

  \item \textbf{Convexité en strike (pas d’arbitrage butterfly)} :
  \[
  \frac{\partial^2 C}{\partial K^2}(K,T) \ge 0.
  \]
  Cette condition assure l’existence d’une densité risque--neutre positive :
  \[
  \frac{\partial^2 C}{\partial K^2}(K,T)
  = e^{-rT} f^{\mathbb{Q}}_{S_T}(K).
  \]
\end{itemize}

Ces contraintes doivent être vérifiées par toute surface interpolée/paramétrée
avant calibration.

%%%%%%%%%%%%%%%%%%%%%%%%%%%%%%%%%%%%%%%%%%%%%%%%%%%%%%%%%%%%%%
\paragraph{Calibration de mod\`eles sur la surface de volatilit\'e.}

Un mod\`ele param\'etr\'e $\theta$ (par exemple Heston, SABR, rough volatility, etc.)
génère une surface implicite de volatilité :
\[
\sigma_{\mathrm{mod}}(K,T;\theta).
\]

La calibration consiste à choisir $\theta$ pour que cette surface reproduise au
mieux la surface observée sur le marché. Le problème d’optimisation s’écrit
classiquement :
\[
\theta^\*
=
\arg\min_{\theta}
\sum_{(K,T)\in \mathcal{G}}
w_{K,T}\,
\bigl(
\sigma_{\mathrm{mod}}(K,T;\theta)
-
\sigma_{\mathrm{impl}}^{\mathrm{mkt}}(K,T)
\bigr)^2,
\]
où :
\begin{itemize}
  \item $\mathcal{G}$ est la grille de strikes et de maturités cotés,
  \item $w_{K,T}$ sont des poids reflétant volumes, liquidité ou importance stratégique.
\end{itemize}

%%%%%%%%%%%%%%%%%%%%%%%%%%%%%%%%%%%%%%%%%%%%%%%%%%%%%%%%%%%%%%
\paragraph{Calibration de mod\`eles stochastiques (exemple : Heston).}

Dans le mod\`ele de Heston, les param\`etres
\[
\theta = (\kappa,\theta,\xi,\rho,v_0)
\]
sont ajustés en évaluant les prix d’options via la formule semi--fermée basée
sur la fonction caractéristique. Pour chaque $(K,T)$, on en déduit
\[
\sigma_{\mathrm{mod}}(K,T;\theta)
\quad\text{telle que}\quad
C^{BS}(S_0,K,r,\sigma_{\mathrm{mod}},T)
= C^{\mathrm{Heston}}(S_0,K,r,\theta,T).
\]

La calibration est un problème non convexe : on utilise généralement des méthodes
de type multi-start, recuit simulé, CMA-ES, ou une descente locale après une
initialisation robuste.

%%%%%%%%%%%%%%%%%%%%%%%%%%%%%%%%%%%%%%%%%%%%%%%%%%%%%%%%%%%%%%
\paragraph{Calibration rough volatility / rBergomi.}

Dans les mod\`eles de volatilité rugueuse, la variance suit une dynamique 
fractionnaire. Les paramètres typiques sont $(H,\eta,\xi_0(t))$, avec $H<1/2$.
La valorisation s'effectue via Monte Carlo rapide (hybrid schemes), puis une
minimisation analogue au cas Heston :
\[
\theta^\*
=
\arg\min_{\theta}
\sum_{(K,T)\in \mathcal{G}}
\bigl(
\sigma_{\mathrm{mod}}(K,T;\theta)
-
\sigma_{\mathrm{impl}}^{\mathrm{mkt}}(K,T)
\bigr)^2.
\]

%%%%%%%%%%%%%%%%%%%%%%%%%%%%%%%%%%%%%%%%%%%%%%%%%%%%%%%%%%%%%%
\paragraph{Synth\`ese.}

\begin{itemize}
  \item La surface de volatilité implicite est la représentation empirique de la
  distribution future du sous--jacent perçue par le marché.
  \item Les contraintes $\partial_{K}^2 C \ge 0$ et $C(K,T_2)\ge C(K,T_1)$ doivent
  être respectées pour éviter l’arbitrage statique.
  \item La calibration consiste à résoudre un problème d’optimisation sur les paramètres
  du modèle afin de reproduire au mieux la surface de marché.
  \item Les modèles stochastiques (Heston, rough volatility) sont calibrés via des
  méthodes globales puis locales, en exploitant leur fonction caractéristique ou
  des schémas Monte Carlo rapides.
\end{itemize}

\paragraph{Surrogates (r\'eseaux) pour acc\'el\'erer la calibration.}
Dans de nombreux mod\`eles, le goulot d'\'etranglement est le \emph{co\^ut} d'\'evaluation de la carte
$(K,T)\mapsto C_{\theta}(K,T)$ (ou de $\sigma_{\mathrm{mod}}(K,T;\theta)$) \`a chaque it\'eration d'optimisation.
L'id\'ee d'un \textbf{surrogate} est de construire une approximation rapide
\[
(K,T,\theta)\;\longmapsto\; \widehat C(K,T;\theta)
\quad \text{ou} \quad
(K,T,\theta)\;\longmapsto\;\widehat\sigma_{\mathrm{imp}}(K,T;\theta),
\]
par apprentissage supervis\'e sur des donn\'ees synth\'etiques g\'en\'er\'ees par le mod\`ele lui-m\^eme.
Le fondement th\'eorique est l'approximation universelle : sous des hypoth\`eses de r\'egularit\'e,
un r\'eseau suffisamment riche peut approximer une fonction continue sur un compact.

\paragraph{Int\'er\^et, limites et validation.}
Un surrogate peut servir \`a (i) initialiser une calibration (warm start), (ii) explorer rapidement
la sensibilit\'e \`a des param\`etres, (iii) d\'eployer des diagnostics en quasi temps r\'eel. En revanche,
les risques sont conceptuellement clairs : biais d'approximation, extrapolation hors du domaine d'apprentissage,
et possible violation d'absence d'arbitrage. Une bonne pratique th\'eorique est d'imposer des \emph{contraintes}
(par exemple convexit\'e en $K$) ou de projeter l'approximation sur un ensemble arbitrage-free, et de valider
sur un jeu out-of-sample en contr\^olant (a) l'erreur de prix/vol, (b) la stabilit\'e des param\`etres obtenus,
et (c) la robustesse \`a des perturbations de donn\'ees (stress, bid--ask, maturit\'es manquantes).


\subsubsection{Mod\`eles avanc\'es : SABR et \'Equations de Volterra}

Dans le cadre de la mod\'elisation des smiles de volatilité, deux familles de
mod\`eles occupent une place particuli\`ere : les mod\`eles SABR et les
mod\`eles \`a noyaux de Volterra. Ils fournissent des repr\'esentations plus
flexibles des dynamiques de volatilité observées sur les march\'es.

%%%%%%%%%%%%%%%%%%%%%%%%%%%%%%%%%%%%%%%%%%%%%%%%%%%%%%%%%%%%%%
\paragraph{Mod\`ele SABR.}

Le mod\`ele SABR (Stochastic Alpha Beta Rho), introduit par Hagan et al., est
largement utilis\'e en taux d’int\'er\^et pour reproduire les smiles de
volatilit\'e des swaptions. Sous la mesure risque--neutre, le forward $F_t$
\'evolue selon :
\[
\begin{cases}
dF_t = \alpha_t F_t^{\beta}\, dW_t^{(1)}, \\[4pt]
d\alpha_t = \nu \alpha_t\, dW_t^{(2)}, \\[4pt]
dW_t^{(1)} dW_t^{(2)} = \rho\, dt.
\end{cases}
\]

Les param\`etres ont les interpr\'etations suivantes :
\begin{itemize}
  \item $\beta \in [0,1]$ : contr\^ole la forme du smile (lognormal, CEV, normal),
  \item $\nu > 0$ : vol-of-vol, responsable de la convexit\'e du smile,
  \item $\rho$ : corr\'elation, pilote le skew,
  \item $\alpha_t$ : volatilité instantanée du forward.
\end{itemize}

Une approximation asymptotique (formule de Hagan) fournit une expression fermée de la
volatilité implicite :
\[
\sigma_{\mathrm{SABR}}^{\mathrm{impl}}(K)
=
\frac{\alpha_0}{(F_0 K)^{\frac{1-\beta}{2}}}
\left[ 1 + \text{Corrections}_1(\nu,\rho,\beta,K) \right],
\]
ce qui permet une calibration directe du smile. Cette formule est obtenue par expansion
asymptotique en $\nu$ et les d\'erivations compl\`etes sont fournies en \textbf{ANNEXE}
(voir fichiers associ\'es au mod\`ele SABR).

%%%%%%%%%%%%%%%%%%%%%%%%%%%%%%%%%%%%%%%%%%%%%%%%%%%%%%%%%%%%%%
\paragraph{Mod\`eles de Volterra et volatilité rugueuse.}

Les mod\`eles Volterra-g\'en\'eriques consid\`erent des dynamiques \`a m\'emoire
longue :
\[
dX_t 
= \left( \int_0^t K(t-s)\,g(X_s)\,ds \right) dt
+ \left( \int_0^t K(t-s)\,h(X_s)\,ds \right) dW_t,
\]
o\`u $K$ est un noyau de Volterra causal. Le choix du noyau d\'etermine la
r\'egularité et la structure de m\'emoire du mod\`ele.

Un cas fondamental est la \textbf{volatilité rugueuse (rough volatility)}, fondée 
sur un \emph{fractional Brownian motion} (fBM) de param\`etre de Hurst $H<1/2$ :
\[
B_t^H = \frac{1}{\Gamma(H+\frac12)}
\int_0^t (t-s)^{H-\frac12} dW_s.
\]

Dans le mod\`ele \textbf{rBergomi}, la variance suit :
\[
v_t = \xi_0(t)\,
\exp\!\left(\eta B_t^H - \tfrac12 \eta^2 t^{2H}\right),
\]
o\`u $\eta$ est la vol-of-vol rugueuse et $H$ r\'egit la r\'egularit\'e du
processus.

Les mod\`eles rough permettent de reproduire fid\`element :
\begin{itemize}
  \item le skew court terme des march\'es equity,
  \item la persistance et la m\'emoire longue de la variance,
  \item les d\'eformations fines de la surface de vol.
\end{itemize}

Les d\'etails de construction du mouvement brownien fractionnaire, des noyaux de Volterra et des sch\'emas de simulation
peuvent \^etre rel\'egu\'es en annexes (ou \`a la litt\'erature) ; on se limite ici aux d\'efinitions, \`a l'intuition et aux implications
pour la g\'eom\'etrie court terme du smile.


\subsubsection{Lien avec le projet \emph{PaperTradingApp}}

Le projet s'appuie sur ces briques pour articuler : (i) un pricing ``de r\'ef\'erence'' (Black--Scholes) qui sert de point de comparaison,
(ii) des m\'ethodes num\'eriques (arbres, Monte Carlo, Fourier) n\'ecessaires d\`es que l'on traite des exotiques, des exercices anticip\'es
ou des mod\`eles non r\'esolubles en forme ferm\'ee, et (iii) un continuum entre \emph{prix} et \emph{surface} via l'inversion en volatilité implicite.
Le lien essentiel avec la plateforme est la capacit\'e \`a justifier le choix d'une m\'ethode (pr\'ecision, co\^ut, hypoth\`eses) et \`a contr\^oler
les erreurs (discr\'etisation, troncature, instabilit\'e de calibration) de mani\`ere th\'eorique.

\subsection{Hedging et Greeks}


%%%%%%%%%%%%%%%%%%%%%%%%%%%%%%%%%%%%%%%%%%%%%%%%%%%%%%%%%%%%%%
\subsubsection{Greeks}

Les Greeks quantifient les sensibilités du prix d’un dérivé aux différents
paramètres de marché. Pour une option de prix $V(t,S)$ :

\paragraph{Delta.}
\[
\Delta = \frac{\partial V}{\partial S}.
\]
Sensibilité directionnelle à une variation du sous-jacent.
Pour un call : $0 < \Delta < 1$ ; pour un put : $-1 < \Delta < 0$.

\paragraph{Gamma.}
\[
\Gamma = \frac{\partial^2 V}{\partial S^2}.
\]
Mesure la convexité du prix.
Un portefeuille long gamma bénéficie des variations du sous-jacent.

\paragraph{Vega.}
\[
\text{Vega} = \frac{\partial V}{\partial \sigma}.
\]
Sensibilité à la volatilité implicite. Maximum autour de l’ATM.

\paragraph{Theta.}
\[
\Theta = \frac{\partial V}{\partial t}.
\]
Sensibilité à l’écoulement du temps. Les options vanilles longues sont
généralement $\Theta < 0$.

\paragraph{Rho.}
\[
\rho_{\mathrm{IR}} = \frac{\partial V}{\partial r}.
\]
Sensibilité au taux sans risque, importante pour FX / taux ; on la note $\rho_{\mathrm{IR}}$ pour éviter la confusion avec la corrélation $\rho$.

%%%%%%%%%%%%%%%%%%%%%%%%%%%%%%%%%%%%%%%%%%%%%%%%%%%%%%%%%%%%%%
\subsubsection{Hedging}

\paragraph{Delta--hedging.}

On construit un portefeuille :
\[
\Pi_t = V(t,S_t) - \Delta_t S_t.
\]
En choisissant $\Delta_t = V_S(t,S_t)$, le terme aléatoire $dW_t$ disparaît de
la dynamique de $\Pi_t$.
Cela conduit au portefeuille localement sans risque utilisé dans la dérivation
de Black--Scholes.

\paragraph{Gamma--Vega hedging.}

Une seule option ne permet pas de neutraliser simultanément $\Gamma$ et Vega.
On combine deux options $V^{(1)}$ et $V^{(2)}$ via :
\[
\begin{cases}
x \Gamma^{(1)} + y \Gamma^{(2)} = 0, \\
x \text{Vega}^{(1)} + y \text{Vega}^{(2)} = 0,
\end{cases}
\]
pour construire un portefeuille insensible aux variations de volatilité et de
convexité.

%%%%%%%%%%%%%%%%%%%%%%%%%%%%%%%%%%%%%%%%%%%%%%%%%%%%%%%%%%%%%%
\subsubsection{D\'erivation du mod\`ele de Black--Scholes par delta--hedging}

On part de la dynamique du sous-jacent :
\[
dS_t = \mu S_t\, dt + \sigma S_t\, dW_t.
\]

Le prix $V(t,S_t)$ d'une option européenne vérifie (par développement stochastique) :
\[
dV_t 
=
\left(
V_t 
+ \mu S_t V_S 
+ \frac12 \sigma^2 S_t^2 V_{SS}
\right) dt
+ \sigma S_t V_S\, dW_t.
\]

%%%%%%%%%%%%%%%%%%%%%%%%%%%%%%%%%%%%%%%%%%%%%%%%%%%%%%%%%%%%%%
\paragraph{Portefeuille de couverture.}

On considère :
\[
\Pi_t = V(t,S_t) - \Delta_t S_t.
\]

Sa variation :
\[
d\Pi_t = dV_t - \Delta_t\, dS_t.
\]

Après substitution :
\[
\begin{aligned}
d\Pi_t
&=
\left(
V_t 
+ \mu S_t V_S
+ \frac12 \sigma^2 S_t^2 V_{SS}
\right) dt
- \Delta_t \mu S_t\, dt \\
&\quad
+ \sigma S_t( V_S - \Delta_t )\, dW_t.
\end{aligned}
\]

%%%%%%%%%%%%%%%%%%%%%%%%%%%%%%%%%%%%%%%%%%%%%%%%%%%%%%%%%%%%%%
\paragraph{Élimination du risque.}

On impose que le terme en $dW_t$ disparaisse :
\[
V_S - \Delta_t = 0 
\quad\Longrightarrow\quad
\Delta_t = V_S.
\]

Le portefeuille devient alors localement sans risque :
\[
d\Pi_t
=
\left(
V_t 
+ \frac12 \sigma^2 S_t^2 V_{SS}
\right) dt.
\]

%%%%%%%%%%%%%%%%%%%%%%%%%%%%%%%%%%%%%%%%%%%%%%%%%%%%%%%%%%%%%%
\paragraph{Absence d'arbitrage : taux sans risque.}

Un portefeuille sans risque doit croître au taux $r$ :
\[
d\Pi_t = r\, \Pi_t\, dt = r( V - S V_S )\, dt.
\]

On obtient donc :
\[
V_t + \frac12 \sigma^2 S^2 V_{SS}
= r(V - S V_S).
\]

%%%%%%%%%%%%%%%%%%%%%%%%%%%%%%%%%%%%%%%%%%%%%%%%%%%%%%%%%%%%%%
\paragraph{PDE de Black--Scholes.}

En réarrangeant :
\[
V_t 
+ \frac12 \sigma^2 S^2 V_{SS}
+ r S V_S 
- r V
= 0.
\]

Avec condition terminale :
\[
V(T,S) = (S-K)^+,
\]
on retrouve la formule fermée de Black--Scholes :
\[
C_0 
= S_0 N(d_1) - K e^{-rT} N(d_2),
\]
avec
\[
d_{1,2}
=
\frac{
\ln(S_0/K) + ( r \pm \tfrac12 \sigma^2 )T
}{
\sigma\sqrt{T}
}.
\]

\noindent
Cette dérivation illustre que la formule de Black--Scholes provient directement
du principe fondamental : \emph{un portefeuille delta--hedgé ne doit pas permettre d’arbitrage}.

\subsubsection{Couverture et d\'ecision s\'equentielle : erreurs, co\^uts et formulation MDP}

\paragraph{Erreurs de couverture (discr\'etisation et risque de mod\`ele).}
Dans Black--Scholes frictionless, une couverture continue en delta r\'eplique le payoff. En pratique, la couverture est
\emph{discr\`ete} et le mod\`ele est \emph{approximatif}. L'erreur de couverture \`a maturit\'e peut se d\'efinir comme
\[
\varepsilon_T := V_T^{\text{couvert}} - \Phi(S_T),
\]
et s'interpr\`ete comme la partie non r\'epliqu\'ee (variance r\'ealis\'ee diff\'erente de la variance mod\`ele, sauts,
volatilité stochastique, smile dynamics). Un message cl\'e est que l'exposition \`a la convexit\'e ($\Gamma$) rend la
couverture sensible \`a la variation quadratique r\'eelle du sous-jacent.

\paragraph{Smile dynamics et couverture multi-Greeks.}
Hedger seulement le delta suppose implicitement que la dynamique de volatilité est ``int\'egr\'ee'' dans le mod\`ele.
D\`es que l'on observe un smile/skew, la question de la dynamique du smile devient centrale (hypoth\`eses du type
``sticky strike'' vs ``sticky delta''). Dans un cadre incomplet, une couverture plus riche vise \`a g\'erer des
expositions $\Delta/\Gamma/\nu$ en combinant sous-jacent et options/liquidit\'es, au prix d'un risque r\'esiduel
et d'une d\'ependance au mod\`ele.

\paragraph{Co\^uts de transaction.}
Les co\^uts (bid--ask, impact) rendent la couverture plus ``inertielle'' : un ajustement fr\'equent r\'eduit le risque
mais augmente les co\^uts. Th\'eoriquement, cela se traduit par un compromis \`a optimiser, et par une modification de la
valeur ``r\'epliquante'' (par exemple via des corrections effectives de volatilité dans certains cadres id\'ealis\'es).

\paragraph{Formulation MDP du hedging.}
On peut formaliser la couverture comme un \emph{Markov Decision Process} : l'\'etat $s_t$ r\'esume l'information utile
(temps restant, spot, indicateurs de volatilité, position et Greeks), l'action $a_t$ d\'efinit l'ajustement de portefeuille,
et la r\'ecompense p\'enalise simultan\'ement l'erreur et les co\^uts, par exemple sous la forme
\[
R_{t\to t+\Delta}(s_t,a_t) \;=\; \mathrm{PnL}_{t\to t+\Delta} \;-\; \text{co\^uts} \;-\; \lambda_{\mathrm{risk}}\,\text{risque}.
\]
Cette \'ecriture rend explicite le compromis \emph{risque/co\^ut} et la non-stationnarit\'e (changements de r\'egime, d\'eformation du smile).

\paragraph{Ouverture RL et validation.}
Des algorithmes de RL (Q-learning, DQN, PPO) peuvent alors approximer une politique de couverture. Au niveau th\'eorique, la convergence
est bien comprise dans le cas tabulaire sous hypoth\`eses fortes ; avec approximation de fonctions, la stabilit\'e devient une question
centrale (sur-apprentissage, d\'erives de distribution, non-stationnarit\'e). Une validation rigoureuse doit \^etre out-of-sample,
inclure des stress tests (spot/vol/taux), et rapporter la sensibilit\'e aux co\^uts et \`a la qualit\'e des donn\'ees.

\paragraph{Lien avec le projet \emph{PaperTradingApp}.}
Cette formalisation relie directement : (i) les Greeks et erreurs de couverture, (ii) les contraintes de march\'e (co\^uts, liquidit\'e),
et (iii) une boucle de d\'ecision s\'equentielle. Elle fournit un cadre conceptuel pour comparer une couverture ``r\`egle'' (delta/vega)
\`a une couverture issue d'une politique optimis\'ee, sans pr\'esumer de r\'esultats empiriques.

\subsection{Taux \& courbes de taux}

\subsubsection{Taux z\'ero--coupon et courbe des taux}

Le prix d’un z\'ero--coupon de maturit\'e $T$ est not\'e $P(0,T)$ : c’est la valeur
aujourd’hui d’un instrument qui versera exactement $1$ unit\'e \`a la date $T$.

Le taux z\'ero--coupon (capitalisation continue) est d\'efini par :
\[
P(0,T) = e^{-R(0,T)\,T},
\qquad
R(0,T) := -\frac{1}{T}\ln P(0,T).
\]

La fonction :
\[
T \longmapsto R(0,T)
\]
est appel\'ee \emph{courbe des taux spots}.

%%%%%%%%%%%%%%%%%%%%%%%%%%%%%%%%%%%%%%%%%%%%%%%%%%%%%%%%%%%%%%
\subsubsection{Facteurs d’actualisation (Discount Factors)}

Le facteur d’actualisation associ\'e \`a la date $T$ est :
\[
D(0,T) = P(0,T) = e^{-R(0,T)\,T}.
\]

Pour un paiement certain $C$ re\c{c}u \`a la date $T$, la valeur actuelle est :
\[
\text{VA} = C\, D(0,T).
\]

Ces facteurs constituent la base du calcul de la valeur de tous les flux futurs
(obligations, swaps, cash-flows divers).

%%%%%%%%%%%%%%%%%%%%%%%%%%%%%%%%%%%%%%%%%%%%%%%%%%%%%%%%%%%%%%
\subsubsection{Taux forward instantan\'e}

Le \emph{taux forward instantan\'e} $f(0,T)$ est le taux d’int\'er\^et implicite
pour un emprunt infinitésimal entre les dates $T$ et $T + dT$, d\'eduit \`a partir
du prix des z\'ero-coupons.

Il est d\'efini par :
\[
f(0,T) := -\frac{\partial}{\partial T} \ln P(0,T).
\]

On en d\'eduit :
\[
P(0,T) 
= 
\exp\!\left( -\int_0^T f(0,u)\,du \right).
\]

La quantit\'e $f(0,T)$ peut \^etre vue comme la \emph{densit\'e instantan\'ee}
du taux d’actualisation. Elle joue un r\^ole central dans les mod\`eles de taux
de type Heath--Jarrow--Morton (HJM), o\`u la dynamique de $f(t,T)$ d\'etermine
l’ensemble de la structure par terme.

\medskip

Dans la pratique, le taux forward instantan\'e est approch\'e en diff\'erenciant
num\'eriquement les discount factors bootstrapp\'es.
%%%%%%%%%%%%%%%%%%%%%%%%%%%%%%%%%%%%%%%%%%%%%%%%%%%%%%%%%%%%%%
\subsubsection{Construction de la courbe (bootstrapping)}

La courbe des taux z\'ero-coupon n’est pas directement observable sur le march\'e.
Les march\'es cotent des instruments liquides \`a diff\'erentes maturit\'es
(dépôts interbancaires, FRAs, futures, swaps IRS).  
Le \emph{bootstrapping} consiste \`a reconstruire de mani\`ere cohérente les 
facteurs d’actualisation $P(0,T)$ à partir de ces instruments.

%%%%%%%%%%%%%%%%%%%%%%%%%%%%%%%%%%%%%%%%%%%%%%%%%%%%%%%%%%%%%%
\paragraph{1. Instruments courts (dép\^ots).}

Pour les d\'ep\^ots de maturit\'e $T$ (typiquement $\leq 1$ an), on conna\^it le
taux simple $L(0,T)$.  
On en d\'eduit directement :
\[
P(0,T) = \frac{1}{1 + L(0,T)\,T}.
\]

%%%%%%%%%%%%%%%%%%%%%%%%%%%%%%%%%%%%%%%%%%%%%%%%%%%%%%%%%%%%%%
\paragraph{2. FRAs et futures.}

Les contrats FRA (Forward Rate Agreement) donnent les taux forward implicites
sur une p\'eriode future $[T_1,T_2]$.  
Leur prix impose la relation :
\[
1 = P(0,T_1)\,\bigl(1 + F(0;T_1,T_2)(T_2-T_1)\bigr)\, P(0,T_2)^{-1}.
\]
Ainsi :
\[
P(0,T_2)
=
\frac{P(0,T_1)}{1 + F(0;T_1,T_2)(T_2-T_1)}.
\]

Les futures sur taux court (Eurodollar, SOFR futures, Euribor futures) permettent
d’extraire des taux forward similaires via des conventions d’ajustement.

%%%%%%%%%%%%%%%%%%%%%%%%%%%%%%%%%%%%%%%%%%%%%%%%%%%%%%%%%%%%%%
\paragraph{3. Swaps IRS (maturit\'es longues).}

Pour un swap de maturit\'e $T_n$ et flux semestriels, on observe le taux swap $S_n$
tel que le swap ait une valeur initiale nulle.

La jambe fixe vaut :
\[
\text{Fixe}
= S_n \sum_{k=1}^{n} \alpha_k\, P(0,T_k),
\]
où $\alpha_k$ est la fraction de jour (Act/360, 30/360, etc.).

La jambe flottante vaut :
\[
\text{Flottante} = 1 - P(0,T_n).
\]

Comme $\text{Fixe} = \text{Flottante}$, on obtient :
\[
1 - P(0,T_n)
=
S_n \sum_{k=1}^{n} \alpha_k\, P(0,T_k).
\]

Les facteurs $P(0,T_k)$ pour $k<n$ sont d\'ej\`a connus par bootstrapping
(récursivité), ce qui permet d’obtenir $P(0,T_n)$.

%%%%%%%%%%%%%%%%%%%%%%%%%%%%%%%%%%%%%%%%%%%%%%%%%%%%%%%%%%%%%%
\paragraph{4. Interpolation.}

Une fois les $P(0,T_k)$ obtenus sur un ensemble discret de maturit\'es,
on construit une courbe continue via une interpolation :

\begin{itemize}
  \item sur les taux forwards instants $f(0,T)$,
  \item ou sur les discount factors,
  \item en utilisant des splines monotones ou Hermite préservant la convexité,
  \item en veillant à éviter les violations d’arbitrage 
        (monotonie des discount factors et convexité des prix d’option).
\end{itemize}

%%%%%%%%%%%%%%%%%%%%%%%%%%%%%%%%%%%%%%%%%%%%%%%%%%%%%%%%%%%%%%
\paragraph{5. Courbe finale.}

La courbe compl\`ete des facteurs d’actualisation est donn\'ee par :
\[
P(0,T) = \exp\!\left( -\int_0^T f(0,u)\,du \right),
\]
o\`u $f(0,T)$ est obtenu par dérivation numérique ou interpolation.

Cette courbe sert de base \`a toute valorisation d’instruments \`a taux fixe,
flottant ou dérivés (IRS, cap/floor, swaption, CMS, etc.).


%%%%%%%%%%%%%%%%%%%%%%%%%%%%%%%%%%%%%%%%%%%%%%%%%%%%%%%%%%%%%%
\subsubsection{Param\'etrisation Nelson--Siegel : niveau, pente, courbure}

\paragraph{Objectif.}
Le bootstrapping fournit des facteurs $P(0,T_k)$ sur des maturit\'es discr\`etes. Pour interpoler et surtout \emph{extrapoler}
de fa\c{c}on lisse et interpr\'etable, on utilise souvent des familles param\'etriques. Le mod\`ele de Nelson--Siegel (NS) est un standard
car il d\'ecompose la courbe en composantes ``niveau/pente/courbure''.

\paragraph{Forme (taux z\'ero).}
Une \'ecriture courante du taux z\'ero continu $y(0,T)$ est :
\[
y(0,T)
=
\beta_0

+
\beta_1\,
\frac{1-\ee^{-T/\tau}}{T/\tau}

+
\beta_2\,
\left(
\frac{1-\ee^{-T/\tau}}{T/\tau}
- \ee^{-T/\tau}
\right),
\]
o\`u $\tau>0$ contr\^ole la vitesse de d\'ecroissance des facteurs.

\paragraph{Interpr\'etation.}
\begin{itemize}
  \item $\beta_0$ (\textbf{niveau}) gouverne le niveau de long terme (quand $T\to\infty$).
  \item $\beta_1$ (\textbf{pente}) affecte surtout le court terme ; son effet d\'ecroit avec $T$.
  \item $\beta_2$ (\textbf{courbure}) cr\'ee une bosse \`a maturit\'e interm\'ediaire, pilot\'ee par $\tau$.
\end{itemize}
La version Svensson ajoute un deuxi\`eme couple $(\beta_3,\tau_2)$ pour plus de flexibilit\'e.

\paragraph{Calibration et limites.}
La calibration se fait en minimisant une erreur entre prix observ\'es (ou taux) et ceux impliqu\'es par $P(0,T)$, avec pond\'erations
refl\'etant liquidit\'e et maturit\'e. Les limites sont connues : (i) flexibilit\'e restreinte face \`a des courbes tr\`es non monotones,
(ii) risque de produire des forwards instantan\'es non plausibles (par exemple fortement n\'egatifs) selon les param\`etres,
(iii) d\'ependance \`a la convention de calibration (prix vs taux). Une v\'erification d'absence d'arbitrage (monotonie des $P(0,T)$,
raisonnabilit\'e de $f(0,T)$) est indispensable.


%%%%%%%%%%%%%%%%%%%%%%%%%%%%%%%%%%%%%%%%%%%%%%%%%%%%%%%%%%%%%%
\subsubsection{Lien avec le projet \emph{PaperTradingApp}}

La courbe de taux est un ingr\'edient transversal : elle intervient dans l'actualisation, dans les forwards et donc dans la d\'efinition
de la moneyness et des variables de changement de num\'eraire. Dans le projet, elle est mobilis\'ee (au niveau conceptuel) pour garantir
la coh\'erence ``prix $\leftrightarrow$ taux $\leftrightarrow$ forwards'' et pour \'eviter des incoh\'erences d'absence d'arbitrage li\'ees \`a une interpolation
mal choisie. Elle sert aussi de base \`a des stress tests de taux et \`a des analyses de sensibilit\'e ($\rho_{\mathrm{IR}}$).

\subsection{Forwards, Futures et Swaps}

Les produits linéaires (forwards, futures, swaps) sont fondamentaux en finance
de taux et de devises. Leur valorisation repose exclusivement sur la courbe des
taux d’actualisation et sur les taux forward extraits de celle-ci.

%%%%%%%%%%%%%%%%%%%%%%%%%%%%%%%%%%%%%%%%%%%%%%%%%%%%%%%%%%%%%%
\subsubsection{Forwards}

Un contrat forward fixe aujourd’hui un prix de transaction à une date future $T$.
L’actif n’échange aucun flux avant la maturité.

Sous absence d’arbitrage, le prix forward $F(0,T)$ est donné par :
\[
F(0,T) = S_0 \,\frac{1}{P(0,T)}
       = S_0\, e^{R(0,T)\,T},
\]
où $P(0,T)$ est le facteur d’actualisation.

Avec dividende continu (ou taux de portage $q$) :
\[
F(0,T) = S_0\, e^{(R(0,T)-q)T}.
\]

%%%%%%%%%%%%%%%%%%%%%%%%%%%%%%%%%%%%%%%%%%%%%%%%%%%%%%%%%%%%%%
\subsubsection{Futures}

Les futures sont des forwards \emph{standardisés} négociés en bourse, avec
ajustement quotidien (marche-to-market).

Dans un cadre frictionless, le prix du future vaut en général le forward :
\[
F_{\mathrm{fut}}(0,T) \approx F(0,T),
\]
mais l’ajustement quotidien introduit une différence lorsque $S_t$ est
corrélé au taux court (effet convexité futures).

%%%%%%%%%%%%%%%%%%%%%%%%%%%%%%%%%%%%%%%%%%%%%%%%%%%%%%%%%%%%%%
\subsubsection{Interest Rate Swaps (IRS)}

Un swap de taux échange :
\begin{itemize}
  \item une \emph{jambe fixe} : paiements $S_n \alpha_k$ aux dates $T_k$,
  \item une \emph{jambe flottante} : taux variables (LIBOR, Euribor, SOFR) ajustés à chaque période.
\end{itemize}

La valeur de la jambe fixe est :
\[
\mathrm{Fixe}(0) 
= S_n \sum_{k=1}^{n} \alpha_k\, P(0,T_k).
\]

La jambe flottante vaut :
\[
\mathrm{Flottante}(0) = 1 - P(0,T_n),
\]
car un taux variable réinitialisé assure la parité avec un instrument zéro-coupon.

Le taux swap $S_n$ d’un IRS (valeur initiale nulle) est défini par :
\[
\mathrm{Fixe}(0) = \mathrm{Flottante}(0)
\quad\Longrightarrow\quad
S_n
= 
\frac{1 - P(0,T_n)}{\sum_{k=1}^{n} \alpha_k\, P(0,T_k)}.
\]

%%%%%%%%%%%%%%%%%%%%%%%%%%%%%%%%%%%%%%%%%%%%%%%%%%%%%%%%%%%%%%
\subsubsection{Interprétation économique}

\begin{itemize}
  \item Le \textbf{forward} verrouille un prix d'achat/vente futur.
  \item Le \textbf{future} ajoute la liquidité, standardisation, et marche-to-market.
  \item Le \textbf{swap de taux} échange flux fixes et variables, permettant de
        transformer exposition fixe ↔ variable.
\end{itemize}

Ces produits linéaires servent aussi de brique fondamentale pour le pricing
d’instruments plus complexes (swaption, cap/floor, CMS, etc.).


%%%%%%%%%%%%%%%%%%%%%%%%%%%%%%%%%%%%%%%%%%%%%%%%%%%%%%%%%%%%%%
\subsubsection{Lien avec le projet \emph{PaperTradingApp}}

M\^eme dans un projet centr\'e options, ces instruments lin\'eaires structurent le \emph{portage} : ils clarifient les relations
spot/forward, les conventions de taux et l'actualisation. Ils fournissent donc le cadre conceptuel pour : (i) d\'efinir une moneyness
coh\'erente (via $F(t,T)$), (ii) relier les payoffs \`a des prix forward dans les formules, et (iii) proposer des stress tests
sur les taux (et leur impact sur prix/Greeks) sans ambig\"uit\'e de convention.

\iffalse % Hors p\'erim\`etre principal : section retir\'ee du rendu final
\subsection{Cr\'edit}

La valorisation des instruments de cr\'edit repose sur la notion de risque de
d\'efaut et sur l’intensit\'e (hazard rate) associ\'ee. Les Credit Default Swaps
(CDS) constituent l’instrument de march\'e le plus standard pour extraire ces
informations.

%%%%%%%%%%%%%%%%%%%%%%%%%%%%%%%%%%%%%%%%%%%%%%%%%%%%%%%%%%%%%%
\subsubsection{CDS : d\'efinition et structure}

Un Credit Default Swap est un contrat d’assurance contre le d\'efaut d’un
emprunteur (corporate, souverain). Il \'{e}change :
\begin{itemize}
  \item une \emph{jambe de protection} : paiement $(1 - R)$ en cas de d\'efaut,
  \item une \emph{jambe de prime} : paiements r\'eguliers du spread $s$ tant que
        le d\'efaut n’a pas eu lieu.
\end{itemize}

Le payoff de protection, si le d\'efaut survient en $\tau \in (0,T]$, est :
\[
\mathrm{Protection}
= (1 - R)\, P(0,\tau),
\]
o\`u $R$ est le taux de recouvrement (recovery).

La jambe de prime vaut :
\[
\mathrm{Premium}
=
s \sum_{k=1}^{n}
\alpha_k\, P(0,T_k)\, \mathbb{Q}(\tau > T_k),
\]
avec $\alpha_k$ la fraction de jour.

%%%%%%%%%%%%%%%%%%%%%%%%%%%%%%%%%%%%%%%%%%%%%%%%%%%%%%%%%%%%%%
\subsubsection{Intensit\'e de d\'efaut (hazard rate)}

On mod\'elise le d\'efaut par un temps d’arriv\'ee al\'eatoire $\tau$ dont la
loi est caract\'eris\'ee par une intensit\'e $\lambda(t)$ telle que :
\[
\mathbb{Q}(\tau > t)
= e^{-\int_0^t \lambda(u)\, du}.
\]

Lorsque $\lambda(t) = \lambda$ est constante :
\[
\mathbb{Q}(\tau > t) = e^{-\lambda t}.
\]

La densit\'e de d\'efaut est alors :
\[
f_\tau(t) = \lambda e^{-\lambda t}.
\]

%%%%%%%%%%%%%%%%%%%%%%%%%%%%%%%%%%%%%%%%%%%%%%%%%%%%%%%%%%%%%%
\subsubsection{Spread de CDS}

Un CDS ayant une valeur initiale nulle impose :
\[
\mathrm{Protection}(0) = \mathrm{Premium}(0),
\]
soit :
\[
(1 - R)\int_0^T P(0,t)\, \lambda(t)\, e^{-\int_0^t \lambda(u)\,du}\, dt
=
s
\sum_{k=1}^{n}
\alpha_k P(0,T_k) e^{-\int_0^{T_k} \lambda(u)\,du}.
\]

Dans le cas stationnaire ($\lambda$ constant, taux fixe $r$) :
\[
s
=
\frac{ (1-R) \lambda }{ \sum_{k=1}^{n} \alpha_k e^{-(r+\lambda)T_k} }.
\]

En pratique, les spreads de CDS observés sur le marché permettent d’en déduire
l’intensit\'e $\lambda(t)$ via une calibration similaire à celle de la courbe
des taux.

%%%%%%%%%%%%%%%%%%%%%%%%%%%%%%%%%%%%%%%%%%%%%%%%%%%%%%%%%%%%%%
\subsubsection{Interprétations}

\begin{itemize}
  \item Le spread $s$ augmente lorsque le risque de d\'efaut per\c{c}u augmente.
  \item La probabilit\'e de survie est 
  \[
    \mathbb{Q}(\tau > T) = e^{-\int_0^T \lambda(u)\,du}.
  \]
  \item Le taux de recouvrement $R$ influence directement la jambe de protection.
\end{itemize}

Les CDS sont essentiels pour :
\begin{itemize}
  \item la gestion du risque de cr\'edit,
  \item la construction de courbes de spread,
  \item la valorisation de produits hybrides (equity--credit, FX--credit).
\end{itemize}


%%%%%%%%%%%%%%%%%%%%%%%%%%%%%%%%%%%%%%%%%%%%%%%%%%%%%%%%%%%%%%
\fi
\subsection{Mod\`eles de volatilit\'e avanc\'es : rough, sauts et filtrage}

Cette section a pour objectif d’esquisser quelques directions modernes en finance
quantitative, sans chercher \`a d\'evelopper les d\'erivations compl\`etes, qui sont
report\'ees en annexe.

%%%%%%%%%%%%%%%%%%%%%%%%%%%%%%%%%%%%%%%%%%%%%%%%%%%%%%%%%%%%%%
\subsubsection{Rough volatility et SDE de Volterra}

Les mod\`eles de \emph{volatilit\'e rugueuse} constituent aujourd’hui l’un des cadres
les plus avancés pour reproduire les smiles de volatilité observés sur les
marchés actions. Ils reposent sur des SDE de type Volterra dont la dynamique
est non markovienne :
\[
dX_t
=
\left(
 \int_0^t K(t-s) g(X_s)\, ds
\right) dt
+
\left(
 \int_0^t K(t-s) h(X_s)\, ds
\right) dW_t,
\]
o\`u $K$ est un noyau singulier près de $0$.

Dans le cas rough (Hurst $H<1/2$), la variance est souvent modélisée via un
fractional Brownian motion :
\[
v_t = 
\xi_0(t)\,
\exp\!\left(
\eta B_t^H - \tfrac12 \eta^2 t^{2H}
\right).
\]

Ces mod\`eles permettent :
\begin{itemize}
  \item de reproduire un skew de court terme marqu\'e (stylized fact largement document\'e),
  \item de d\'ecrire une structure de corr\'elation temporelle cohérente,
  \item d’expliquer certains ph\'enom\`enes de la microstructure.
\end{itemize}

\emph{La d\'erivation compl\`ete des SDE de Volterra, du fBM, et du mod\`ele rBergomi
est fournie en annexe.}

%%%%%%%%%%%%%%%%%%%%%%%%%%%%%%%%%%%%%%%%%%%%%%%%%%%%%%%%%%%%%%
\subsubsection{Mod\`ele de Heston et variants affines}

Le mod\`ele de Heston constitue un mod\`ele de r\'ef\'erence en volatilité stochastique.
Il est affine, ce qui permet une expression ferm\'ee de la fonction caract\'eristique :
\[
\phi(u;T) = 
\exp \bigl(C(u,T) + D(u,T) v_0 + i u \ln S_0 \bigr).
\]

Les variants (Heston à corrélation locale, Heston multifactor, hybrid Heston--Hull--White)
permettent :
\begin{itemize}
  \item une meilleure reproduction de la term-structure du skew,
  \item la prise en compte de facteurs de taux ou de FX,
  \item une calibration robuste via transformée de Fourier.
\end{itemize}

Les d\'erivations Riccati compl\`etes (équations pour $C$ et $D$) sont report\'ees en
annexe.

%%%%%%%%%%%%%%%%%%%%%%%%%%%%%%%%%%%%%%%%%%%%%%%%%%%%%%%%%%%%%%
\subsubsection{Mod\`ele SABR et volatilité inspirée (SVI)}

Pour les marchés de taux, le mod\`ele SABR, caractérisé par
\[
dF_t = \alpha_t F_t^\beta\, dW_t^{(1)}, 
\qquad
d\alpha_t = \nu \alpha_t\, dW_t^{(2)},
\]
remains l’un des standards de calibration en raison de son expression 
asymptotique fermée de la volatilité implicite.

En parall\`ele, les mod\`eles SVI (Stochastic Volatility Inspired) fournissent une
paramétrisation analytique de la surface de volatilité, contr\^olant :
\begin{itemize}
  \item le skew,
  \item la convexité,
  \item les conditions d’arbitrage.
\end{itemize}

Ces deux mod\`eles sont des outils incontournables pour reproduire des surfaces
de volatilité de marché de façon flexible et stable.

%%%%%%%%%%%%%%%%%%%%%%%%%%%%%%%%%%%%%%%%%%%%%%%%%%%%%%%%%%%%%%
\subsubsection{Diffusions \`a sauts : Merton et Bates}

\paragraph{Motivation.}
Les smiles et wings de volatilité, en particulier \`a tr\`es courte maturit\'e, traduisent souvent des probabilit\'es
de queues plus \'elev\'ees que celles induites par des diffusions continues. Une mani\`ere classique d'introduire
des queues plus lourdes et des discontinuit\'es est d'ajouter un processus de sauts (compound Poisson).

\paragraph{Mod\`ele de Merton (jump--diffusion).}
Sous $\QQ$ et dans une convention \`a dividende continu $q$, on \'ecrit typiquement :
\[
\frac{\dd S_t}{S_{t^-}}

=
(r-q-\lambda\kappa_J)\,\dd t
+
\sigma\,\dd W_t
+
(J-1)\,\dd N_t,
\]
o\`u $N_t$ est un Poisson d'intensit\'e $\lambda$, $J>0$ le multiplicateur de saut, et
$\kappa_J := \E[J-1]$ (sous $\QQ$) assure la correction de drift (martingalit\'e de $S_t\,\ee^{-\int_0^t (r_u-q_u)\dd u}$).
Un choix standard est $\ln J \sim \mathcal{N}(\mu_J,\delta_J^2)$.

\paragraph{Cons\'equence sur le pricing.}
Le log-prix est somme d'une partie gaussienne et d'une partie compound Poisson ; la fonction caract\'eristique
de $\ln S_T$ s'\'ecrit explicitement, ce qui rend les m\'ethodes de Fourier naturelles. Une repr\'esentation classique
exprime aussi le prix comme une \emph{moyenne} (sur le nombre de sauts) de prix Black--Scholes \`a param\`etres ajust\'es,
ce qui illustre l'impact des sauts sur la courbure des wings.

\paragraph{Mod\`ele de Bates (Heston + sauts).}
Le mod\`ele de Bates combine volatilité stochastique (Heston) et sauts :
\[
\frac{\dd S_t}{S_{t^-}}
= (r-q-\lambda\kappa_J)\,\dd t + \sqrt{v_t}\,\dd W_t^{(1)} + (J-1)\,\dd N_t,
\qquad
\dd v_t = \kappa(\theta-v_t)\,\dd t + \xi\sqrt{v_t}\,\dd W_t^{(2)},
\]
avec $\mathrm{corr}(\dd W_t^{(1)},\dd W_t^{(2)})=\rho\,\dd t$.
Intuitivement, $v_t$ explique la persistance de volatilité et le skew ``diffusif'', tandis que les sauts pilotent
les queues et l'accentuation du smile \`a court terme.

\paragraph{Calibration et limites.}
La calibration vise \`a identifier $(\lambda,\mu_J,\delta_J)$ conjointement aux param\`etres de volatilité.
En pratique, l'identifiabilit\'e est d\'elicate : certaines combinaisons de $(\xi,\rho,v_0)$ et de param\`etres de sauts
peuvent produire des smiles proches. Il est donc naturel d'utiliser des pond\'erations (vega), des contraintes
(bornes, conditions de positivité) et des diagnostics de stabilit\'e (sensibilit\'e, out-of-sample).

%%%%%%%%%%%%%%%%%%%%%%%%%%%%%%%%%%%%%%%%%%%%%%%%%%%%%%%%%%%%%%
\subsubsection{Filtrage : mod\`eles en espace d'\'etat et filtre de Kalman}

\paragraph{Id\'ee.}
De nombreuses quantit\'es pertinentes (volatilité instantan\'ee, variance latente, niveau de bruit) ne sont pas observables
directement. Les \emph{mod\`eles en espace d'\'etat} s\'eparent une dynamique latente (\'etat) et une \'equation d'observation
(mesures bruit\'ees), puis utilisent un filtre pour estimer l'\'etat conditionnellement aux observations.

\paragraph{Cadre lin\'eaire gaussien.}
Dans sa forme standard, on postule :
\[
x_{t+1} = A x_t + w_t,
\qquad
y_t = H x_t + v_t,
\]
o\`u $(w_t)$ et $(v_t)$ sont gaussiens centr\'es, de covariances $Q$ et $R$, ind\'ependants. Le \textbf{filtre de Kalman}
fournit alors l'estimateur optimal (au sens des moindres carr\'es) de $x_t$ sachant $(y_0,\dots,y_t)$, avec une r\'ecurrence
sur la moyenne conditionnelle et sa covariance.

\paragraph{Extensions et limites.}
Pour des dynamiques non lin\'eaires (volatilité stochastique, contraintes de positivité), on recourt \`a des extensions
(EKF/UKF) ou \`a des filtres particulaires. Les limites principales viennent de la non-gaussianit\'e (queues, sauts),
de la non-stationnarit\'e (r\'egimes), et du risque de mod\`ele : un filtre est un outil d'estimation conditionnelle,
pas une preuve de validit\'e du mod\`ele.

\paragraph{Lien avec le projet \emph{PaperTradingApp}.}
Un filtrage de type Kalman s'interpr\`ete comme une brique th\'eorique d'estimation d'\'etats latents (variance/volatilité)
et de lissage, utilisable pour : (i) produire des s\'eries de volatilité plus robustes que des estimateurs bruts,
(ii) fournir des entr\'ees d'\'etat pour des analyses de r\'egimes ou des politiques de couverture, (iii) quantifier
une incertitude d'estimation via la covariance filtr\'ee.

%%%%%%%%%%%%%%%%%%%%%%%%%%%%%%%%%%%%%%%%%%%%%%%%%%%%%%%%%%%%%%
\subsubsection{Perspectives}

Les mod\`eles modernes combinent :
\begin{itemize}
  \item volatilité stochastique,
  \item volatilité rugueuse,
  \item noyaux de Volterra,
  \item dynamiques multi-actifs,
  \item facteurs de taux ou de crédit.
\end{itemize}

La recherche actuelle explore des structures hybrides (Heston--Volterra,
rough Heston, SABR--Volterra) et des méthodes numériques accélérées (FFT,
hybrid schemes, quadratures rapides).

Ces diff\'erentes directions visent \`a construire des mod\`eles plus expressifs, susceptibles de mieux prendre en compte
plusieurs effets empiriques : skew dynamique, d\'eformations fines de smile, et ph\'enom\`enes de m\'emoire.





%%%%%%%%%%%%%%%%%%%%%%%%%%%%%%%%%%%%%%%%%%%%%%%%%%%%%%%%%%%%%%
\subsubsection{Lien avec le projet \emph{PaperTradingApp}}

Cette section fournit le panorama des mod\`eles effectivement pertinents pour une plateforme d'options : (i) mod\`eles affines (Heston) pour
leur tractabilit\'e via fonctions caract\'eristiques, (ii) mod\`eles \`a sauts (Merton/Bates) pour traiter des queues et de la dynamique court terme,
(iii) mod\`eles \`a m\'emoire (Volterra/rough) pour la g\'eom\'etrie du smile \`a court terme, et (iv) outils d'estimation (Kalman) pour relier variables
latentes et observations bruit\'ees. Le projet vise \`a comparer ces cadres, \`a discuter leurs hypoth\`eses et limites, et \`a articuler mod\`eles et
surfaces de march\'e dans une logique de calibration et de gestion du risque.

\subsection{Gestion du risque de portefeuille et allocation}

La valorisation (\emph{pricing}) et la mod\'elisation de la volatilit\'e n'ont de sens op\'erationnel que si elles
s'int\`egrent dans un cadre de gestion du risque. Dans une plateforme de trading, ce cadre se traduit par :
(i) une mesure claire des expositions, (ii) une d\'ecomposition du PnL, (iii) des limites et alertes, et
(iv) des r\`egles d'allocation et de r\'e\'equilibrage coh\'erentes avec le profil de risque.

%%%%%%%%%%%%%%%%%%%%%%%%%%%%%%%%%%%%%%%%%%%%%%%%%%%%%%%%%%%%%%
\subsubsection{PnL, expositions et drawdowns}

On distingue typiquement :
\begin{itemize}
  \item le \textbf{PnL r\'ealis\'e} (positions ferm\'ees) et le \textbf{PnL latent} (mark-to-market),
  \item l'\textbf{exposition brute} (somme des valeurs absolues) et l'\textbf{exposition nette} (somme sign\'ee),
  \item des sensibilit\'es de premier ordre (\textbf{delta}) et de second ordre (\textbf{gamma}) lorsque le portefeuille
        contient des options.
\end{itemize}

Le \textbf{drawdown} mesure la perte relative depuis un plus-haut historique :
\[
\mathrm{DD}(t)=\frac{V_t-\max_{u\le t}V_u}{\max_{u\le t}V_u},
\qquad
\mathrm{MDD}=\min_t \mathrm{DD}(t),
\]
où $V_t$ est la valeur du portefeuille. Le drawdown est une mesure intuitive de risque ``v\'ecu'' et compl\`ete utilement
des mesures bas\'ees sur la variance.

%%%%%%%%%%%%%%%%%%%%%%%%%%%%%%%%%%%%%%%%%%%%%%%%%%%%%%%%%%%%%%
\subsubsection{Value-at-Risk (VaR) et Expected Shortfall (ES)}

Pour un horizon $\Delta t$ et un niveau de confiance $\alpha\in(0,1)$, la \textbf{VaR} est d\'efinie comme un quantile
de la distribution des pertes $L$ :
\[
\mathrm{VaR}_\alpha(L)=\inf\{\ell\in\mathbb{R}:\mathbb{P}(L\le \ell)\ge \alpha\}.
\]

L'\textbf{Expected Shortfall} (ES) ou \emph{Conditional VaR} est l'esp\'erance des pertes au-del\`a de la VaR :
\[
\mathrm{ES}_\alpha(L)=\mathbb{E}\big[L\mid L\ge \mathrm{VaR}_\alpha(L)\big].
\]

La VaR est simple \`a communiquer mais n'est pas toujours coh\'erente au sens d'Artzner et al. (elle peut violer la
subadditivit\'e). L'ES corrige ce point et est souvent pr\'ef\'er\'ee pour piloter des limites de risque.

%%%%%%%%%%%%%%%%%%%%%%%%%%%%%%%%%%%%%%%%%%%%%%%%%%%%%%%%%%%%%%
\subsubsection{Allocation moyenne--variance (Markowitz)}

Dans le cadre moyenne--variance, on mod\'elise les rendements d'un vecteur d'actifs $R\in\mathbb{R}^n$ par sa moyenne
$\mu$ et sa matrice de covariance $\Sigma$. L'allocation (poids) $w\in\mathbb{R}^n$ est choisie pour :
\begin{itemize}
  \item minimiser la variance pour un rendement cible,
  \item ou maximiser un ratio de Sharpe sous contraintes (poids positifs, somme \`a 1, etc.).
\end{itemize}

Le probl\`eme ``minimum variance'' (sans contrainte de signe) admet une solution ferm\'ee :
\[
w^\star \propto \Sigma^{-1}\mathbf{1},\qquad
\mathbf{1}=(1,\dots,1)^\top.
\]
En pratique, on impose des contraintes r\'ealistes et on r\'egularise $\Sigma$ (estimation bruit\'ee, fen\^etre courte,
actifs corr\'el\'es).

%%%%%%%%%%%%%%%%%%%%%%%%%%%%%%%%%%%%%%%%%%%%%%%%%%%%%%%%%%%%%%
\subsubsection{Risk parity et contributions au risque}

Le \emph{risk parity} vise une allocation o\`u chaque actif contribue \`a parts \'egales au risque du portefeuille.
Pour une volatilité de portefeuille $\sigma_P=\sqrt{w^\top\Sigma w}$, la contribution marginale au risque est
$\partial\sigma_P/\partial w_i = (\Sigma w)_i/\sigma_P$, et la contribution totale de l'actif $i$ est
\[
\mathrm{RC}_i = w_i \frac{(\Sigma w)_i}{\sigma_P}.
\]
Le risk parity cherche $\mathrm{RC}_i=\sigma_P/n$ pour tous $i$ (sous contraintes de signe/somme), ce qui d\'efinit un
probl\`eme non lin\'eaire r\'esolu num\'eriquement.

%%%%%%%%%%%%%%%%%%%%%%%%%%%%%%%%%%%%%%%%%%%%%%%%%%%%%%%%%%%%%%
\subsubsection{PCA et eigen-portfolios}

La d\'ecomposition spectrale $\Sigma = Q\Lambda Q^\top$ (valeurs propres $\Lambda$, vecteurs propres $Q$) fournit une
lecture factorielle des risques : les grandes valeurs propres correspondent aux facteurs dominants. Les
\emph{eigen-portfolios} exploitent ces vecteurs propres pour construire des portefeuilles align\'es sur des facteurs
ou, \`a l'inverse, moins expos\'es au facteur principal.

\subsubsection{Lien avec le projet \emph{PaperTradingApp}}

Le projet mobilise ces notions pour rendre \emph{mesurable} et \emph{attribuable} le risque : expositions brutes/nettes, PnL r\'ealis\'e vs latent,
drawdowns, VaR/ES, et principes d'allocation (Markowitz, risk parity, d\'ecomposition factorielle). L'objectif est de relier les briques de pricing
et de calibration \`a des d\'ecisions de portefeuille (taille de position, limites, r\'e\'equilibrage) dans un langage quantitatif coh\'erent.

\subsection{R\'egimes et analytics de volatilit\'e}

L'analyse de volatilité ne se limite pas \`a une surface implicite. Dans une perspective de gestion du risque et de
d\'ecision, on s'int\'eresse \`a la \textbf{volatilité r\'ealis\'ee}, \`a des diagnostics de \textbf{retour \`a la moyenne},
et \`a des mod\`eles de \textbf{r\'egimes} capables de capturer des changements de dynamique (calme/crise, compression/expansion de vol).

%%%%%%%%%%%%%%%%%%%%%%%%%%%%%%%%%%%%%%%%%%%%%%%%%%%%%%%%%%%%%%
\subsubsection{Volatilit\'e r\'ealis\'ee : d\'efinitions, annualisation, biais}

\paragraph{D\'efinition (variance r\'ealis\'ee).}
Soit une partition $t=t_0<\dots<t_n=t+\Delta$ et des rendements log
$r_i=\ln(S_{t_i}/S_{t_{i-1}})$. La \emph{variance r\'ealis\'ee} sur $[t,t+\Delta]$ est
\[
\mathrm{RV}_{t,t+\Delta}=\sum_{i=1}^n r_i^2.
\]
Sous des hypoth\`eses de diffusion continue, $\mathrm{RV}$ converge (quand la maille tend vers $0$) vers la \emph{variation quadratique}
du log-prix, et constitue un estimateur naturel de la variance int\'egr\'ee.

\paragraph{Annualisation.}
Une volatilité annualis\'ee (sur une fen\^etre) est souvent d\'efinie par $\widehat\sigma=\sqrt{\mathrm{RV}}\times\sqrt{A}$,
o\`u $A$ d\'epend de la convention (nombre de jours de bourse, etc.). Le point cl\'e est que l'annualisation est une convention
et doit \^etre coh\'erente avec les maturit\'es d'options et la courbe de taux.

\paragraph{Biais et pi\`eges.}
En pratique, des ph\'enom\`enes de microstructure (bid--ask bounce, non-synchronisme), des heures de cotation, et des sauts
peuvent biaiser $\mathrm{RV}$. Le choix de la fr\'equence et de la fen\^etre r\'esulte d'un compromis biais/variance ; on peut
compl\'eter par des estimateurs robustes (r\'ealis\'ee bipower, trimming) selon l'objectif.

%%%%%%%%%%%%%%%%%%%%%%%%%%%%%%%%%%%%%%%%%%%%%%%%%%%%%%%%%%%%%%
\subsubsection{Retour \`a la moyenne : diagnostics et interpr\'etation}

\paragraph{Intuition.}
De nombreux indicateurs de volatilité (implicite/realized) montrent une tendance \`a revenir vers un niveau ``normal'' \`a horizon
de quelques semaines \`a quelques mois. Ce ph\'enom\`ene motive des dynamiques de type OU/CIR pour $v_t$ ou $\sigma_t$.

\paragraph{R\'esultat cl\'e (OU discretis\'e).}
Pour un processus d'Ornstein--Uhlenbeck $\dd x_t=\kappa(\theta-x_t)\dd t + \eta\,\dd W_t$, la discr\'etisation \`a pas $\Delta$ donne :
\[
x_{t+\Delta}-\theta = \ee^{-\kappa\Delta}(x_t-\theta) + \text{bruit}.
\]
On en d\'eduit une notion de \emph{demi-vie} $t_{1/2}=\ln 2/\kappa$ : horizon \`a partir duquel l'\'ecart \`a $\theta$ est divis\'e par deux en moyenne.

\paragraph{Limites.}
Les diagnostics de mean reversion sont sensibles \`a la non-stationnarit\'e, aux changements de r\'egime, et \`a la d\'efinition
de la volatilité (r\'ealis\'ee, implicite, filtr\'ee). Une conclusion robuste exige des tests sur plusieurs horizons, et des stress
sur les choix de fen\^etre et d'\'echantillonnage.

%%%%%%%%%%%%%%%%%%%%%%%%%%%%%%%%%%%%%%%%%%%%%%%%%%%%%%%%%%%%%%
\subsubsection{Cha\^ines de Markov de r\'egimes : matrice de transition et lecture probabiliste}

\paragraph{D\'efinition.}
On introduit une variable latente $Z_t\in\{1,\dots,m\}$ (r\'egime) telle que
\[
\PP(Z_{t+\Delta}=j \mid Z_t=i)=\Pi_{ij},
\]
o\`u $\Pi$ est une matrice de transition (lignes sommant \`a $1$). Conditionnellement au r\'egime, les observations (rendements, volatilité)
suivent une loi param\'etr\'ee (par exemple gaussienne avec variance r\'egime-d\'ependante).

\paragraph{R\'esultat cl\'e (stationnarit\'e).}
Si la cha\^ine est ergodique, il existe une distribution stationnaire $\pi$ telle que $\pi^\top \Pi=\pi^\top$ ; elle fournit la proportion
de temps de long terme pass\'ee dans chaque r\'egime (interpr\'etation probabiliste du ``climat'' de march\'e).

\paragraph{Limites.}
Les r\'egimes sont une \emph{discr\'etisation} : le choix de $m$ et de l'\'equation d'observation peut induire de l'instabilit\'e et de la
sur-param\'etrisation. La validation doit inclure des tests hors \'echantillon, des contraintes d'interpr\'etabilit\'e (r\'egimes stables),
et des analyses de sensibilit\'e aux donn\'ees.

%%%%%%%%%%%%%%%%%%%%%%%%%%%%%%%%%%%%%%%%%%%%%%%%%%%%%%%%%%%%%%
\subsubsection{Lien avec le projet \emph{PaperTradingApp}}

Ces outils fournissent un langage commun pour : (i) contextualiser les smiles (r\'egimes de stress vs de compression),
(ii) moduler des limites de risque (VaR/ES, drawdowns) en fonction d'un r\'egime, (iii) proposer des diagnostics de mean reversion sur la
volatilité r\'ealis\'ee/implicite, et (iv) d\'efinir des \'etats synth\'etiques utilisables dans une formulation MDP de la couverture
(\`a d\'efaut d'hypoth\`eses stationnaires fortes).

%==========================================================
%                 PARTIE 2 : APPRENTISSAGE \& D\'ECISION
%==========================================================
\section{Apprentissage supervis\'e sous non-stationnarit\'e : validation, leakage et robustesse}

Cette partie introduit des notions de \emph{m\'ethodologie quantitative} issues de l'apprentissage supervis\'e, particuli\`erement utiles
lorsque les donn\'ees sont structur\'ees dans le temps et potentiellement \emph{non-stationnaires}. Dans un contexte de pr\'ediction en production
ou de comp\'etition structur\'ee en p\'eriodes (\emph{eras}), les enjeux ne sont pas uniquement la performance moyenne, mais la \textbf{stabilit\'e}
du signal et la \textbf{pr\'evention des biais} (en particulier le \emph{leakage}).

\subsection{Cadre : apprentissage supervis\'e et g\'en\'eralisation}

\subsubsection{D\'efinitions}
On observe un ensemble annot\'e $\{(x_i,y_i)\}_{i=1}^n$ o\`u $x_i\in\mathcal{X}$ (vecteur de variables explicatives) et $y_i\in\mathcal{Y}$ (cible).
Un mod\`ele param\'etr\'e $f_\theta:\mathcal{X}\to\mathcal{Y}$ est ajust\'e en minimisant une perte empirique
\[
\widehat\theta\in\argmin_{\theta}\ \frac{1}{n}\sum_{i=1}^n \ell\bigl(y_i,f_\theta(x_i)\bigr).
\]
La question centrale est la \textbf{g\'en\'eralisation} : les performances sur des donn\'ees futures (ou hors \'echantillon) seront-elles coh\'erentes
avec celles observ\'ees sur l'\'echantillon d'apprentissage ?

\subsubsection{Intuition : biais--variance et non-stationnarit\'e}
Le compromis biais--variance fournit une lecture utile : un mod\`ele trop flexible sur-ajuste le bruit (variance forte), un mod\`ele trop rigide
manque le signal (biais fort). La non-stationnarit\'e (changement de r\'egime, d\'erive de distribution) complique encore l'analyse : la distribution
des $(x,y)$ peut varier avec le temps, ce qui rend les estimations i.i.d.\ (\"echantillon al\'eatoire) inad\'equates.

\subsection{Eras, r\'egimes et \emph{leakage} : risques sp\'ecifiques}

\subsubsection{D\'efinition d'une structure en eras}
On suppose que les observations sont regroup\'ees par p\'eriodes $e\in\{1,\dots,E\}$ (``eras''), chaque era correspondant \`a un bloc temporel
ou \`a un sous-ensemble relativement homog\`ene. On note $\mathcal{I}_e$ l'ensemble des indices appartenant \`a l'era $e$.
Cette structure impose une \'evaluation \emph{par p\'eriode} : un mod\`ele utile doit produire un signal \textbf{r\'ep\'etable} au fil des eras.

\subsubsection{Leakage : d\'efinition et cons\'equences}
On appelle \emph{leakage} toute utilisation, directe ou indirecte, d'une information indisponible au moment r\'eel de la pr\'ediction
(information ``future'', fuite entre train et validation, pr\'etraitements calcul\'es sur un ensemble trop large).

\paragraph{Cons\'equence.}
Le leakage conduit \`a une performance hors \'echantillon artificiellement optimiste : l'\'evaluation mesure alors une capacit\'e \`a exploiter une fuite,
pas \`a g\'en\'eraliser. Dans les donn\'ees temporelles, le leakage est souvent plus subtil qu'un simple doublon : il peut provenir d'un choix de split,
d'une s\'election de variables, ou d'une normalisation utilisant des statistiques calcul\'ees ``dans le futur''.

\subsubsection{R\'esultat cl\'e (principe de validation temporelle)}
Pour obtenir une estimation non biais\'ee de performance, la validation doit respecter la causalit\'e : les ensembles d'entra\^{i}nement doivent \^etre
ant\'erieurs (ou, \`a minima, disjoints et sans information partag\'ee) par rapport aux ensembles de validation. Ce principe se formalise par des d\'ecoupages
par blocs temporels :
\begin{itemize}
  \item \textbf{fen\^etre croissante (expanding window)} : on entra\^{i}ne sur le pass\'e, on valide sur un bloc futur ;
  \item \textbf{fen\^etre glissante (rolling window)} : on entra\^{i}ne sur une fen\^etre r\'ecente, on valide juste apr\`es.
\end{itemize}
Le choix refl\`ete un compromis : stabilit\'e (beaucoup de donn\'ees) versus pertinence (r\'ecence du r\'egime).

\paragraph{Holdout final et risque d'optimisation.}
En pratique, on r\'eserve un dernier bloc temporel (\emph{holdout}) comme \textbf{test final}, consult\'e le plus tard possible.
Le consulter it\'erativement pour choisir variables ou hyperparam\`etres le transforme en ensemble de d\'eveloppement et r\'eintroduit un biais
(\"optimisation sur le test\"). Une validation prudente multiplie les fen\^etres et contr\^ole l'\'ecart entre performances sur CV temporelle et sur holdout.

\subsection{M\'etriques de performance par era : corr\'elation et stabilit\'e}

\subsubsection{Corr\'elation de Pearson par era}
Une m\'etrique fr\'equente dans les probl\`emes de pr\'ediction ``faible signal'' est la corr\'elation entre pr\'ediction $\widehat y$ et cible $y$,
calcul\'ee \emph{au sein de chaque era}. Pour une era $e$ de taille $n_e$, avec $\{\widehat y_i\}_{i\in\mathcal{I}_e}$ et $\{y_i\}_{i\in\mathcal{I}_e}$,
\[
\rho_e(\widehat y,y)=
\frac{\sum_{i\in\mathcal{I}_e}(\widehat y_i-\overline{\widehat y}_e)(y_i-\overline y_e)}
{\sqrt{\sum_{i\in\mathcal{I}_e}(\widehat y_i-\overline{\widehat y}_e)^2}\sqrt{\sum_{i\in\mathcal{I}_e}(y_i-\overline y_e)^2}}.
\]
On obtient une s\'erie $\{\rho_e\}_{e=1}^E$, dont l'analyse renseigne sur la stabilit\'e inter-eras.

\subsubsection{Agr\'egation : moyenne, dispersion, ratio type Sharpe}
On r\'esume souvent la s\'erie par sa moyenne et sa dispersion :
\[
\mu = \frac{1}{E}\sum_{e=1}^E\rho_e,
\qquad
\sigma = \sqrt{\frac{1}{E-1}\sum_{e=1}^E(\rho_e-\mu)^2},
\qquad
\mathrm{SR}_{\rho}=\frac{\mu}{\sigma}.
\]
Le ratio $\mathrm{SR}_{\rho}$ (analogue \`a un Sharpe) mesure un compromis ``niveau'' versus ``stabilit\'e'' : un signal faible mais r\'egulier peut \^etre
pr\'ef\'erable \`a un signal moyen plus \'elev\'e mais instable.

\subsubsection{Limites}
La corr\'elation de Pearson est sensible aux non-lin\'earit\'es et aux queues ; elle peut varier fortement si la taille $n_e$ est faible.
Une validation robuste doit donc (i) analyser les distributions de $\rho_e$ (asym\'etries, outliers), (ii) r\'ep\'eter sur plusieurs p\'eriodes,
et (iii) compl\'eter par des tests de stress (changements de r\'egime, d\'erive de features).

\subsection{R\'eduction de dimension et s\'election de variables}

\subsubsection{Filtre de valeurs manquantes}
Pour une variable $f$, on note le ratio de valeurs manquantes
\[
\mathrm{NA}(f)=\frac{\#\{i: x_{i,f}\ \text{manquant}\}}{n}.
\]
Un filtre $\mathrm{NA}(f)\le \tau$ est un contr\^ole simple visant \`a supprimer des variables trop instables.

\subsubsection{Score univari\'e par association \`a la cible}
Une s\'election univari\'ee consiste \`a calculer un score $s(f)$ (corr\'elation, information mutuelle approxim\'ee, etc.) et \`a retenir les top-$K$
variables. L'intuition est de limiter le bruit et le co\^ut de calcul en conservant des variables portant un signal ``individuel''.

\paragraph{Limites.}
Un filtre univari\'e ignore les interactions et la redondance entre variables ; il peut aussi sur-ajuster si le score est calcul\'e en ``regardant''
la validation (leakage de s\'election). Une pratique robuste consiste \`a figer la s\'election sur une p\'eriode d'entra\^{i}nement, ou \`a la r\'e\'evaluer
par fold de validation temporelle.

\subsection{Mod\`eles et r\'egularisation (vue th\'eorique)}

\subsubsection{Gradient boosting sur arbres (GBDT)}
Le gradient boosting construit un mod\`ele additif
\[
\widehat y(x)=\sum_{t=1}^T \eta\, f_t(x),
\]
o\`u chaque $f_t$ est un arbre de d\'ecision (faible apprenant) et $\eta$ un pas d'apprentissage. L'id\'ee est d'ajouter it\'erativement des arbres
qui corrigent les erreurs des pr\'edictions pr\'ec\'edentes (descente de gradient dans l'espace des fonctions).
La complexit\'e est contr\^ol\'ee par des param\`etres de profondeur/feuilles, des \'echantillonnages (lignes/variables) et de la r\'egularisation.

\paragraph{Arr\^et anticip\'e (early stopping).}
L'arr\^et anticip\'e consiste \`a interrompre l'apprentissage lorsque la performance sur validation n'am\'eliore plus : il agit comme une forme
de r\'egularisation (on \'evite d'ajouter des composantes qui n'apprennent que le bruit) et r\'eduit le co\^ut de calcul.
Dans un cadre structur\'e en eras, il est coh\'erent de surveiller une m\'etrique align\'ee avec l'objectif de stabilit\'e (par exemple une corr\'elation agr\'eg\'ee par eras).

\subsubsection{Ridge (r\'egression lin\'eaire r\'egularis\'ee)}
La Ridge minimise
\[
\min_{w\in\R^p}\ \frac{1}{n}\|y-Xw\|_2^2 + \alpha\|w\|_2^2,
\]
avec $\alpha\ge 0$. Elle est utile comme baseline stable, et comme ``meta-mod\`ele'' pour combiner des pr\'edictions (stacking) lorsque la validation est stricte.

\subsubsection{R\'eseaux (MLP) : principe et risques}
Un perceptron multi-couches (MLP) approxime une fonction non-lin\'eaire par compositions affines + activations. Il est expressif, mais sensible \`a la
normalisation, \`a la r\'egularisation, et au risque de sur-apprentissage si le protocole temporel n'est pas rigoureux.

\subsection{Ensembles, OOF et post-traitements monotones}

\subsubsection{Ensembles : r\'eduction de variance}
L'agr\'egation de plusieurs mod\`eles ``semblables'' (initialisations diff\'erentes, sous-\'echantillonnages) par moyenne
\[
\widehat y^{ens}(x)=\frac{1}{M}\sum_{m=1}^M \widehat y^{(m)}(x)
\]
r\'eduit typiquement la variance et stabilise la performance par era, au prix d'un co\^ut de calcul plus \'elev\'e.

\subsubsection{OOF et stacking : prudence m\'ethodologique}
Le stacking consiste \`a apprendre un combinateur sur des pr\'edictions \emph{out-of-fold} (OOF), produites par des mod\`eles de base sur des donn\'ees non vues
lors de leur entra\^{i}nement. Le point critique est que le meta-mod\`ele doit \^etre valid\'e avec le m\^eme soin temporel, faute de quoi il sur-ajuste
la validation.

\subsubsection{Ranking (transform\'ee par CDF empirique)}
Un post-traitement fr\'equent est la transformation par rang (CDF empirique) : si $r_i$ est le rang de $\widehat y_i$ parmi $N$ observations,
\[
\widehat y_i'=\frac{r_i}{N}.
\]
Cette transformation impose une distribution marginale uniforme (stabilit\'e d'\'echelle). Elle pr\'eserve l'ordre (donc la corr\'elation de rang),
mais modifie en g\'en\'eral la corr\'elation lin\'eaire ; son int\'er\^et est surtout de limiter des effets de d\'erive de distribution et d'outliers.

\subsubsection{Contr\^oles de coh\'erence (sanity checks)}
Avant d'utiliser un mod\`ele pour produire des d\'ecisions, des contr\^oles de coh\'erence limitent les erreurs grossi\`eres : (i) absence de valeurs manquantes
dans les sorties, (ii) coh\'erence des dimensions (m\^emes variables qu'en entra\^{i}nement), (iii) d\'etection d'effondrement de distribution (scores constants,
valeurs extr\^emes), et (iv) tra\c{c}abilit\'e des versions de donn\'ees et d'hypoth\`eses.

\subsection{Lien avec le projet \emph{PaperTradingApp}}

Ces notions apparaissent dans le projet \`a deux niveaux : (i) comme \textbf{m\'ethodologie} de validation pour toute brique ``apprenante''
(surrogates de pricing/calibration, classification de r\'egimes, diagnostics de volatilité), et (ii) comme \textbf{garde-fous} contre l'illusion de
performance (leakage temporel, sur-ajustement d'hyperparam\`etres, d\'erive de donn\'ees). Elles compl\`etent les chapitres de finance quantitative :
un bon mod\`ele de smile ou de r\'egime n'est pas seulement expressif, il doit aussi \^etre \emph{robuste} et \emph{test\'e} hors \'echantillon selon
un protocole coh\'erent avec la temporalit\'e des donn\'ees de march\'e.

\section{D\'ecision s\'equentielle : MDP et apprentissage par renforcement}

%----------------------------------------------------------
\iffalse % Hors p\'erim\`etre : ML g\'en\'eral (conserv\'e hors rendu final)
\subsection{Supervised Learning}

Le \emph{supervised learning} consiste à apprendre une fonction
\[
f_\theta : \mathcal{X} \to \mathcal{Y}
\]
à partir d’un ensemble d’observations annotées
$\{(x_i,y_i)\}_{i=1}^n$.  
L’objectif est de déterminer les paramètres $\theta$ minimisant une fonction de
perte :
\[
\theta^\*
=
\arg\min_\theta 
\sum_{i=1}^{n}
L\bigl(y_i, f_\theta(x_i)\bigr).
\]
Cette approche regroupe la régression, la classification, les modèles linéaires,
les modèles arborescents, les SVM, ainsi que les méthodes probabilistes
(Gaussian Processes).

%%%%%%%%%%%%%%%%%%%%%%%%%%%%%%%%%%%%%%%%%%%%%%%%%%%%%%%%%%%%%%
\subsubsection{Fondamentaux}

\paragraph{Tâches usuelles.}

\begin{itemize}
  \item \textbf{Régression} : prédiction de variables continues ($y \in \mathbb{R}$).
  \item \textbf{Classification} : prédiction de labels discrets 
        ($y \in \{1,\dots,K\}$).
\end{itemize}

\paragraph{Séparation du dataset.}

Les données sont divisées en :
\[
\text{train} \longrightarrow \text{validation} \longrightarrow \text{test}.
\]
Cela permet d’évaluer la capacité de généralisation du modèle.

\paragraph{Overfitting / Underfitting.}

\begin{itemize}
  \item \textbf{Overfitting} : modèle trop complexe, apprend le bruit.
  \item \textbf{Underfitting} : modèle trop simple, structure insuffisante.
\end{itemize}

\paragraph{Biais--variance.}

Le risque prédit se décompose en :
\[
\text{Erreur} = \text{Biais}^2 + \text{Variance} + \text{Bruit}.
\]
L’objectif est de trouver un compromis optimal entre biais et variance.

\paragraph{Fonctions de perte.}

\begin{itemize}
  \item Régression : 
  \[
    \text{MSE},\quad \text{MAE},\quad \text{Huber}.
  \]
  \item Classification :
  \[
    \text{cross-entropy},\quad \text{log-loss}.
  \]
\end{itemize}

\paragraph{Optimisation.}

La mise à jour des paramètres suit la descente de gradient :
\[
\theta \leftarrow \theta - \eta \nabla_\theta L.
\]
Avec variantes :
\[
\text{Momentum},\quad \text{Nesterov},\quad \text{Adam},\quad \text{AdamW},
\quad \text{LR Scheduling}.
\]

%%%%%%%%%%%%%%%%%%%%%%%%%%%%%%%%%%%%%%%%%%%%%%%%%%%%%%%%%%%%%%
\subsubsection{Modèles linéaires \& factoriels}

\paragraph{Régression linéaire.}
\[
\hat y = x^\top w.
\]

\paragraph{Régression logistique.}
\[
p(y=1|x) = \sigma(w^\top x),
\qquad
\sigma(z) = \frac{1}{1+e^{-z}}.
\]

\paragraph{PCA (réduction dimensionnelle).}
On cherche les vecteurs propres maximisant :
\[
\max_v v^\top \Sigma v,
\]
où $\Sigma$ est la matrice de covariance.

\paragraph{SVD.}
Décomposition :
\[
X = U \Sigma V^\top,
\]
fondamentale pour la factorisation matricielle et certains algorithmes de ML.

\paragraph{Gaussian Processes (GP).}
Modèle bayésien non-paramétrique :
\[
f(x) \sim \mathcal{GP}(m(x), k(x,x')),
\]
où $k$ est un noyau (RBF, Matérn, etc.).
Prédiction :
\[
f(x_*) | X,y \sim \mathcal{N}(\mu_*, \sigma_*^2).
\]

\paragraph{Bases de Bernstein.}
Bases polynomiales utilisées pour approximer des fonctions :
\[
B_{k,n}(x) = \binom{n}{k} x^k (1-x)^{n-k}.
\]

\paragraph{Régularisation.}
\begin{itemize}
  \item Ridge : $\|w\|_2^2$
  \item Lasso : $\|w\|_1$
  \item ElasticNet : combinaison des deux
\end{itemize}

%%%%%%%%%%%%%%%%%%%%%%%%%%%%%%%%%%%%%%%%%%%%%%%%%%%%%%%%%%%%%%
\subsubsection{Arbres \& Ensembles}

\paragraph{Arbres de décision.}

Un arbre partitionne l’espace des features :
\[
x \mapsto \text{Région}_i \mapsto \hat y_i.
\]

\paragraph{Bagging.}
Bootstrap + vote majoritaire :
\[
f(x) = \mathrm{Majority}(h_1(x),\dots,h_T(x)).
\]

\paragraph{Random Forest.}
Bagging + sélection aléatoire de variables à chaque split.

\paragraph{Boosting.}
Construction séquentielle :
\[
f_T(x) = f_{T-1}(x) + \eta h_T(x),
\]
où $h_T$ corrige les erreurs du modèle précédent.

Méthodes principales : AdaBoost, GradientBoosting, XGBoost, LightGBM, CatBoost.

\paragraph{Stacking.}
Combinaison de plusieurs modèles via un méta-modèle (logistique, réseau
neural, ridge).

%%%%%%%%%%%%%%%%%%%%%%%%%%%%%%%%%%%%%%%%%%%%%%%%%%%%%%%%%%%%%%
\subsection{Unsupervised Learning}

L’apprentissage non supervisé vise à identifier des structures, regrouper des
observations ou réduire la dimensionnalité d’un jeu de données sans disposer
de labels explicites. Ces méthodes sont essentielles pour la détection de
régimes, la segmentation de marché, la compression de données et la découverte
de facteurs latents.

%%%%%%%%%%%%%%%%%%%%%%%%%%%%%%%%%%%%%%%%%%%%%%%%%%%%%%%%%%%%%%
\subsubsection{Clustering}

\paragraph{K-Means.}

Le K-Means vise à partitionner un ensemble de points 
$\{x_i\}_{i=1}^n \subset \mathbb{R}^d$ 
en $K$ groupes en minimisant l’inertie intra-cluster :
\[
\min_{C_1,\dots,C_K}
\sum_{k=1}^K \sum_{x_i \in C_k} \|x_i - \mu_k\|^2,
\]
o\`u $\mu_k$ est le barycentre du cluster $C_k$.

L’algorithme alterne entre :
\begin{itemize}
  \item \textit{assignation} des points aux centres les plus proches,
  \item \textit{mise à jour} des centres.
\end{itemize}

La convergence est rapide mais dépend de l’initialisation.

%%%%%%%%%%%%%%%%%%%%%%%%%%%%%%%%%%%%%%%%%%%%%%%%%%%%%%%%%%%%%%
\subsubsection{Anomaly / Novelty Detection}

\paragraph{One-Class SVM.}

Le One-Class SVM estime la frontière d’un ensemble de données non étiquetées.
On cherche une fonction $f$ telle que la majorité des points vérifient $f(x) \ge 0$.

Formulation :
\[
\min_{w,\rho,\xi_i}
\frac12 \|w\|^2 + \frac{1}{\nu n} \sum_i \xi_i - \rho
\quad
\text{s.c.}
\quad
w^\top \phi(x_i) \ge \rho - \xi_i.
\]

Les anomalies sont les points pour lesquels $f(x)<0$.

%%%%%%%%%%%%%%%%%%%%%%%%%%%%%%%%%%%%%%%%%%%%%%%%%%%%%%%%%%%%%%
\subsubsection{Réduction dimensionnelle}

\paragraph{PCA (Principal Component Analysis).}

On diagonalise la matrice de covariance :
\[
\Sigma = \frac{1}{n}\sum_{i=1}^n (x_i - \bar x)(x_i - \bar x)^\top
= Q \Lambda Q^\top.
\]

Les vecteurs propres $q_k$ correspondant aux plus grandes valeurs propres 
capturent les directions de variance maximale.

\paragraph{SVD (Singular Value Decomposition).}

Pour une matrice $X \in \mathbb{R}^{n \times d}$ :
\[
X = U \Sigma V^\top.
\]

La SVD est la base de :
\begin{itemize}
  \item la compression,
  \item la factorisation,
  \item la résolution de systèmes linéaires,
  \item les méthodes de recommandation.
\end{itemize}

\paragraph{t-SNE.}

Méthode non linéaire conservant des distances locales via une minimisation de
divergence de Kullback–Leibler entre les voisinages dans l’espace initial et
l’espace réduit.

%%%%%%%%%%%%%%%%%%%%%%%%%%%%%%%%%%%%%%%%%%%%%%%%%%%%%%%%%%%%%%
\subsubsection{Séries temporelles classiques}

\paragraph{Stationnarité.}

Une série $(X_t)$ est stationnaire si :
\[
\mathbb{E}[X_t] = \mu,
\qquad
\mathrm{Var}(X_t) = \sigma^2,
\qquad
\mathrm{Cov}(X_t, X_{t+h}) = \gamma(h).
\]

Les modèles ARMA supposent la stationnarité.

\paragraph{AR(p).}
\[
X_t = \sum_{i=1}^p \phi_i X_{t-i} + \varepsilon_t.
\]

\paragraph{MA(q).}
\[
X_t = \sum_{i=1}^q \theta_i \varepsilon_{t-i} + \varepsilon_t.
\]

\paragraph{ARMA(p,q).}
\[
X_t 
= \sum_{i=1}^p \phi_i X_{t-i}
+ \sum_{j=1}^q \theta_j \varepsilon_{t-j}
+ \varepsilon_t.
\]

\paragraph{ARIMA(p,d,q).}

Pour gérer la non-stationnarité :
\[
\Delta^d X_t \sim \mathrm{ARMA}(p,q).
\]

Les modèles ARIMA sont très utilisés pour la détection de tendances, la prévision
de volatilité, ou l'analyse de séries financières en basse fréquence.


%%%%%%%%%%%%%%%%%%%%%%%%%%%%%%%%%%%%%%%%%%%%%%%%%%%%%%%%%%%%%%
\subsection{Deep Learning}

Le deep learning repose sur les réseaux de neurones artificiels, capables
d’approximer des fonctions complexes par composition de transformations
linéaires et non linéaires. Il permet de traiter des données structurées
(images, séries temporelles, texte) via des architectures spécialisées.

%%%%%%%%%%%%%%%%%%%%%%%%%%%%%%%%%%%%%%%%%%%%%%%%%%%%%%%%%%%%%%
\subsubsection{Bases}

\paragraph{Perceptron.}

Le neurone élémentaire applique :
\[
z = w^\top x + b,
\qquad
a = \sigma(z),
\]
où $\sigma$ est une fonction d’activation (ReLU, sigmoid, tanh, GELU).

\paragraph{MLP (Multi-Layer Perceptron).}

Un réseau profond est une composition :
\[
f(x) = \sigma_L(W_L \sigma_{L-1}(W_{L-1}\cdots \sigma_1(W_1 x))).
\]

\paragraph{Problème du vanishing/exploding gradient.}

Lorsque la profondeur augmente, les gradients peuvent diverger ou s’annuler.
Solutions :
\begin{itemize}
  \item initialisation Xavier/He,
  \item BatchNorm / LayerNorm,
  \item activations adaptées (ReLU, GELU).
\end{itemize}

\paragraph{Normalisation et régularisation.}

\begin{itemize}
  \item \textbf{BatchNorm} : normalisation des activations au sein du batch.
  \item \textbf{LayerNorm} : normalisation par feature.
  \item \textbf{Dropout} : régularisation par mise à zéro aléatoire d’activations.
\end{itemize}

%%%%%%%%%%%%%%%%%%%%%%%%%%%%%%%%%%%%%%%%%%%%%%%%%%%%%%%%%%%%%%
\subsubsection{Convolutional Neural Networks (CNN)}

Les CNN sont adaptés aux données spatiales (images, cartes de features).

\paragraph{Convolution.}

\[
(\text{Conv}(x))_{i,j}
=
\sum_{u,v} K_{u,v}\, x_{i+u,j+v}.
\]

Les noyaux (kernels) extraient des motifs locaux (bords, textures).

\paragraph{Stride, padding.}

Le \emph{stride} contrôle le sous-échantillonnage.  
Le \emph{padding} permet de conserver la taille spatiale.

\paragraph{Pooling.}

\[
\text{max-pool}(x)_{i,j} = \max_{u,v \in \text{patch}} x_{i+u,j+v}.
\]

Enfin, les couches convolutionnelles sont suivies de denses :
\[
\text{Conv} \to \text{ReLU} \to \cdots \to \text{Flatten} \to \text{Dense} \to \text{Softmax}.
\]

%%%%%%%%%%%%%%%%%%%%%%%%%%%%%%%%%%%%%%%%%%%%%%%%%%%%%%%%%%%%%%
\subsubsection{RNN, LSTM, GRU}

Les RNN sont adaptés aux séries temporelles.

\paragraph{RNN simple.}
\[
h_t = \sigma(W_h h_{t-1} + W_x x_t + b).
\]

Limites : pas parallélisable, gradient vanishing.

\paragraph{LSTM.}

Le LSTM introduit des portes :
\[
\begin{aligned}
f_t &= \sigma(W_f x_t + U_f h_{t-1} + b_f),\\
i_t &= \sigma(W_i x_t + U_i h_{t-1} + b_i),\\
\tilde c_t &= \tanh(W_c x_t + U_c h_{t-1} + b_c),\\
c_t &= f_t \odot c_{t-1} + i_t \odot \tilde c_t,\\
h_t &= o_t \odot \tanh(c_t).
\end{aligned}
\]

\paragraph{GRU.}

Structure simplifiée :
\[
\begin{aligned}
z_t &= \sigma(W_z x_t + U_z h_{t-1}),\\
r_t &= \sigma(W_r x_t + U_r h_{t-1}),\\
h_t &= (1-z_t)\odot \tilde h_t + z_t h_{t-1}.
\end{aligned}
\]

%%%%%%%%%%%%%%%%%%%%%%%%%%%%%%%%%%%%%%%%%%%%%%%%%%%%%%%%%%%%%%
\subsubsection{Transformers}

Le Transformer repose sur le \emph{self-attention}.

\paragraph{Tokenisation.}
Les entrées sont découpées en jetons (BPE, WordPiece).

\paragraph{Embeddings.}
\[
x_i \mapsto e_i \in \mathbb{R}^d,
\]
plus un encodage positionnel :
\[
p_t = (\sin(t/10000^{2k/d}),\ \cos(t/10000^{2k/d})).
\]

\paragraph{Attention Q/K/V.}
Chaque token produit trois vecteurs :
\[
Q = W_Q e,\qquad K = W_K e,\qquad V = W_V e.
\]

On calcule :
\[
\mathrm{Attention}(Q,K,V)
=
\mathrm{softmax}\!\left(
\frac{QK^\top}{\sqrt{d}}
\right)V.
\]

Intuition :
- $Q$ “cherche” des informations pertinentes,  
- $K$ dit quelles informations chaque token “porte”,  
- $V$ apporte la valeur combinée.

\paragraph{Entraînement.}

\begin{itemize}
  \item GPT : \textit{next token prediction},
  \item BERT : \textit{masked language modeling}.
\end{itemize}

%%%%%%%%%%%%%%%%%%%%%%%%%%%%%%%%%%%%%%%%%%%%%%%%%%%%%%%%%%%%%%
\subsubsection{Modèles génératifs}

\paragraph{Autoencodeurs.}

Un AE apprend une compression :
\[
z = f_{\theta}(x),
\qquad
\hat x = g_{\theta}(z).
\]

\paragraph{VAE (Variational Autoencoder).}
\[
z \sim \mathcal{N}(\mu(x),\Sigma(x)),
\qquad
\mathcal{L} = \mathbb{E}[\|x-\hat x\|^2] + \mathrm{KL}(q(z|x)\,\|\,p(z)).
\]

\paragraph{VQ-VAE.}
Quantification discrète via dictionnaire de codes.

\paragraph{GAN.}

\[
\min_G \max_D\ 
\mathbb{E}[\log D(x)] + \mathbb{E}[\log(1 - D(G(z)))].
\]

\paragraph{VQ-GAN.}
GAN + quantisation vectorielle (utilisé dans Stable Diffusion).

\paragraph{Diffusion Models.}

\[
x_T \sim \mathcal{N}(0,I),\qquad
x_{t-1} = \mathrm{denoise}_\theta(x_t).
\]
Ce sont les modèles sous-jacents à DALL·E 2, Stable Diffusion, Imagen.



%%%%%%%%%%%%%%%%%%%%%%%%%%%%%%%%%%%%%%%%%%%%%%%%%%%%%%%%%%%%%%
\fi
\subsection{Reinforcement Learning}

Le Reinforcement Learning (RL) étudie l’interaction entre un agent et un
environnement.  
À chaque étape, l’agent observe un état $s_t$, choisit une action $a_t$, reçoit
une récompense $\mathcal{R}_t$, et l’environnement évolue selon une dynamique inconnue.
L’objectif est de maximiser la récompense cumulée.

%%%%%%%%%%%%%%%%%%%%%%%%%%%%%%%%%%%%%%%%%%%%%%%%%%%%%%%%%%%%%%
\subsubsection{Concepts fondamentaux}

Un processus de décision est modélisé par un MDP
(Markov Decision Process) :
\[
(\mathcal{S},\mathcal{A},\mathcal{P},\mathcal{R},\gamma_{\mathrm{RL}}),
\]
où
\begin{itemize}
  \item $\mathcal{S}$ : espace des états,
  \item $\mathcal{A}$ : espace des actions,
  \item $\mathcal{P}(s'|s,a)$ : dynamique (transition),
  \item $\mathcal{R}(s,a)$ : récompense immédiate,
  \item $\gamma_{\mathrm{RL}} \in (0,1)$ : facteur d’escompte (sans lien avec les taux $r_t$).
\end{itemize}

La politique $\pi(a|s)$ définit le comportement de l’agent.

%%%%%%%%%%%%%%%%%%%%%%%%%%%%%%%%%%%%%%%%%%%%%%%%%%%%%%%%%%%%%%
\subsubsection{Fonctions de valeur et équations de Bellman}

La \textbf{valeur d’une politique} $\pi$ est :
\[
V^\pi(s) = 
\mathbb{E}_\pi\!\left[
 \sum_{t=0}^\infty \gamma_{\mathrm{RL}}^t \mathcal{R}(s_t,a_t)
 \,\bigg|\, s_0 = s
\right].
\]

La \textbf{valeur d’état-action} :
\[
Q^\pi(s,a)
=
\mathbb{E}_\pi\!\left[
 \sum_{t=0}^\infty \gamma_{\mathrm{RL}}^t \mathcal{R}(s_t,a_t)
 \,\bigg|\, s_0=s, a_0=a
\right].
\]

Les équations de Bellman caractérisent ces valeurs :
\[
V^\pi(s)
= 
\sum_{a} \pi(a|s)
\sum_{s'} \mathcal{P}(s'|s,a)\bigl(
 \mathcal{R}(s,a) + \gamma_{\mathrm{RL}} V^\pi(s')
 \bigr),
\]
\[
Q^\pi(s,a)
=
 \mathcal{R}(s,a) + \gamma_{\mathrm{RL}} 
\sum_{s'} \mathcal{P}(s'|s,a)
\sum_{a'} \pi(a'|s') Q^\pi(s',a').
\]

La \textbf{valeur optimale} satisfait :
\[
V^\*(s) = \max_{a} Q^\*(s,a).
\]

%%%%%%%%%%%%%%%%%%%%%%%%%%%%%%%%%%%%%%%%%%%%%%%%%%%%%%%%%%%%%%
\subsubsection{Monte Carlo (MC)}

Le MC estime la valeur par \emph{retours échantillonnés} :
\[
G_t = \sum_{k=0}^\infty \gamma_{\mathrm{RL}}^k \mathcal{R}_{t+k}.
\]

Mise à jour :
\[
V(s) \leftarrow V(s) + \alpha (G - V(s)).
\]

Méthodes :
\begin{itemize}
  \item First-Visit MC,
  \item Every-Visit MC,
  \item MC à $\alpha$ constant.
\end{itemize}

%%%%%%%%%%%%%%%%%%%%%%%%%%%%%%%%%%%%%%%%%%%%%%%%%%%%%%%%%%%%%%
\subsubsection{Temporal Difference (TD)}

Les méthodes TD mettent à jour la valeur à partir d'une seule transition :
\[
V(s_t) \leftarrow V(s_t)
+ \alpha \bigl( \mathcal{R}_t + \gamma_{\mathrm{RL}} V(s_{t+1}) - V(s_t) \bigr).
\]

\paragraph{TD(0).}
\[
\delta_t = \mathcal{R}_t + \gamma_{\mathrm{RL}} V(s_{t+1}) - V(s_t).
\]

\paragraph{TD($n$)-step.}
\[
G_t^{(n)} = \sum_{k=0}^{n-1} \gamma_{\mathrm{RL}}^k \mathcal{R}_{t+k} + \gamma_{\mathrm{RL}}^n V(s_{t+n}).
\]

\paragraph{TD($\lambda$).}
Combinaison pondérée des retours $n$-step.

%%%%%%%%%%%%%%%%%%%%%%%%%%%%%%%%%%%%%%%%%%%%%%%%%%%%%%%%%%%%%%
\subsubsection{SARSA, Expected SARSA, Q-Learning}

\paragraph{SARSA (on-policy).}
\[
Q(s_t,a_t)
\leftarrow
Q(s_t,a_t)
+ \alpha
\bigl(
\mathcal{R}_t 
+ \gamma_{\mathrm{RL}} Q(s_{t+1},a_{t+1})
- Q(s_t,a_t)
\bigr).
\]

\paragraph{Expected SARSA.}
\[
Q(s_t,a_t)
\leftarrow
Q(s_t,a_t)
+
\alpha\left(
\mathcal{R}_t + \gamma_{\mathrm{RL}} \sum_{a'} \pi(a'|s_{t+1}) Q(s_{t+1},a')
- Q(s_t,a_t)
\right).
\]

\paragraph{Q-Learning (off-policy).}
\[
Q(s_t,a_t)
\leftarrow
Q(s_t,a_t)
+
\alpha\bigl(
\mathcal{R}_t 
+ \gamma_{\mathrm{RL}} \max_{a'} Q(s_{t+1},a')
- Q(s_t,a_t)
\bigr).
\]

%%%%%%%%%%%%%%%%%%%%%%%%%%%%%%%%%%%%%%%%%%%%%%%%%%%%%%%%%%%%%%
\subsubsection{On-policy vs Off-policy}

\begin{itemize}
  \item \textbf{On-policy} : on apprend la valeur de la politique suivie (SARSA).
  \item \textbf{Off-policy} : on apprend la politique optimale tout en explorant 
        différemment (Q-Learning).
\end{itemize}

Importance sampling permet de corriger les biais off-policy.

%%%%%%%%%%%%%%%%%%%%%%%%%%%%%%%%%%%%%%%%%%%%%%%%%%%%%%%%%%%%%%
\subsubsection{Function Approximation}

Lorsque l’espace des états est trop grand, on approxime :
\[
V(s;\theta),\qquad Q(s,a;\theta).
\]

Mise à jour TD :
\[
\theta \leftarrow \theta 
+ \alpha\, \delta_t \, \nabla_\theta V(s_t;\theta),
\qquad
\delta_t = \mathcal{R}_t + \gamma_{\mathrm{RL}} V(s_{t+1}) - V(s_t).
\]

Bases :
\begin{itemize}
  \item features manuelles (polynômes, splines, bases de Bernstein),
  \item réseaux neuronaux (Deep RL).
\end{itemize}

%%%%%%%%%%%%%%%%%%%%%%%%%%%%%%%%%%%%%%%%%%%%%%%%%%%%%%%%%%%%%%
\subsubsection{Policy Gradient Methods}

Les méthodes de gradient de politique optimisent $\pi_\theta(a|s)$ directement.

\paragraph{Objectif.}
\[
J(\theta) = \mathbb{E}_{\pi_\theta}[G_t].
\]

\paragraph{REINFORCE.}
\[
\theta \leftarrow \theta
+ \alpha\, G_t\, \nabla_\theta \log \pi_\theta(a_t|s_t).
\]

\paragraph{Actor--Critic.}
\[
\theta_{\mathrm{policy}} \leftarrow 
\theta_{\mathrm{policy}}
+ \alpha\, \delta_t\, \nabla_\theta \log \pi_\theta(a_t|s_t),
\]
\[
\theta_{\mathrm{value}} \leftarrow
\theta_{\mathrm{value}}
+ \alpha\, \delta_t\, \nabla_\theta V(s_t;\theta_{\mathrm{value}}).
\]

\paragraph{PPO (Proximal Policy Optimization).}
Méthode stable basée sur une pénalisation du ratio de probabilité :
\[
L^{\mathrm{PPO}}(\theta)
=
\mathbb{E}\Big[
\min\big(
\omega_t(\theta)\,\hat A_t,\,
\mathrm{clip}(\omega_t(\theta),1-\epsilon,1+\epsilon)\hat A_t
\big)
\Big].
\]

%%%%%%%%%%%%%%%%%%%%%%%%%%%%%%%%%%%%%%%%%%%%%%%%%%%%%%%%%%%%%%
\iffalse % Hors p\'erim\`etre : non requis (et trop ``proc\'edural'') pour ce m\'emoire
\subsubsection{Algorithmes évolutionnaires}

Les algorithmes évolutionnaires optimisent une population de politiques :
\[
\theta^{(i)} \leftarrow 
\text{Mutation/Croisement/Selection}(\{\theta^{(j)}\}).
\]

Ils sont utilisés en Deep RL pour explorer de grands espaces de politiques
et éviter les minima locaux.
\fi

%%%%%%%%%%%%%%%%%%%%%%%%%%%%%%%%%%%%%%%%%%%%%%%%%%%%%%%%%%%%%%
\paragraph{Synthèse.}

Le RL combine :
\begin{itemize}
  \item estimation de valeurs (MC, TD),
  \item apprentissage d’actions optimales (SARSA, Q-Learning),
  \item optimisation directe de politiques (Policy Gradient, Actor--Critic, PPO),
  \item approximation de fonctions (Deep RL),
  \item compromis exploration/exploitation et stabilit\'e d'apprentissage.
\end{itemize}

Ces techniques sont essentielles pour les stratégies séquentielles,
l’optimisation de portefeuilles, la gestion dynamique du risque,
et les modèles décisionnels non paramétriques.

\section{M\'ethodologie ``plateforme'' (sans technique) : responsabilit\'es, tra\c{c}abilit\'e et validation}

Ce chapitre d\'ecrit une \emph{m\'ethodologie d'ing\'enierie quantitative} : comment organiser un projet de valorisation,
de calibration et de gestion du risque de fa\c{c}on coh\'erente, audit\-able et r\'ep\'etable, sans entrer dans des choix
techniques. L'enjeu est de rendre explicites les hypoth\`eses et les limites, et de proposer des crit\`eres de validation
adapt\'es \`a des donn\'ees de march\'e imparfaites et \`a des mod\`eles incomplets.

%%%%%%%%%%%%%%%%%%%%%%%%%%%%%%%%%%%%%%%%%%%%%%%%%%%%%%%%%%%%%%
\subsection{D\'efinitions : ``briques'' quantitatives et responsabilit\'es}

\paragraph{Briques conceptuelles.}
On distingue classiquement : (i) \textbf{donn\'ees} (quotes, courbes, conventions), (ii) \textbf{mod\`eles} (dynamique sous-jacente, param\`etres),
(iii) \textbf{m\'ethodes} (PDE, arbres, Monte Carlo, Fourier), (iv) \textbf{d\'ecision} (couverture, allocation, limites), (v) \textbf{reporting}
(PnL, expositions, diagnostics). Une s\'eparation claire de ces responsabilit\'es (analogue \`a une architecture conceptuelle type MVC)
limite les couplages implicites : un changement de convention (taux, dividendes) ou de mod\`ele (Heston $\rightarrow$ Bates) doit avoir un impact
localis\'e et tra\c{c}able.

\paragraph{Hypoth\`ese et gouvernance.}
Une \emph{hypoth\`ese} est un objet explicite : choix de num\'eraire, choix de mesure, conventions de dividendes, param\'etrisation d'une surface,
fonction de perte de calibration, et r\`egles de gestion du risque. La gouvernance consiste \`a documenter ces hypoth\`eses, \`a contr\^oler leur
coh\'erence (absence d'arbitrage statique), et \`a en mesurer la sensibilit\'e.

%%%%%%%%%%%%%%%%%%%%%%%%%%%%%%%%%%%%%%%%%%%%%%%%%%%%%%%%%%%%%%
\subsection{Intuition : tra\c{c}abilit\'e et reproductibilit\'e comme r\'eduction du risque de mod\`ele}

Dans un environnement non stationnaire, la question n'est pas seulement ``quel est le bon mod\`ele ?'', mais aussi
``quelles sont les cons\'equences d'un choix d'hypoth\`ese ?''. La tra\c{c}abilit\'e (donn\'ees utilis\'ees, param\`etres calibr\'es,
pond\'erations, contraintes) permet d'expliquer un PnL, de diagnostiquer une instabilit\'e (donn\'ees vs mod\`ele), et d'\'eviter des erreurs
de type ``m\^emes entr\'ees, sorties diff\'erentes''. Elle est donc une composante directe de la ma\^{i}trise du risque de mod\`ele.

%%%%%%%%%%%%%%%%%%%%%%%%%%%%%%%%%%%%%%%%%%%%%%%%%%%%%%%%%%%%%%
\subsection{R\'esultat cl\'e : analyse de sensibilit\'e et stress tests}

Un sch\'ema minimal formalise l'impact des incertitudes :
\[
\text{Sorties} = g(\text{donn\'ees},\text{hypoth\`eses}),
\]
et, \`a premier ordre, une perturbation $(\delta x,\delta\theta)$ induit
\[
\delta g \approx \nabla_x g \cdot \delta x + \nabla_\theta g \cdot \delta\theta.
\]
Cette d\'ecomposition justifie th\'eoriquement : (i) les \textbf{analyses de sensibilit\'e} (Greeks, sensibilit\'e \`a la courbe, \`a la surface),
et (ii) les \textbf{stress tests} (sc\'enarios sur spot/vol/taux, d\'eformation de smile, widening de bid--ask). Elle rappelle aussi que
la stabilit\'e d'une calibration n'est pas une question d'optimisation seulement, mais de conditionnement du probl\`eme inverse.

%%%%%%%%%%%%%%%%%%%%%%%%%%%%%%%%%%%%%%%%%%%%%%%%%%%%%%%%%%%%%%
\subsection{Crit\`eres de validation : out-of-sample, stabilit\'e, coh\'erence}

\paragraph{Validation out-of-sample.}
Pour un mod\`ele de volatilité, la validation consiste \`a tester la capacit\'e \`a reproduire des quotes hors \'echantillon (autres dates,
autres maturit\'es) et \`a stabiliser des param\`etres. Pour une surface interpol\'ee, on v\'erifie l'absence d'arbitrage (butterfly/calendar)
après interpolation et sous stress.

\paragraph{Validation de couverture.}
Sans r\'esultats empiriques fournis, on formalise la validation par un protocole : simuler/observer un sous-jacent, appliquer une strat\'egie
de couverture \`a fr\'equence donn\'ee, et \'etudier la distribution de l'erreur de couverture (moyenne, variance, queues) sous co\^uts de transaction.

\paragraph{Contr\^oles et coh\'erence.}
Exemples de contr\^oles conceptuels : respect des parit\'es (call--put), bornes et monotonicit\'es, coh\'erence taux/forward,
stabilit\'e des param\`etres \`a l'\'echelle journali\`ere, et coh\'erence entre vol r\'ealis\'ee (historique) et vol implicite (march\'e).

%%%%%%%%%%%%%%%%%%%%%%%%%%%%%%%%%%%%%%%%%%%%%%%%%%%%%%%%%%%%%%
\subsection{Limites}

Une ``plateforme'' ne supprime pas l'incertitude : elle l'organise. La non-stationnarit\'e, la raret\'e d'observations extr\^emes, et les effets de microstructure
limitent la port\'ee des validations. L'objectif est donc la robustesse (sensibilit\'e contr\^ol\'ee), l'explicabilit\'e (attribution, diagnostics) et
la discipline de mod\'elisation (hypoth\`eses document\'ees, stress tests syst\'ematiques).

%%%%%%%%%%%%%%%%%%%%%%%%%%%%%%%%%%%%%%%%%%%%%%%%%%%%%%%%%%%%%%
\subsection{Lien avec le projet \emph{PaperTradingApp}}

Le projet s'inscrit dans cette logique : articuler pricing, calibration, courbe de taux, risque et couverture dans un cadre coh\'erent, en rendant explicites
les hypoth\`eses (conventions, mod\`eles, surfaces) et en proposant des crit\`eres de validation (out-of-sample, stress, sensibilit\'e) compatibles avec un
contexte de paper trading o\`u l'objectif est l'apprentissage et la ma\^{i}trise explicative des briques quantitatives.




%==========================================================
%                 PARTIE 3 : PROJETS PERSONNELS
%==========================================================
\section{Projet : PaperTradingApp}

Cette partie synth\'etise le projet \emph{PaperTradingApp} en explicitant le lien entre les modules fonctionnels d'une
plateforme de paper trading d'options et les notions th\'eoriques pr\'esent\'ees dans les parties pr\'ec\'edentes.
L'enjeu est double : (i) montrer que chaque brique fonctionnelle repose sur des hypoth\`eses de mod\`ele identifiables,
et (ii) rendre claires les limites (donn\'ees, non-stationnarit\'e, co\^uts, risque de mod\`ele).

%----------------------------------------------------------
\subsection{Objectifs, p\'erim\`etre et hypoth\`eses}

Le projet vise \`a fournir un environnement unifi\'e permettant :
\begin{itemize}
  \item d'explorer des \textbf{strat\'egies options} (vanilles et structures simples) ;
  \item de construire une \textbf{lecture coh\'erente des surfaces de volatilit\'e} et de les relier \`a des mod\`eles ;
  \item de comparer des \textbf{m\'ethodes de pricing} (ferm\'ees, semi-ferm\'ees, num\'eriques) et d'en analyser les biais ;
  \item de suivre des \textbf{expositions} (Greeks) et des \textbf{m\'etriques de risque} au niveau portefeuille ;
  \item d'\'etudier des \textbf{r\'egimes de volatilit\'e} et des heuristiques (mean reversion, Markov) utiles \`a la d\'ecision.
\end{itemize}

\subsection{Donn\'ees de march\'e : qualit\'e, horizons et coh\'erence}

La qualit\'e des conclusions d\'epend directement de la qualit\'e des donn\'ees : traitements des valeurs manquantes,
coh\'erence des dates de maturit\'e, gestion des \emph{corporate actions}, granularit\'e (intraday vs daily), et choix de
l'horizon (pricing instantan\'e vs backtest historique). Le m\'emoire insiste sur les \textbf{sources de biais} :
\begin{itemize}
  \item non-synchronisation des quotes (spot vs options),
  \item microstructure et effets de bid--ask,
  \item d\'erives de r\'egime (\emph{non-stationarity}) dans la volatilité et les corr\'elations.
\end{itemize}

\subsection{Options : pricing, volatilit\'e implicite et surfaces}

Le module ``Options'' s'appuie sur :
\begin{itemize}
  \item la d\'efinition rigoureuse des payoffs et des conventions (discounting, dividendes),
  \item l'inversion prix $\leftrightarrow$ volatilit\'e (\emph{implied volatility}),
  \item la construction et la visualisation de surfaces (maturit\'e/strike ou maturit\'e/log-moneyness),
  \item la discussion des \textbf{arbitrages statiques} (monotonie, convexit\'e, calendar).
\end{itemize}

\subsection{Calibration avanc\'ee : du smile au mod\`ele}

La calibration est un probl\`eme inverse : \`a partir d'une surface observ\'ee, on cherche des param\`etres $\theta$
minimisant une erreur (en prix ou en vol implicite). Un sch\'ema g\'en\'erique est :
\[
\theta^\star=\arg\min_{\theta\in\Theta}\sum_{i} w_i\,
\bigl( \sigma^{\text{mkt}}_i-\sigma^{\text{mod}}_i(\theta)\bigr)^2,
\]
avec pond\'erations $w_i$ et contraintes $\Theta$ (bornes, conditions de Feller, etc.).
Les mod\`eles discut\'es incluent les structures affines (Heston), les param\'etrisations asymptotiques (SABR),
les mod\`eles \`a m\'emoire (Volterra/rough) et les diffusions \`a sauts (Merton/Bates). Pour acc\'el\'erer le pricing,
on recourt aux fonctions caract\'eristiques et aux m\'ethodes de Fourier (Carr--Madan/FFT).
Une extension naturelle est l'utilisation de \textbf{surrogates} (r\'eseaux) pour approximer la carte
surface $\mapsto$ param\`etres, afin d'obtenir des initialisations robustes.

\subsection{Courbe des taux : discounting, forwards, param\'etrisation}

La courbe z\'ero-coupon fournit le lien entre taux, discount factors et taux forwards. Le module ``Yield Curve'' met en
sc\`ene (i) la conversion yield $\leftrightarrow$ discount factor, (ii) l'interpolation/extrapolation, (iii) le
bootstrapping, et (iv) des formes param\'etriques (Nelson--Siegel) permettant une repr\'esentation lisse et interpr\'etable
(niveau, pente, courbure).

\subsection{Couverture : Greeks, erreur de hedging, d\'ecision s\'equentielle}

La couverture (delta/gamma/vega) se comprend comme un compromis entre erreur de r\'eplication et co\^uts de transaction.
Le projet met en regard :
\begin{itemize}
  \item une couverture ``classique'' issue de la th\'eorie Black--Scholes,
  \item des extensions sous volatilit\'e stochastique (greeks locaux vs greeks de mod\`ele),
  \item une formulation \`a base de \textbf{MDP} o\`u l'action est un ajustement de couverture et la r\'ecompense p\'enalise
        une exposition r\'esiduelle et/ou des co\^uts.
\end{itemize}

\subsection{Portefeuille \& risque : m\'etriques, attribution, allocation}

Au niveau portefeuille, le projet articule :
\begin{itemize}
  \item mesures de risque (expositions brutes/nettes, drawdowns, VaR/ES),
  \item d\'ecomposition du PnL (r\'ealis\'e, latent, par symbole),
  \item allocations de r\'ef\'erence (Markowitz, risk parity, eigen-portfolios) et questions d'estimation de covariance.
\end{itemize}

\subsection{Strat\'egies et r\'egimes de volatilité : straddle, IV crush, Markov}

Le module ``Bots'' sert de laboratoire d'id\'ees : straddle (sensibilit\'e \`a la vol), \emph{IV crush} (choc de vol
implicite autour d'\'ev\'enements), r\'egimes de vol r\'ealis\'ee (percentiles), diagnostics de mean reversion via
r\'egressions, et cha\^ines de Markov sur r\'egimes discrets (matrice de transition, probabilit\'es conditionnelles).

\subsection{Assistance et limites}

Un assistant textuel peut aider \`a synth\'etiser un \'etat de portefeuille, mais ne remplace ni un mod\`ele ni une
proc\'edure de contr\^ole : biais d'interpr\'etation, incertitudes non quantifi\'ees, et risques de sur-confiance. Le
m\'emoire insiste sur la n\'ecessit\'e d'encadrer ces outils par des r\`egles de prudence (validation, tra\c{c}abilit\'e,
absence d'action non contr\^ol\'ee).

%==========================================================
%                 PARTIE 4 : CONCLUSION
%==========================================================
\section{Conclusion et ouverture}

Ce document a parcouru les fondements de la finance quantitative moderne en prenant pour fil conducteur le projet
\emph{PaperTradingApp} : produits d\'eriv\'es, mod\'elisation stochastique, pricing et Greeks, volatilit\'e implicite et
calibration, construction de courbe de taux, gestion du risque et allocation, ainsi qu'une ouverture vers des outils
d'apprentissage utiles \`a la calibration (surrogates) et \`a la d\'ecision s\'equentielle (MDP/RL).

L’ensemble constitue une base conceptuelle solide pour analyser un syst\`eme de trading d\'eriv\'es de bout en bout :
donn\'ees, hypoth\`eses de mod\`eles, incertitudes, limites num\'eriques et contr\^oles de risque.

%%%%%%%%%%%%%%%%%%%%%%%%%%%%%%%%%%%%%%%%%%%%%%%%%%%%%%%%%%%%%%
\subsection{Synthèse}

Les points suivants peuvent \^etre retenus :

\begin{itemize}
  \item Les produits d\'eriv\'es (vanilles et exotiques) s’analysent via la mesure risque--neutre, les EDS et des
        m\'ethodes num\'eriques (arbres, PDE, Monte Carlo, Fourier/FFT).
  \item Les Greeks structurent la lecture du risque et la couverture (delta/gamma/vega), et permettent d’anticiper les
        erreurs de hedging sous co\^uts de transaction et risque de mod\`ele.
  \item Les surfaces de volatilit\'e implicite r\'esument l'information options ; la calibration relie cette surface \`a des
        mod\`eles (Heston, SABR, rough/Volterra, sauts) et met en \'evidence les compromis de param\'etrisation.
  \item La courbe de taux (discounting/forwards) est un ingr\'edient transversal : coh\'erence absence d’arbitrage,
        choix d'interpolation/extrapolation, param\'etrisations lisses (Nelson--Siegel).
  \item La gestion du risque au niveau portefeuille n\'ecessite des indicateurs actionnables (PnL, drawdowns, VaR/ES),
        ainsi que des r\`egles d'allocation (\`a la Markowitz, risk parity, eigen-portfolios).
  \item L'apprentissage apporte des outils compl\'ementaires : surrogates pour acc\'el\'erer la calibration et formulation
        MDP/RL pour des probl\`emes de d\'ecision s\'equentielle (couverture, rebalancement).
\end{itemize}

%%%%%%%%%%%%%%%%%%%%%%%%%%%%%%%%%%%%%%%%%%%%%%%%%%%%%%%%%%%%%%
\subsection{Ouverture}

Plusieurs directions naturelles s’ouvrent à la suite de ce travail :

\begin{itemize}
  \item \textbf{Qualit\'e des donn\'ees et microstructure} :
        mod\'eliser explicitement bid--ask, co\^uts de transaction et non-synchronisme spot/options, et mesurer l'impact sur
        surfaces, calibration, backtests et risque.
  \item \textbf{Volatilité rugueuse et SDE de Volterra} : 
        extension non markovienne des mod\`eles de variance, offrant une 
        flexibilit\'e accrue pour ajuster des smiles courts (selon les donn\'ees et la pond\'eration),
        et une mod\'elisation explicite de la m\'emoire longue.
  \item \textbf{Modèles multi-facteurs et modélisation affine avancée} :
        dynamiques jointes des variances, taux et spreads, incluant corr\'elations
        stochastiques.
  \item \textbf{Acc\'el\'eration num\'erique et stabilit\'e} :
        r\'eduction de variance, sch\'emas de discr\'etisation stables, FFT (Carr--Madan) et approximations markoviennes
        des mod\`eles \`a m\'emoire.
  \item \textbf{Surrogates pour le pricing et la calibration} :
        approximations supervis\'ees de cartes $(K,T,\theta)\mapsto C$ ou $\sigma_{\mathrm{imp}}$, utilisables pour initialiser
        des calibrations et acc\'el\'erer des diagnostics, sous contraintes arbitrage-free.
  \item \textbf{Reinforcement Learning pour la gestion dynamique du risque} :
        agents apprenant des stratégies de hedging optimales 
        (delta-gamma-vega), optimisation de portefeuilles, ou exécution optimale.
  \item \textbf{Validation out-of-sample et contr\^ole de risque de mod\`ele} :
        protocoles de validation, stress tests, r\'egularisation, et mesure de la stabilit\'e des param\`etres calibr\'es.
\end{itemize}

%%%%%%%%%%%%%%%%%%%%%%%%%%%%%%%%%%%%%%%%%%%%%%%%%%%%%%%%%%%%%%
\subsection{Conclusion générale}

Cette synthèse met en évidence l’interconnexion entre :

\begin{itemize}
  \item maths appliquées (EDS, variation quadratique, PDE),
  \item finance quantitative (pricing, surfaces de volatilité, calibrations),
  \item machine learning (apprentissage supervisé, non supervisé, deep learning),
  \item reinforcement learning (prise de décision séquentielle).
\end{itemize}

L’objectif est de fournir une vision cohérente et moderne des outils quantitatifs
utilisables en finance de marché.  
Les annexes rassemblent des compl\'ements (Heston, SABR, Rough Heston, m\'ethodes Fourier/PDE) pour servir de base \`a des travaux plus avanc\'es.


%==========================================================
%                 PARTIE 5 : BIBLIOGRAPHIE
%==========================================================
\section{Bibliographie comment\'ee (ouvrages et articles)}

Cette bibliographie est \emph{align\'ee avec le contenu du m\'emoire}. Chaque r\'ef\'erence est accompagn\'ee d'un bref commentaire indiquant
ce qu'elle apporte (cadre th\'eorique, mod\`ele, m\'ethode) et \`a quels chapitres elle se rattache. Elle n'a pas vocation \`a \^etre exhaustive.

%%%%%%%%%%%%%%%%%%%%%%%%%%%%%%%%%%%%%%%%%%%%%%%%%%%%%%%%%%%%%%
\subsection{Absence d'arbitrage, mesures martingales, calcul stochastique}

\begin{description}[style=nextline,leftmargin=2.8cm]
  \item[\textbf{Björk, T. (2009).}] \emph{Arbitrage Theory in Continuous Time}.\\
  R\'ef\'erence structurante sur l'absence d'arbitrage, les changements de num\'eraire et les mesures martingales (cadre du pricing risque--neutre).

  \item[\textbf{Delbaen, F., Schachermayer, W. (2006).}] \emph{The Mathematics of Arbitrage}.\\
  Traitement approfondi du th\'eor\`eme fondamental de l'asset pricing (NFLVR) et de ses hypoth\`eses techniques (chapitre A).

  \item[\textbf{Shreve, S. (2004).}] \emph{Stochastic Calculus for Finance II}.\\
  D\'eveloppe It\^o, Girsanov, Feynman--Kac et les PDE de pricing ; utile pour relier dynamique sous-jacente et valorisation (chapitre C).
\end{description}

%%%%%%%%%%%%%%%%%%%%%%%%%%%%%%%%%%%%%%%%%%%%%%%%%%%%%%%%%%%%%%
\subsection{Options, pricing et m\'ethodes num\'eriques (arbres, MC, Fourier)}

\begin{description}[style=nextline,leftmargin=2.8cm]
  \item[\textbf{Hull, J. (2022).}] \emph{Options, Futures, and Other Derivatives}.\\
  Panorama des produits, parit\'es, bornes, et introduction aux pricers (BS, arbres, MC) (chapitres B--D, H).

  \item[\textbf{Cox, J., Ross, S., Rubinstein, M. (1979).}] \emph{Option Pricing: A Simplified Approach}.\\
  Article fondateur des arbres binomiaux CRR et de leur logique risque--neutre (chapitre D).

  \item[\textbf{Glasserman, P. (2004).}] \emph{Monte Carlo Methods in Financial Engineering}.\\
  Ouvrage de r\'ef\'erence pour estimateurs, r\'eduction de variance et Greeks Monte Carlo (chapitre D).

  \item[\textbf{Longstaff, F., Schwartz, E. (2001).}] \emph{Valuing American Options by Simulation}.\\
  M\'ethode LSMC : r\'egression pour l'exercice anticip\'e (am\'ericaines) ; utile pour comprendre l'approximation et ses biais (chapitre D).

  \item[\textbf{Carr, P., Madan, D. (1999).}] \emph{Option Valuation Using the Fast Fourier Transform}.\\
  Cadre Carr--Madan : damping, transform\'ee et FFT pour prix europ\'eens \`a partir de $\phi_T$ (chapitres D, F ; annexe F).

  \item[\textbf{Lewis, A. (2001).}] \emph{A Simple Option Formula for General Jump-Diffusion and Other Exponential L\'evy Processes}.\\
  Repr\'esentations Fourier g\'en\'erales en pr\'esence de sauts ; utile pour Merton/Bates (chapitre F).
\end{description}

%%%%%%%%%%%%%%%%%%%%%%%%%%%%%%%%%%%%%%%%%%%%%%%%%%%%%%%%%%%%%%
\subsection{Volatilit\'e implicite, surfaces et arbitrage statique}

\begin{description}[style=nextline,leftmargin=2.8cm]
  \item[\textbf{Gatheral, J. (2006).}] \emph{The Volatility Surface: A Practitioner's Guide}.\\
  R\'ef\'erence sur smile/skew, param\'etrisations (moneyness), et contraintes d'arbitrage statique (chapitre E).

  \item[\textbf{Breeden, D., Litzenberger, R. (1978).}] \emph{Prices of State-Contingent Claims Implicit in Option Prices}.\\
  R\'esultat central liant $\partial_{KK}C$ \`a la densit\'e risque--neutre (chapitre E).

  \item[\textbf{Gatheral, J., Jacquier, A. (2014).}] \emph{Arbitrage-free SVI volatility surfaces}.\\
  Conditions d'absence d'arbitrage (butterfly/calendar) pour SVI ; utile si l'on veut une surface param\'etrique robuste (chapitre E, annexes propos\'ees).
\end{description}

%%%%%%%%%%%%%%%%%%%%%%%%%%%%%%%%%%%%%%%%%%%%%%%%%%%%%%%%%%%%%%
\subsection{Mod\`eles de volatilit\'e : Heston, SABR, rough, sauts}

\begin{description}[style=nextline,leftmargin=2.8cm]
  \item[\textbf{Heston, S. (1993).}] \emph{A Closed-Form Solution for Options with Stochastic Volatility}.\\
  Mod\`ele affine de r\'ef\'erence ; fonction caract\'eristique et pricing semi-analytique (chapitre F ; annexe A).

  \item[\textbf{Hagan, P., Kumar, D., Lesniewski, A., Woodward, D. (2002).}] \emph{Managing Smile Risk}.\\
  Formule asymptotique SABR (Hagan) et interpr\'etation des param\`etres ; standard de calibration (chapitre F ; annexe D).

  \item[\textbf{Merton, R. (1976).}] \emph{Option pricing when underlying stock returns are discontinuous}.\\
  Jump--diffusion de Merton : motivation (queues) et pricing (chapitre F).

  \item[\textbf{Bates, D. (1996).}] \emph{Jumps and stochastic volatility: exchange rate processes implicit in option prices}.\\
  Combinaison volatilité stochastique + sauts (Bates) ; utile pour smiles courts et wings (chapitre F).

  \item[\textbf{Gatheral, J., Jaisson, T., Rosenbaum, M. (2018).}] \emph{Volatility is rough}.\\
  Motivation rough volatility (Hurst $<1/2$) et cons\'equences sur le skew court terme (chapitre F).

  \item[\textbf{Bayer, C., Friz, P., Gatheral, J. (2016).}] \emph{Pricing under rough volatility}.\\
  Questions num\'eriques et mod\'elisation associ\'ees \`a la rough volatility (chapitre F, annexe E).

  \item[\textbf{El Euch, O., Rosenbaum, M. (2019).}] \emph{The characteristic function of rough Heston models}.\\
  Rough Heston : cadre Volterra, fonction caract\'eristique et pricing (chapitre F ; annexe E).
\end{description}

%%%%%%%%%%%%%%%%%%%%%%%%%%%%%%%%%%%%%%%%%%%%%%%%%%%%%%%%%%%%%%
\subsection{Courbe de taux et param\'etrisations}

\begin{description}[style=nextline,leftmargin=2.8cm]
  \item[\textbf{Brigo, D., Mercurio, F. (2006).}] \emph{Interest Rate Models: Theory and Practice}.\\
  Cadre complet pour taux, discounting, bootstrapping et mod\`eles de taux (chapitre H).

  \item[\textbf{Nelson, C., Siegel, A. (1987).}] \emph{Parsimonious Modeling of Yield Curves}.\\
  Mod\`ele Nelson--Siegel : d\'ecomposition niveau/pente/courbure (chapitre H).

  \item[\textbf{Svensson, L. (1994).}] \emph{Estimating and Interpreting Forward Interest Rates}.\\
  Extension Svensson (plus flexible) et lecture \'economique des composantes (chapitre H).
\end{description}

%%%%%%%%%%%%%%%%%%%%%%%%%%%%%%%%%%%%%%%%%%%%%%%%%%%%%%%%%%%%%%
\subsection{Gestion du risque, VaR/ES, allocation}

\begin{description}[style=nextline,leftmargin=2.8cm]
  \item[\textbf{Artzner, P., Delbaen, F., Eber, J.-M., Heath, D. (1999).}] \emph{Coherent Measures of Risk}.\\
  Cadre axiomatique de coh\'erence (subadditivit\'e, etc.) ; utile pour discuter VaR vs ES (chapitre I).

  \item[\textbf{McNeil, A., Frey, R., Embrechts, P. (2015).}] \emph{Quantitative Risk Management}.\\
  R\'ef\'erence sur mod\'elisation des queues, VaR/ES, backtesting et d\'ependance (chapitre I).

  \item[\textbf{Markowitz, H. (1952).}] \emph{Portfolio Selection}.\\
  Texte fondateur de la moyenne--variance et de la fronti\`ere efficiente (chapitre I).

  \item[\textbf{Ledoit, O., Wolf, M. (2004).}] \emph{A well-conditioned estimator for large-dimensional covariance matrices}.\\
  Shrinkage de covariance : r\'egularisation essentielle pour stabiliser Markowitz et risk parity (chapitre I).
\end{description}

%%%%%%%%%%%%%%%%%%%%%%%%%%%%%%%%%%%%%%%%%%%%%%%%%%%%%%%%%%%%%%
\subsection{Filtrage, r\'egimes et d\'ecision s\'equentielle (MDP/RL)}

\begin{description}[style=nextline,leftmargin=2.8cm]
  \item[\textbf{Kalman, R. (1960).}] \emph{A New Approach to Linear Filtering and Prediction Problems}.\\
  Filtre de Kalman : estimation r\'ecursive d'\'etats latents en cadre lin\'eaire gaussien (chapitre F).

  \item[\textbf{Durbin, J., Koopman, S. (2012).}] \emph{Time Series Analysis by State Space Methods}.\\
  Approche moderne \'etat-espace, lissage (RTS), estimation par vraisemblance (chapitres F, K).

  \item[\textbf{Hamilton, J. (1989).}] \emph{A new approach to the economic analysis of nonstationary time series and the business cycle}.\\
  Mod\`eles de r\'egimes (Markov switching) ; utile pour la lecture probabiliste des phases de march\'e (chapitre K).

  \item[\textbf{Sutton, R., Barto, A. (2018).}] \emph{Reinforcement Learning: An Introduction}.\\
  R\'ef\'erence g\'en\'erale sur MDP, Bellman, MC/TD et Q-learning (chapitre J).

  \item[\textbf{Bertsekas, D. (2012).}] \emph{Dynamic Programming and Optimal Control}.\\
  Point de vue contr\^ole/DP compl\'ementaire, utile pour relier couverture et optimisation (chapitre J).

  \item[\textbf{Mnih, V. et al. (2015).}] \emph{Human-level control through deep reinforcement learning}.\\
  DQN : approximation de $Q$ par r\'eseau et stabilisation (experience replay, targets) (chapitre J).

  \item[\textbf{Schulman, J. et al. (2017).}] \emph{Proximal Policy Optimization Algorithms}.\\
  PPO : stabilisation des policy gradients via clipping (chapitre J).
\end{description}

%%%%%%%%%%%%%%%%%%%%%%%%%%%%%%%%%%%%%%%%%%%%%%%%%%%%%%%%%%%%%%
\subsection{Pistes d'annexes (au-del\`a de celles incluses)}

\begin{itemize}
  \item \textbf{Arbitrage de surface} : d\'erivation d\'etaill\'ee des contraintes butterfly/calendar et conditions SVI arbitrage-free.
  \item \textbf{Bates} : fonction caract\'eristique et formule Carr--Madan (mise en parall\`ele avec Heston).
  \item \textbf{Kalman} : r\'ecursions filtre/lisseur (RTS) et log-vraisemblance pour calibration de param\`etres latents.
  \item \textbf{R\'egimes} : estimation EM d'un Markov switching gaussien, filtres (probabilit\'es a posteriori) et limites.
  \item \textbf{Risque} : backtesting VaR/ES (tests de couverture) et sensibilit\'e \`a la non-stationnarit\'e.
  \item \textbf{Couverture} : d\'erivation d'une approximation de l'erreur de delta-hedging en fonction de $\Gamma$ et de la variance r\'ealis\'ee.
\end{itemize}

\iffalse % Ancienne bibliographie (ressources g\'en\'eralistes) conserv\'ee hors rendu final

%%%%%%%%%%%%%%%%%%%%%%%%%%%%%%%%%%%%%%%%%%%%%%%%%%%%%%%%%%%%%%
\subsection{Finance quantitative}

\begin{itemize}
  \item \textbf{Hull, J.}  
        \emph{Options, Futures and Other Derivatives}.  
        Référence classique couvrant les dérivés, les modèles de taux, les
        produits structurés et la gestion du risque.

  \item \textbf{Natenberg, S.}  
        \emph{Option Volatility and Pricing}.  
        Ouvrage fondamental sur le smile, les Greeks, les stratégies de trading
        et l’arbitrage.

  \item \textbf{Gatheral, J.}  
        \emph{The Volatility Surface: A Practitioner’s Guide}.  
        Référence sur la surface de volatilité, le smile/skew, les conditions d’arbitrage statique et les paramétrisations.

  \item \textbf{Heston, S.} (1993).  
        \emph{A Closed-Form Solution for Options with Stochastic Volatility}.  
        Article fondateur sur la volatilité stochastique affine et le pricing via fonction caractéristique.

  \item \textbf{Carr, P., Madan, D.} (1999).  
        \emph{Option Valuation Using the Fast Fourier Transform}.  
        Méthode FFT (type Carr--Madan) pour valoriser rapidement des options européennes sur une grille de strikes.

  \item \textbf{Hagan, P., Kumar, D., Lesniewski, A., Woodward, D.} (2002).  
        \emph{Managing Smile Risk}.  
        Formule asymptotique SABR et interprétation des paramètres (alpha, beta, rho, nu).

  \item \textbf{Glasserman, P.}  
        \emph{Monte Carlo Methods in Financial Engineering}.  
        Ouvrage de référence sur Monte Carlo, réduction de variance et applications au pricing/greeks.

  \item \textbf{Markowitz, H.} (1952).  
        \emph{Portfolio Selection}.  
        Texte fondateur de l’allocation moyenne--variance et de la frontière efficiente.

  \item \textbf{Artzner, P., Delbaen, F., Eber, J.-M., Heath, D.} (1999).  
        \emph{Coherent Measures of Risk}.  
        Cadre axiomatique des mesures de risque (subadditivité, homogénéité, monotonicité, invariance par translation).

  \item \textbf{Bouchaud, J.-P., Potters, M.}  
        \emph{Theory of Financial Risk and Derivative Pricing}.  
        Approche statistique et empirique de la volatilité, du risque et des
        phénomènes de queue.

  \item \textbf{Darbyshire, M.}  
        \emph{Pricing and Trading Interest Rate Derivatives}.  
        Référence sur les dérivés de taux, IRS, caps/floors, swaptions.

  \item \textbf{Wooldridge, J.}  
        \emph{Econometrics}.  
        Fondements en économétrie, indispensable pour les modèles statistiques
        appliqués à la finance.
\end{itemize}

%%%%%%%%%%%%%%%%%%%%%%%%%%%%%%%%%%%%%%%%%%%%%%%%%%%%%%%%%%%%%%
\subsection{FinTech et Data Science appliquée à la finance}

\begin{itemize}
  \item \textbf{Kaabar, M.}  
        \emph{Deep Learning for Finance}.  

  \item \textbf{Hilpisch, Y.}  
        \emph{Reinforcement Learning for Finance}.  

  \item \textbf{Kanungo, D.}  
        \emph{Probabilistic Machine Learning for Finance and Investing}.  

  \item \textbf{Tatsat, R.}  
        \emph{Machine Learning \& Data Science Blueprints for Finance}.  

  \item \textbf{Hanson, G.}  
        \emph{Modern C++ for Finance}.  

  \item \textbf{Lopez de Prado, M.}  
        \emph{Advances in Financial Machine Learning}.  

  \item \textbf{Chan, E.}  
        \emph{Algorithmic Trading: Winning Strategies and Their Rationale}.  
\end{itemize}

%%%%%%%%%%%%%%%%%%%%%%%%%%%%%%%%%%%%%%%%%%%%%%%%%%%%%%%%%%%%%%
\subsection{Mathématiques et probabilités}

\begin{itemize}
  \item \textbf{Shreve, S.}  
        \emph{Stochastic Calculus for Finance I \& II}.  
        Références incontournables pour le calcul stochastique appliqué aux dérivés.

  \item \textbf{Ségonne, R.}  
        \emph{Quant Basics}.  
        Introduction rigoureuse aux EDS, au pricing et aux outils quantitatifs.

  \item \textbf{Cohen, N.}  
        \emph{Practical Linear Algebra for Data Science}.  

  \item \textbf{Nelson, R.}  
        \emph{Essential Math for AI}.  
\end{itemize}

%%%%%%%%%%%%%%%%%%%%%%%%%%%%%%%%%%%%%%%%%%%%%%%%%%%%%%%%%%%%%%
\subsection{Machine Learning}

\begin{itemize}
  \item \textbf{Goodfellow, Bengio, Courville}.  
        \emph{Deep Learning}.  
        Ouvrage fondateur pour le deep learning moderne.

  \item \textbf{Sutton, Barto}.  
        \emph{Reinforcement Learning: An Introduction}.  
        Livre de référence sur le RL, dynamique des MDP et méthodes TD.

  \item \textbf{Gillard, D.}  
        \emph{Computer Vision}.  

  \item \textbf{Foster, D.}  
        \emph{Generative Deep Learning}.  

  \item \textbf{Tunstall, von Werra, Wolf}.  
        \emph{NLP with Transformers}.  

  \item \textbf{Bishop, C.}  
        \emph{Pattern Recognition and Machine Learning}.  

  \item \textbf{Hastie, T., Tibshirani, R., Friedman, J.}  
        \emph{The Elements of Statistical Learning}.  

  \item \textbf{Kirill Eremenko}.  
        \emph{Hands-On Machine Learning with C++}.  

  \item \textbf{Moses, Dezy}.  
        \emph{Data Quality Fundamentals}.  

  \item \textbf{Auffarth, B.}  
        \emph{Machine Learning for Time Series}.  
\end{itemize}

%%%%%%%%%%%%%%%%%%%%%%%%%%%%%%%%%%%%%%%%%%%%%%%%%%%%%%%%%%%%%%
\subsection{Vidéos de Reinforcement Learning}

\vspace{0.15cm}
Chaîne YouTube \emph{Reinforcement Learning — by the Book} :
\begin{itemize}
  \item Bellman Equations, Dynamic Programming — Part 2
  \item Monte Carlo \& Off-Policy Methods — Part 3
  \item Temporal Difference Learning — Part 4
  \item Function Approximation — Part 5
  \item Policy Gradient Methods — Part 6
\end{itemize}

%%%%%%%%%%%%%%%%%%%%%%%%%%%%%%%%%%%%%%%%%%%%%%%%%%%%%%%%%%%%%%
\subsection{Ressources vidéo : mathématiques}

\begin{itemize}
  \item \emph{Practical Linear Algebra} — Introduction and Motivation  
  \item Fourier Analysis — Overview  
  \item Differential Equations and Dynamical Systems — Overview  
  \item Probability and Statistics — Overview  
  \item Introduction to Statistics and Data Analysis  
  \item Laplace Transform — Generalized Fourier Transform  
  \item Inner Products in Hilbert Space  
  \item Properties of the Expected Value  
  \item Bayesian MAP Estimation  
  \item Lecture 1 (Motivation and Introduction)  
\end{itemize}

%%%%%%%%%%%%%%%%%%%%%%%%%%%%%%%%%%%%%%%%%%%%%%%%%%%%%%%%%%%%%%
\subsection{Ressources vidéo : statistiques}

\begin{itemize}
  \item Maximum Likelihood (Normal Distribution)  
  \item Why $\pi$ appears in the Gaussian distribution  
  \item Multivariate Gaussian — definition and intuition  
  \item The Kelly Criterion  
  \item Learning Probability Distributions  
  \item Chi-Square Distribution  
  \item Jensen’s Inequality  
  \item Accept-Reject Sampling  
  \item Importance Sampling  
  \item Softmax — intuition  
  \item Linear Regression — 12 minutes  
  \item Bernstein Basis  
  \item Maximum Likelihood — explained  
  \item Exponential Family (Part 1 \& 2)  
  \item Gaussian Processes — overview  
  \item Bias–Variance Tradeoff  
  \item Fisher Information  
  \item Factor Analysis \& Probabilistic PCA  
  \item Maximum Entropy Principle  
\end{itemize}


\fi
%==========================================================
%                 ANNEXES
%==========================================================
\appendix

\subsection*{Annexe A - Mod\`ele de Heston : d\'erivation compl\`ete}
\addcontentsline{toc}{section}{Annexe A - Mod\`ele de Heston}



\subsection{Mod\`ele de Heston et log--prix}

Sous la mesure risque--neutre $Q$, le prix de l'actif $S_t$ et sa variance
instantan\'ee $v_t$ suivent
\begin{equation}
\begin{cases}
dS_t = r S_t\,dt + \sqrt{v_t}\,S_t\,dW^{(1)}_t, \\[3pt]
dv_t = \kappa(\theta - v_t)\,dt + \xi\sqrt{v_t}\,dW^{(2)}_t, \\[3pt]
dW^{(1)}_t\,dW^{(2)}_t = \rho\,dt,
\end{cases}
\end{equation}
avec $v_0 > 0$, $\kappa>0$, $\theta>0$, $\xi>0$, $|\rho|\le 1$.

On pose le log--prix
\begin{equation}
X_t = \ln S_t.
\end{equation}
Par la formule d'It\^o,
\begin{equation}
dX_t = \left(r - \frac{v_t}{2}\right)dt + \sqrt{v_t}\,dW^{(1)}_t.
\end{equation}

On consid\`ere une option call europ\'eenne de strike $K>0$ et maturit\'e $T>0$,
de payoff $(S_T-K)^+$. Son prix au temps $0$ est
\begin{equation}
C(K,T) = e^{-rT}\,\mathbb{E}_Q\big[(S_T - K)^+\big].
\end{equation}

%%%%%%%%%%%%%%%%%%%%%%%%%%%%%%%%%%%%%%%%%%%%%%%%%%%%%%%%%%%%
\subsection{Fonction caract\'eristique du log--prix sous Heston}

On d\'efinit la fonction caract\'eristique du log--prix $X_T = \ln S_T$ par
\begin{equation}
\varphi_T(u) := \mathbb{E}\!\left[ e^{i u X_T} \right], \qquad u\in\mathbb{C}.
\end{equation}
On veut une expression explicite de $\varphi_T(u)$.

%%%%%%%%%%%%%%%%%%%%%%%%%%%%%%%%%%%%
\subsection{PDE de Feynman--Kac et ansatz affine}

Pour $0\le t\le T$, $x\in\mathbb{R}$, $v\ge0$, on introduit
\begin{equation}
p(t,x,v;u) := \mathbb{E}\big[ e^{i u X_T} \,\big|\, X_t=x, v_t=v \big].
\end{equation}
Alors $p(T,x,v;u) = e^{iux}$ et
\begin{equation}
\varphi_T(u) = p(0,\ln S_0, v_0;u).
\end{equation}

Le couple $(X_t,v_t)$ a pour g\'en\'erateur infinit\'esimal
\begin{equation}
\mathcal{L} f =
\left(r - \frac{v}{2}\right)\partial_x f
+ \kappa(\theta - v)\partial_v f
+ \frac{1}{2} v\,\partial_{xx} f
+ \rho\xi v\,\partial_{xv} f
+ \frac{1}{2}\xi^2 v\,\partial_{vv} f.
\end{equation}
Par Feynman--Kac, $p$ satisfait
\begin{equation}
\partial_t p + \mathcal{L}p = 0, \qquad p(T,x,v;u) = e^{iux}.
\end{equation}

On cherche une solution de la forme affine en $v$:
\begin{equation}
p(t,x,v;u) = \exp\!\big(C(\tau;u) + D(\tau;u)v + iux\big),
\qquad \tau := T-t.
\end{equation}
On impose $C(0;u)=0$, $D(0;u)=0$, de sorte que $p(T,x,v;u)=e^{iux}$.

On calcule
\[
\partial_t p = -C'(\tau;u)p - D'(\tau;u)v p,\quad
\partial_x p = iup,\quad
\partial_{xx}p = -u^2p,
\]
\[
\partial_v p = Dp,\quad
\partial_{vv}p = D^2p,\quad
\partial_{xv}p = i u D p.
\]

En injectant dans la PDE, puis en factorisant par $p\neq0$, on obtient
\begin{align}
0 &= -C' - D'v 
+ \left(r - \frac{v}{2}\right)(iu)
+ \kappa(\theta - v)D
+ \frac{1}{2} v(-u^2)
+ \rho\xi v(iuD)
+ \frac{1}{2}\xi^2 v D^2.
\end{align}
On s\'eparant les termes constants et les termes proportionnels \`a $v$, ce qui donne
\begin{align}
\text{(terme constant)}  &:& -C' + iur + \kappa\theta D &= 0,\\
\text{(terme en $v$)}    &:& -D'
 + \left(-\frac{iu}{2}\right)
 - \kappa D
 - \frac{u^2}{2}
 + \rho\xi iu D
 + \frac{1}{2}\xi^2 D^2 &= 0.
\end{align}
On en d\'eduit les EDO (en $\tau$):
\begin{align}
C'(\tau;u) &= \kappa\theta D(\tau;u) + i u r, \qquad C(0;u)=0,
\label{eq:Cprime}\\[3pt]
D'(\tau;u) &= \frac{1}{2}\xi^2 D(\tau;u)^2
           + (\rho\xi i u - \kappa)D(\tau;u)
           + \frac{1}{2}(u^2 + i u),\qquad D(0;u)=0.
\label{eq:Dprime}
\end{align}

L'\'equation~\eqref{eq:Dprime} est de Riccati.

%%%%%%%%%%%%%%%%%%%%%%%%%%%%%%%%%%%%
\subsection{R\'esolution de l'\'equation de Riccati}

On note
\begin{equation}
a := \frac{\xi^2}{2},\qquad
\beta(u) := \rho\xi i u - \kappa,\qquad
\omega(u) := \frac{1}{2}(u^2+i u).
\end{equation}
Alors
\begin{equation}
D' = a D^2 + \beta(u) D + \omega(u).
\end{equation}

On introduit le discriminant
\begin{equation}
d(u) := \sqrt{\beta(u)^2 + 4a\omega(u)}
      = \sqrt{(\rho\xi i u - \kappa)^2 + \xi^2(u^2+i u)},
\end{equation}
et on prend la racine avec $\Re[d(u)]>0$.

On d\'efinit ensuite
\begin{equation}
g(u) := \frac{\kappa - \rho\xi i u - d(u)}{\kappa - \rho\xi i u + d(u)}.
\end{equation}

La solution classique de l'\'equation de Riccati est alors
\begin{equation}
D(\tau;u)
= \frac{\kappa - \rho\xi i u - d(u)}{\xi^2}\;
  \frac{1 - e^{-d(u)\tau}}{1 - g(u)e^{-d(u)\tau}}.
\end{equation}

%%%%%%%%%%%%%%%%%%%%%%%%%%%%%%%%%%%%
\subsection{Int\'egration pour obtenir $C(\tau;u)$}

En int\'egrant~\eqref{eq:Cprime}, on obtient
\begin{equation}
C(\tau;u) = i u r\,\tau
+ \kappa\theta\int_0^\tau D(s;u)\,ds.
\end{equation}

L'int\'egrale en $s$ se calcule explicitement. Le r\'esultat standard (obtenu
en faisant le changement $y=e^{-d(u)s}$ et en int\'egrant terme \`a terme) est
\begin{equation}
C(\tau;u)
= i u r \tau
+ \frac{\kappa\theta}{\xi^2}
  \left[
     \left(\kappa - \rho\xi i u - d(u)\right)\tau
     - 2\ln\!\left(\frac{1 - g(u)e^{-d(u)\tau}}{1 - g(u)}\right)
  \right].
\end{equation}

%%%%%%%%%%%%%%%%%%%%%%%%%%%%%%%%%%%%
\subsection{Fonction caract\'eristique finale}

En r\'esum\'e, la fonction caract\'eristique du log--prix $X_T=\ln S_T$ sous
le mod\`ele de Heston est
\begin{equation}
\boxed{
\varphi_T(u)
= \mathbb{E}\!\left[ e^{i u X_T} \right]
= \exp\!\Big(
C(T;u) + D(T;u)\,v_0 + i u \ln S_0
\Big),
}
\end{equation}
avec
\begin{align}
d(u) &= \sqrt{(\rho\xi i u - \kappa)^2 + \xi^2(u^2+i u)},\\[3pt]
g(u) &= \frac{\kappa - \rho\xi i u - d(u)}{\kappa - \rho\xi i u + d(u)},\\[3pt]
D(T;u) &= \frac{\kappa - \rho\xi i u - d(u)}{\xi^2}\;
         \frac{1 - e^{-d(u)T}}{1 - g(u)e^{-d(u)T}},\\[3pt]
C(T;u) &= i u r T
+ \frac{\kappa\theta}{\xi^2}
  \left[
     \left(\kappa - \rho\xi i u - d(u)\right)T
     - 2\ln\!\left(\frac{1 - g(u)e^{-d(u)T}}{1 - g(u)}\right)
  \right].
\end{align}

%%%%%%%%%%%%%%%%%%%%%%%%%%%%%%%%%%%%%%%%%%%%%%%%%%%%%%%%%%%%
\subsection{Repr\'esentation de Fourier du prix de la call}

\subsection{Densit\'e du log--prix}

Notons $X_T=\ln S_T$ et $f_{X_T}$ sa densit\'e (sous $Q$). Le prix de la call
en fonction de $k=\ln K$ s'\'ecrit
\begin{equation}
C(k) = C(K,T) = e^{-rT}\int_k^{+\infty}\big(e^x-e^k\big)\,f_{X_T}(x)\,dx,
\qquad k:=\ln K.
\end{equation}

\subsection{Damping et transform\'ee de Fourier}

On choisit un param\`etre de damping $\alpha>0$ tel que
\begin{equation}
\mathbb{E}\big[e^{(\alpha+1)X_T}\big] < +\infty.
\end{equation}
On introduit la quantit\'e
\begin{equation}
c_\alpha(k) := e^{\alpha k}\,C(k),
\end{equation}
et sa transform\'ee de Fourier
\begin{equation}
\psi_T(v) := \int_{-\infty}^{+\infty} e^{ivk} c_\alpha(k)\,dk
           = \int_{-\infty}^{+\infty} e^{(\alpha+iv)k} C(k)\,dk,
           \qquad v\in\mathbb{R}.
\end{equation}

Par un calcul direct (changement d'ordre d'int\'egration et int\'egration en
$k$ d'abord), on obtient
\begin{equation}
\psi_T(v)
= e^{-rT}\,\frac{1}{(\alpha+iv)(\alpha+iv+1)}
      \int_{-\infty}^{+\infty} e^{(\alpha+1+iv)x}f_{X_T}(x)\,dx.
\end{equation}

Or
\begin{equation}
\int_{-\infty}^{+\infty} e^{(\alpha+1+iv)x}f_{X_T}(x)\,dx
= \mathbb{E}\big[e^{(\alpha+1+iv)X_T}\big]
= \varphi_T\big( v - i(\alpha+1)\big),
\end{equation}
puisque
$e^{(\alpha+1+iv)x} = e^{i(v-i(\alpha+1))x}$.

On en d\'eduit
\begin{equation}
\boxed{
\psi_T(v)
= e^{-rT}\,
  \frac{\varphi_T\big(v - i(\alpha+1)\big)}
       {(\alpha+iv)(\alpha+iv+1)}.
}
\end{equation}

\subsection{Formule d'inversion}

Par inversion de Fourier, on r\'ecup\`ere
\begin{equation}
c_\alpha(k) = \frac{1}{2\pi}
\int_{-\infty}^{+\infty} e^{-ivk}\,\psi_T(v)\,dv.
\end{equation}
Comme $C(k)$ est r\'eel, $c_\alpha(k)$ l'est aussi et l'int\'egrale se
simplifie en
\begin{equation}
c_\alpha(k) =
\frac{1}{\pi}
\int_0^{+\infty} \Re\!\bigl(e^{-ivk}\psi_T(v)\bigr)\,dv.
\end{equation}

En rappelant $C(k) = e^{-\alpha k} c_\alpha(k)$, on obtient finalement
\begin{equation}
\boxed{
C(K,T)
= C(k)
= \frac{e^{-\alpha k}}{\pi}
  \int_0^{+\infty}
  \Re\!\left(
    e^{-ivk}\,
    e^{-rT}\,
    \frac{\varphi_T\big(v - i(\alpha+1)\big)}
         {(\alpha+iv)(\alpha+iv+1)}
  \right)\,dv,
\qquad k=\ln K.
}
\end{equation}

%%%%%%%%%%%%%%%%%%%%%%%%%%%%%%%%%%%%%%%%%%%%%%%%%%%%%%%%%%%%
\subsection{Sch\'ema num\'erique (id\'ee g\'en\'erale)}

Pour calculer $C(K,T)$ num\'eriquement :

\begin{enumerate}
\item Choisir $\alpha>0$ tel que
      $\mathbb{E}[e^{(\alpha+1)X_T}]<\infty$ (en pratique $\alpha\in[1,2]$
      est standard).
\item Choisir une grille pour $v$: $v_j = j\Delta v$, $j=0,\dots,N$, avec
      $\Delta v$ et $N$ fixant la pr\'ecision.
\item Pour chaque $v_j$,
      \begin{enumerate}
      \item calculer $u_j := v_j - i(\alpha+1)$,
      \item calculer $d(u_j)$, $g(u_j)$, $D(T;u_j)$, $C(T;u_j)$,
      \item en d\'eduire $\varphi_T(u_j)$,
      \item calculer $\psi_T(v_j)$ via la formule ci--dessus.
      \end{enumerate}
\item Approcher l'int\'egrale par une somme de Riemann ou Simpson:
      \[
      C(K,T) \approx
      \frac{e^{-\alpha k}}{\pi}
      \sum_{j=0}^N
      \Re\!\left( e^{-i v_j k} \psi_T(v_j) \right)\Delta v.
      \]
\end{enumerate}



\subsection*{A.6. Formule finale et sch\'ema num\'erique}


\subsection*{Formule finale de Heston pour $C(K,T)$ et sch\'ema num\'erique}

\subsection*{1. Fonction caract\'eristique sous Heston}

Sous la mesure risque--neutre $Q$ :
\[
\begin{cases}
dS_t = r S_t\,dt + \sqrt{v_t}\,S_t\,dW_t^{(1)}, \\[3pt]
dv_t = \kappa(\theta - v_t)\,dt + \xi\sqrt{v_t}\,dW_t^{(2)}, \\[3pt]
dW_t^{(1)}\,dW_t^{(2)} = \rho\,dt,
\end{cases}
\qquad S_0>0,\; v_0>0.
\]

On pose
\[
X_t := \ln S_t .
\]

La fonction caract\'eristique du log--prix $X_T$ est
\[
\varphi_T(u) := \mathbb{E}\big[e^{iuX_T}\big], \qquad u\in\mathbb{C}.
\]

Elle s'\'ecrit sous forme affine
\[
\boxed{\displaystyle
\varphi_T(u)
= \exp\Big(C(T,u) + D(T,u)\,v_0 + iu\ln S_0\Big),
}
\]
avec
\[
d(u) := \sqrt{(\rho\xi iu - \kappa)^2 + \xi^2(u^2+iu)},
\qquad \Re(d(u))>0,
\]
\[
g(u) := \frac{\kappa - \rho\xi iu - d(u)}{\kappa - \rho\xi iu + d(u)}.
\]

Les fonctions $C$ et $D$ sont
\[
\boxed{\displaystyle
D(T,u)
= \frac{\kappa - \rho\xi iu - d(u)}{\xi^2}
  \,\frac{1 - e^{-d(u)T}}{1 - g(u)\,e^{-d(u)T}},
}
\]
\[
\boxed{\displaystyle
C(T,u)
= iurT
+ \frac{\kappa\theta}{\xi^2}
  \left[
    \big(\kappa - \rho\xi iu - d(u)\big)T
    - 2\ln\!\left(\frac{1 - g(u)e^{-d(u)T}}{1 - g(u)}\right)
  \right].
}
\]

Tous les param\`etres $(\kappa,\theta,\xi,\rho,v_0,r,S_0)$ apparaissent
explicitement dans $d(u)$, $g(u)$, $C(T,u)$, $D(T,u)$ donc dans $\varphi_T(u)$.

%%%%%%%%%%%%%%%%%%%%%%%%%%%%%%%%%%%%%%%%%%%%%%%%%%%%%%%%%%%%
\subsection*{2. Formule de prix via transform\'ee de Fourier (type Carr--Madan)}

On note $K>0$ le strike et
\[
k := \ln K .
\]

On fixe un param\`etre de damping $\alpha>0$ tel que
\[
\mathbb{E}\big[e^{(\alpha+1)X_T}\big] < \infty
\quad\text{(en pratique $\alpha\in[1,2]$).}
\]

Pour $v\in\mathbb{R}$, on pose
\[
A(v) := \alpha + iv.
\]

La transform\'ee de Carr--Madan est
\[
\psi_T(v)
:= \int_{\mathbb{R}} e^{A(v) k'} C(k')\,dk',
\]
et peut se r\'ecrire en fonction de $\varphi_T$ :
\[
\boxed{\displaystyle
\psi_T(v)
= \frac{e^{-rT}}{A(v)\big(A(v)+1\big)}\,
  \varphi_T\big(v - i(\alpha+1)\big).
}
\]

La formule d'inversion donne alors le prix de la call :
\[
\boxed{\displaystyle
C(K,T)
= C(k)
= \frac{e^{-\alpha k}}{\pi}
  \int_0^{+\infty}
    \Re\!\left(
      e^{-ivk}\,\psi_T(v)
    \right)\,dv.
}
\]

En combinant explicitement :
\[
\boxed{\displaystyle
C(K,T)
=
\frac{e^{-\alpha k}}{\pi}
\int_0^{+\infty}
\Re\!\left[
e^{-ivk}\,
\frac{e^{-rT}}{(\alpha+iv)(\alpha+1+iv)}\,
\exp\Big(
C(T,u) + D(T,u)\,v_0 + iu\ln S_0
\Big)
\right]_{u=v-i(\alpha+1)}
dv.
}
\]

Tout est maintenant en fonction de
\[
(\kappa,\theta,\xi,\rho,v_0,r,S_0,K,T,\alpha).
\]

%%%%%%%%%%%%%%%%%%%%%%%%%%%%%%%%%%%%%%%%%%%%%%%%%%%%%%%%%%%%
\subsection*{3. Principes num\'eriques (sans pseudo-code)}

\paragraph{Rappel de la formule.}
Dans le cadre de Carr--Madan, on utilise un call damp\'e et l'on \'evalue une int\'egrale oscillante sur $v\in\R_+$ :
\[
C(K,T)
=
\frac{\ee^{-\alpha k}}{\pi}\int_0^{+\infty}\Re\!\left(\ee^{-\ii v k}\,\psi_T(v)\right)\dd v,
\qquad
k=\ln K,
\]
o\`u $\psi_T(v)$ d\'epend de la fonction caract\'eristique $\varphi_T(u)$ du log-prix \`a maturit\'e $T$ \'evalu\'ee en
\[
u(v)=v-\ii(\alpha+1).
\]

\paragraph{Choix du damping et int\'egrabilit\'e.}
Le param\`etre $\alpha>0$ assure l'int\'egrabilit\'e de la transform\'ee (contr\^ole des queues en $K$). Le choix est un compromis :
un $\alpha$ trop faible peut conduire \`a une int\'egrale mal conditionn\'ee ; un $\alpha$ trop grand amplifie certaines erreurs num\'eriques.

\paragraph{Troncature, discr\'etisation et diagnostics.}
Num\'eriquement, l'int\'egrale est tronqu\'ee \`a $[0,V_{\max}]$ et approch\'ee par quadrature sur une grille $v_n=n\,\Delta v$.
Les erreurs proviennent (i) de la troncature (choix de $V_{\max}$), (ii) de la discr\'etisation ($\Delta v$), et (iii) de l'oscillation.
Une validation th\'eorique consiste \`a v\'erifier la stabilit\'e du prix lorsque l'on augmente $V_{\max}$ et que l'on raffine $\Delta v$,
et \`a croiser avec une m\'ethode ind\'ependante (Monte Carlo ou PDE) sur quelques points de grille.

\paragraph{Branche complexe et stabilit\'e.}
Dans Heston (et Bates), l'\'evaluation de $\varphi_T(u)$ requiert des racines carr\'ees et logarithmes complexes. Le choix de branche
et la continuit\'e en $u$ conditionnent la stabilit\'e num\'erique ; en pratique, on impose une convention coh\'erente sur la racine
et l'on surveille les discontinuit\'es.

\paragraph{Version FFT.}
La version FFT de Carr--Madan s'appuie sur une grille uniforme en fr\'equence $v$ et une grille uniforme en log-strike $k$,
avec une relation de type $\Delta v\,\Delta k = 2\pi/N$. Elle permet d'obtenir un ensemble de prix sur une grille de strikes en un seul
calcul, au prix d'effets d'aliasing et de fen\^etrage \`a contr\^oler (choix du domaine et des pas).

\iffalse % Annexes B--C d\'etaillent des sch\'emas num\'eriques ; retir\'ees pour respecter ``z\'ero pseudo-code''
\subsection*{Annexe B - Mouvement brownien fractionnaire (Davies--Harte)}
\addcontentsline{toc}{section}{Annexe B - fBM et Rough Volatility}



\subsection*{Simulation du mouvement brownien fractionnaire — Méthode de Davies--Harte (dérivation détaillée)}

%%%%%%%%%%%%%%%%%%%%%%%%%%%%%%%%%%%%%%%%%%%%%%%%%%%%%%%%%%%%%%%%%%%%%%%%%%%%%%%
\subsection{Rappels sur le mouvement brownien fractionnaire}

\subsection{Définition}

Pour un paramètre de Hurst $H\in(0,1)$, le \emph{mouvement brownien fractionnaire}
(fBm) $(\bH(t))_{t\ge 0}$ est un processus gaussien centré tel que
\begin{equation}
  \E[\bH(t)] = 0,\qquad t\ge 0,
\end{equation}
et
\begin{equation}
  \Cov\bigl(\bH(t),\bH(s)\bigr)
  = R_H(t,s)
  = \frac{1}{2}\bigl(|t|^{2H} + |s|^{2H} - |t-s|^{2H}\bigr),
  \qquad s,t\ge 0.
\end{equation}
On en déduit deux propriétés clés :
\begin{itemize}
  \item \textbf{incréments stationnaires} :
  \[
    \bH(t+h)-\bH(t)
    \stackrel{d}{=}
    \bH(h)-\bH(0), \qquad \forall t,h\ge 0 ;
  \]
  \item \textbf{autosimilarité} :
  \[
    \bigl(\bH(ct)\bigr)_{t\ge 0}
    \stackrel{d}{=}
    \bigl(c^H \bH(t)\bigr)_{t\ge 0},\qquad \forall c>0.
  \]
\end{itemize}
En particulier,
\begin{equation}
  \Var\bigl(\bH(t)\bigr) = t^{2H}.
\end{equation}

%%%%%%%%%%%%%%%%%%%%%%%%%%%%%%%%%%%%%%%%%%%%%%%%%%%%%%%%%%%%%%%%%%%%%%%%%%%%%%%
\subsection{Discrétisation sur une grille}

On veut simuler une trajectoire de $\bH$ sur une période $[0,T]$.
On fixe un nombre de pas $N\in\mathbb{N}^\ast$ et on définit
\begin{equation}
  \Delta t = \frac{T}{N},\qquad
  t_k = k\,\Delta t,\quad k = 0,\dots,N.
\end{equation}

L'autosimilarité permet d'écrire, pour chaque $k$,
\begin{equation}
  \bH(t_k)
  = \bH(k\Delta t)
  \stackrel{d}{=} (\Delta t)^H \widetilde B_H(k),
\end{equation}
où $(\widetilde B_H(k))_{k=0,\dots,N}$ est un fBm vu aux instants entiers $k$.
Ainsi, on peut :
\begin{enumerate}
  \item simuler un fBm \emph{à pas unitaire} $(\widetilde B_H(0),\dots,\widetilde B_H(N))$ ;
  \item le rescaler ensuite par $(\Delta t)^H$ :
  \[
    \bH(t_k) \approx (\Delta t)^H \widetilde B_H(k).
  \]
\end{enumerate}

On note les incréments unitaires
\begin{equation}
  X_k = \widetilde B_H(k+1) - \widetilde B_H(k),\qquad k=0,\dots,N-1.
\end{equation}
La suite $(X_k)_{k\ge 0}$ est le \emph{fractional Gaussian noise} (fGn).
On a alors
\begin{equation}
  \widetilde B_H(k)
  = \sum_{j=0}^{k-1} X_j,\qquad k\ge 1,
\end{equation}
et
\begin{equation}
  \bH(t_k)
  \approx (\Delta t)^H \sum_{j=0}^{k-1} X_j.
\end{equation}

Toute la difficulté se ramène donc à simuler le vecteur gaussien corrélé
\[
  (X_0,\dots,X_{N-1})
\]
avec la bonne structure de covariance.

%%%%%%%%%%%%%%%%%%%%%%%%%%%%%%%%%%%%%%%%%%%%%%%%%%%%%%%%%%%%%%%%%%%%%%%%%%%%%%%
\subsection{Covariance du bruit fractionnaire}

\subsection{Dérivation explicite de $\gamma(k)$}

On veut la covariance
\begin{equation}
  \gamma(k)
  = \Cov(X_n,X_{n+k}),
\end{equation}
qui, par stationnarité des incréments, ne dépend pas de $n$.

On se place sur le fBm à pas unitaire $\widetilde B_H$ et on note
\begin{equation}
  X_n = \widetilde B_H(n+1)-\widetilde B_H(n).
\end{equation}
En développant la covariance à l'aide de la bilinéarité :
\begin{align}
  \gamma(k)
  &= \Cov\bigl(\widetilde B_H(n+1)-\widetilde B_H(n),
               \widetilde B_H(n+k+1)-\widetilde B_H(n+k)\bigr)
  \\[0.3em]
  &= \Cov\bigl(\widetilde B_H(n+1),\widetilde B_H(n+k+1)\bigr)
   - \Cov\bigl(\widetilde B_H(n+1),\widetilde B_H(n+k)\bigr)
  \notag\\
  &\quad
   - \Cov\bigl(\widetilde B_H(n),\widetilde B_H(n+k+1)\bigr)
   + \Cov\bigl(\widetilde B_H(n),\widetilde B_H(n+k)\bigr).
  \label{eq:gamma-dev}
\end{align}
Par stationnarité des accroissements et invariance par translation,
on peut prendre $n=0$ sans perte de généralité :
\begin{equation}
  \gamma(k)
  = \Cov\bigl(\widetilde B_H(1)-\widetilde B_H(0),
               \widetilde B_H(k+1)-\widetilde B_H(k)\bigr).
\end{equation}
Comme $\widetilde B_H(0)=0$ presque sûrement, on a
\[
  \Cov\bigl(\widetilde B_H(0),\widetilde B_H(t)\bigr)=0,\qquad \forall t\ge 0,
\]
et \eqref{eq:gamma-dev} se simplifie en
\begin{equation}
  \gamma(k)
  = R_H(1,k+1) - R_H(1,k),
\end{equation}
où
\[
  R_H(s,t)
  = \frac{1}{2}\Bigl(|s|^{2H} + |t|^{2H} - |t-s|^{2H}\Bigr).
\]

On calcule explicitement :
\begin{align}
  R_H(1,k+1)
  &= \frac{1}{2}\Bigl(|1|^{2H} + |k+1|^{2H} - |(k+1)-1|^{2H}\Bigr)
   = \frac{1}{2}\Bigl(1 + |k+1|^{2H} - |k|^{2H}\Bigr),
\\
  R_H(1,k)
  &= \frac{1}{2}\Bigl(|1|^{2H} + |k|^{2H} - |k-1|^{2H}\Bigr)
   = \frac{1}{2}\Bigl(1 + |k|^{2H} - |k-1|^{2H}\Bigr).
\end{align}
D'où
\begin{align}
  \gamma(k)
  &= R_H(1,k+1) - R_H(1,k)
\\[0.3em]
  &= \frac{1}{2}\Bigl(1 + |k+1|^{2H} - |k|^{2H}
                      - 1 - |k|^{2H} + |k-1|^{2H}\Bigr)
\\[0.3em]
  &= \frac{1}{2}\Bigl(|k+1|^{2H} + |k-1|^{2H} - 2|k|^{2H}\Bigr).
\end{align}

Ainsi, la covariance théorique du bruit fractionnaire (fGn) est
\begin{equation}
  \boxed{
  \gamma(k)
  = \Cov(X_t,X_{t+k})
  = \E[X_t X_{t+k}]
  = \frac{1}{2}\Bigl(|k+1|^{2H} + |k-1|^{2H} - 2|k|^{2H}\Bigr).
  }
\end{equation}

\subsection{Matrice de covariance Toeplitz}

On regroupe les $N$ premiers incréments dans le vecteur
\[
  X = (X_0,\dots,X_{N-1})^\top.
\]
Sa matrice de covariance est la matrice Toeplitz
\begin{equation}
  \Gamma
  = \Cov(X,X^\top)
  = \bigl(\gamma(|i-j|)\bigr)_{0\le i,j\le N-1}.
\end{equation}
Autrement dit, la première ligne de $\Gamma$ est
\[
  (\gamma(0),\gamma(1),\dots,\gamma(N-1)),
\]
et chaque ligne suivante est une translation vers la droite.

La méthode de Davies--Harte consiste à \emph{embeder} cette matrice Toeplitz
dans une matrice \textbf{circulante} de taille $2N\times 2N$, ce qui
permet d'utiliser la FFT.

%%%%%%%%%%%%%%%%%%%%%%%%%%%%%%%%%%%%%%%%%%%%%%%%%%%%%%%%%%%%%%%%%%%%%%%%%%%%%%%
\subsection{Embedding circulant et vecteur $c$}

On construit le vecteur
\begin{equation}
  c
  = \bigl(
    \gamma(0),\gamma(1),\dots,\gamma(N),\gamma(N-1),\dots,\gamma(1)
    \bigr)^\top
  \in\R^{2N}.
\end{equation}
On en déduit la matrice circulante $C\in\R^{2N\times 2N}$ définie par
\begin{equation}
  C_{jk} = c_{(k-j)\bmod 2N},
  \qquad j,k = 0,\dots,2N-1.
\end{equation}
Ainsi :
\begin{itemize}
  \item la première ligne de $C$ est exactement $c^\top$ ;
  \item la $(j+1)$-ème ligne est obtenue en décalant circulairement la
        précédente d'un cran vers la droite.
\end{itemize}

Par construction, le bloc $N\times N$ en haut à gauche de $C$ coïncide
avec la matrice Toeplitz $\Gamma$.  
L'idée est alors :
\begin{quote}
  simuler un vecteur gaussien $Y\in\R^{2N}$ de covariance $C$,
  puis ne garder que ses $N$ premières composantes ; celles-ci auront
  la covariance désirée $\Gamma$.
\end{quote}

%%%%%%%%%%%%%%%%%%%%%%%%%%%%%%%%%%%%%%%%%%%%%%%%%%%%%%%%%%%%%%%%%%%%%%%%%%%%%%%
\subsection{Diagonalisation de la matrice circulante par la FFT}

\subsection{Matrice de Fourier}

On pose $M=2N$ et on considère la matrice de Fourier \emph{unitaire}
\begin{equation}
  F_{pj} = \frac{1}{\sqrt{M}} e^{-2\pi i pj/M},
  \qquad p,j = 0,\dots,M-1.
\end{equation}
Cette normalisation donne
\[
  F^{-1} = F^\ast,
\]
la transposée conjuguée.

Pour un vecteur $v\in\C^M$, on définit sa transformée de Fourier discrète
(DFT) par
\begin{equation}
  \widehat v = F v,
\end{equation}
et la transformée inverse (IDFT) par
\begin{equation}
  v = F^\ast \widehat v.
\end{equation}
En pratique numérique, $F$ et $F^\ast$ sont calculés via la FFT/IFFT.

\subsection{Diagonalisation d'une matrice circulante}

Un résultat standard est : \emph{toute matrice circulante est diagonalisée
par la matrice de Fourier}.  
Plus précisément, si $C$ est la matrice circulante dont la première ligne
est $c^\top$, alors
\begin{equation}
  \widehat c = F c
  = (\lambda_0,\dots,\lambda_{M-1})^\top
\end{equation}
contient les valeurs propres de $C$, et on a
\begin{equation}
  \boxed{
  C = F^\ast \,\Delta\, F,
  }
\end{equation}
où
\begin{equation}
  \Delta = \mathrm{diag}(\lambda_0,\dots,\lambda_{M-1})
\end{equation}
est diagonale.

Dans le langage numérique, on écrit simplement
\[
  \Lambda = \mathrm{FFT}(c),
\]
où $\Lambda=(\lambda_k)_{0\le k<M}$.

Comme $c$ est réel et symétrique, les $\lambda_k$ sont réels et,
dans le cas idéal, non négatifs.  
En pratique, la méthode de Davies--Harte requiert $\lambda_k\ge 0$ ;
si certains $\lambda_k$ sont négatifs (erreurs numériques mises à part),
l'algorithme n'est pas applicable pour ce couple $(H,N)$.

%%%%%%%%%%%%%%%%%%%%%%%%%%%%%%%%%%%%%%%%%%%%%%%%%%%%%%%%%%%%%%%%%%%%%%%%%%%%%%%
\subsection{Simulation gaussienne via la représentation spectrale}

\subsection{Principe général}

Soit $Y\in\C^M$ un vecteur gaussien centré.  
Si l'on peut écrire la covariance sous la forme
\[
  \Cov(Y,Y^\ast) = C = F^\ast \Delta F,
\]
alors il suffit de prendre
\begin{equation}
  Y = F^\ast Z,
\end{equation}
où $Z$ est un vecteur gaussien complexe tel que
\[
  \E[Z]=0,\qquad \E[Z Z^\ast] = \Delta.
\]
En effet,
\begin{align}
  \Cov(Y,Y^\ast)
  &= \E[YY^\ast]
   = \E[F^\ast Z Z^\ast F]
   = F^\ast \E[Z Z^\ast] F
   = F^\ast \Delta F
   = C.
\end{align}

On a donc transformé un problème de covariance générale $C$
en un problème diagonal beaucoup plus simple.

\subsection{Construction du vecteur $Z$}

On veut un vecteur
\[
  Z = (Z_0,\dots,Z_{M-1})^\top \in\C^M
\]
tel que
\[
  \E[Z]=0,\qquad \E[Z Z^\ast] = \Delta
  = \mathrm{diag}(\lambda_0,\dots,\lambda_{M-1}).
\]
Comme $\Delta$ est diagonale, cela revient à imposer que
\[
  \E[Z_k \overline{Z_\ell}] = 0\quad\text{si }k\ne\ell,
  \qquad
  \E[|Z_k|^2] = \lambda_k.
\]

On tient compte du fait que l'on veut finalement un \emph{vecteur réel}
$Y$, donc on impose en plus la symétrie conjuguée
\begin{equation}
  Z_{M-k} = \overline{Z_k},\qquad 1\le k\le M-1.
\end{equation}
Comme $c$ est réel et symétrique, on a $\lambda_k = \lambda_{M-k}$,
ce qui rend cette construction cohérente.

Une construction explicite standard est :
\begin{itemize}
  \item Générer des gaussiennes réelles indépendantes
        $U_0,\dots,U_N,V_1,\dots,V_{N-1}\sim\N(0,1)$.
  \item Poser
    \begin{align}
      Z_0   &= \sqrt{\lambda_0}\,U_0,\\
      Z_N   &= \sqrt{\lambda_N}\,U_N,\\
      Z_k   &= \sqrt{\frac{\lambda_k}{2}}\,(U_k + iV_k),
              &&k=1,\dots,N-1,\\
      Z_{M-k} &= \sqrt{\frac{\lambda_k}{2}}\,(U_k - iV_k)
              = \overline{Z_k},
              &&k=1,\dots,N-1.
    \end{align}
\end{itemize}
On vérifie alors :
\[
  \E[Z_k] = 0,\qquad
  \E[|Z_k|^2] = \lambda_k,\qquad
  \E[Z_k\overline{Z_\ell}] = 0\ (k\ne\ell),
\]
donc $\E[Z Z^\ast] = \Delta$.

%%%%%%%%%%%%%%%%%%%%%%%%%%%%%%%%%%%%%%%%%%%%%%%%%%%%%%%%%%%%%%%%%%%%%%%%%%%%%%%
\subsection{Génération du bruit fractionnaire par IFFT}

\subsection{Application de l’IFFT}

On définit
\begin{equation}
  Y = F^\ast Z = \mathrm{IFFT}(Z).
\end{equation}
Par ce qui précède, on a
\begin{equation}
  \Cov(Y,Y^\ast) = C.
\end{equation}
De plus, grâce aux relations $Z_{M-k}=\overline{Z_k}$, on montre que
$Y$ est \emph{réel} (à des erreurs numériques près) :
\[
  Y \in \R^M.
\]
Dans l'implémentation, on prend donc
\[
  Y = \Re\bigl(\mathrm{IFFT}(Z)\bigr).
\]

On note
\[
  Y = (Y_0,\dots,Y_{M-1})^\top.
\]
Par construction, la covariance de $Y$ est la matrice circulante $C$.
Le bloc $N\times N$ en haut à gauche de $C$ est la matrice de covariance
$\Gamma$ du fGn.  
Donc, en ne conservant que les $N$ premières composantes de $Y$,
on obtient un vecteur ayant la bonne covariance.

On définit donc le bruit fractionnaire simulé par
\begin{equation}
  (X_0,X_1,\dots,X_{N-1})^\top
  =
  (Y_0,Y_1,\dots,Y_{N-1})^\top
  \approx \text{fGn de covariance }\gamma(k).
\end{equation}

%%%%%%%%%%%%%%%%%%%%%%%%%%%%%%%%%%%%%%%%%%%%%%%%%%%%%%%%%%%%%%%%%%%%%%%%%%%%%%%
\subsection{Reconstruction du mouvement brownien fractionnaire}

On part de
\[
  \widetilde B_H(0) = 0,\qquad
  \widetilde B_H(k) = \sum_{j=0}^{k-1} X_j,\quad k=1,\dots,N,
\]
où $(X_j)$ est le fGn simulé ci-dessus.

Pour revenir aux temps $t_k = k\Delta t$ sur $[0,T]$, on utilise
l'autosimilarité :
\[
  \bH(t_k) \stackrel{d}{=} (\Delta t)^H \widetilde B_H(k).
\]
On pose donc
\begin{equation}
  \boxed{
  \bH(t_k) = (\Delta t)^H \sum_{j=0}^{k-1} X_j,
  \qquad k=0,\dots,N,
  }
\end{equation}
avec la convention $\bH(t_0)=\bH(0)=0$.

On vérifie que, au moins au premier ordre, la variance est correctement
ajustée :
\begin{align}
  \Var\bigl(\bH(t_k)\bigr)
  &= (\Delta t)^{2H}\,\Var\Bigl(\sum_{j=0}^{k-1} X_j\Bigr)
\\[0.3em]
  &= (\Delta t)^{2H}
     \sum_{i=0}^{k-1}\sum_{j=0}^{k-1} \gamma(|i-j|)
\\[0.3em]
  &\approx (k\Delta t)^{2H} = t_k^{2H},
\end{align}
ce qui est cohérent avec la définition de la fBm.

%%%%%%%%%%%%%%%%%%%%%%%%%%%%%%%%%%%%%%%%%%%%%%%%%%%%%%%%%%%%%%%%%%%%%%%%%%%%%%%
\subsection{Récapitulatif de la méthode de Davies--Harte}

\subsection*{Étape 1 : covariance du bruit fractionnaire}
\begin{itemize}
  \item Pour $k=0,\dots,N$, calculer
    \[
      \gamma(k) = \frac{1}{2}\Bigl(|k+1|^{2H} + |k-1|^{2H} - 2|k|^{2H}\Bigr).
    \]
\end{itemize}

\subsection*{Étape 2 : vecteur circulant $c$}
\begin{itemize}
  \item Construire
    \[
      c = \bigl(
        \gamma(0),\gamma(1),\dots,\gamma(N),
        \gamma(N-1),\dots,\gamma(1)
      \bigr)^\top \in\R^{2N}.
    \]
\end{itemize}

\subsection*{Étape 3 : FFT et valeurs propres}
\begin{itemize}
  \item Calculer la FFT :
    \[
      \Lambda = \mathrm{FFT}(c)
              = (\lambda_0,\dots,\lambda_{2N-1})^\top.
    \]
  \item Vérifier que toutes les $\lambda_k$ sont non négatives
        (à un petit $\varepsilon$ près pour les erreurs numériques).
\end{itemize}

\subsection*{Étape 4 : génération du vecteur gaussien complexe}
\begin{itemize}
  \item Générer des $U_0,\dots,U_N,V_1,\dots,V_{N-1}\sim\N(0,1)$ indépendants.
  \item Construire $Z\in\C^{2N}$ par
    \begin{align*}
      Z_0   &= \sqrt{\lambda_0}\,U_0,\\
      Z_N   &= \sqrt{\lambda_N}\,U_N,\\
      Z_k   &= \sqrt{\frac{\lambda_k}{2}}\,(U_k + iV_k),
              &&k=1,\dots,N-1,\\
      Z_{2N-k} &= \sqrt{\frac{\lambda_k}{2}}\,(U_k - iV_k)
              = \overline{Z_k},
              &&k=1,\dots,N-1.
    \end{align*}
\end{itemize}

\subsection*{Étape 5 : IFFT et extraction du fGn}
\begin{itemize}
  \item Calculer
    \[
      Y = \mathrm{IFFT}(Z) = F^\ast Z.
    \]
  \item Prendre la partie réelle et garder les $N$ premières composantes :
    \[
      X_k = \Re(Y_k),\qquad k=0,\dots,N-1.
    \]
    Le vecteur $(X_0,\dots,X_{N-1})$ est le \emph{fractional Gaussian noise}.
\end{itemize}

\subsection*{Étape 6 : reconstruction du fBm sur $[0,T]$}
\begin{itemize}
  \item Fixer $\Delta t = T/N$.
  \item Poser
    \[
      \bH(t_0)=0,\qquad
      \bH(t_k) = (\Delta t)^H \sum_{j=0}^{k-1} X_j,\quad k=1,\dots,N.
    \]
\end{itemize}

On obtient ainsi une trajectoire discrète
\[
  \bigl(\bH(t_0),\bH(t_1),\dots,\bH(t_N)\bigr)
\]
qui suit, en loi, un mouvement brownien fractionnaire de paramètre $H$
sur l'intervalle $[0,T]$.



\subsection*{Annexe C - SDE de Volterra et rough Bergomi}
\addcontentsline{toc}{section}{Annexe C - SDE de Volterra}



\subsection*{Simulation du mouvement brownien fractionnaire — Méthode de Davies--Harte et lien avec le modèle rough Bergomi}

Dans tout ce qui suit, $H\in(0,1)$ désigne le paramètre de Hurst.
On note $(\bH(t))_{t\ge 0}$ un mouvement brownien fractionnaire
normalisé, c'est-à-dire un processus gaussien centré tel que, pour
$t,s\ge 0$,
\begin{equation}
  \label{eq:cov-fBm}
  \Cov\bigl(\bH(t),\bH(s)\bigr)
  = R_H(t,s)
  = \frac{1}{2}\Bigl(t^{2H} + s^{2H} - |t-s|^{2H}\Bigr).
\end{equation}
On veut simuler numériquement la trajectoire
$(\bH(t_k))_{k=0,\dots,N}$ sur une période $[0,T]$.

%-------------------------------------------------------------
\subsection*{I. Discrétisation du fBm et définition du bruit fractionnaire}

On fixe un nombre de pas $N\in\mathbb{N}^\ast$ et on définit
\begin{equation}
  \Delta t = \frac{T}{N},
  \qquad
  t_k = k\,\Delta t,\quad k=0,\dots,N.
\end{equation}

On introduit les \emph{incréments} (ou \emph{fractional Gaussian
noise} / fGn) par
\begin{equation}
  X_k = \bH(t_{k+1}) - \bH(t_k),
  \qquad k = 0,\dots,N-1.
\end{equation}
Les $X_k$ forment une suite gaussienne centrée, stationnaire, mais
corrélée (sauf si $H=\tfrac12$, cas brownien classique).

La fBm discrète est alors reconstruite par sommation :
\begin{equation}
  \bH(t_k)
  = \sum_{j=0}^{k-1} X_j,
  \qquad k=0,\dots,N,
\end{equation}
avec la convention $\bH(t_0)=\bH(0)=0$.

En pratique on simulera d'abord le bruit fractionnaire $(X_0,\dots,X_{N-1})$, puis
on cumulera pour obtenir $\bH$.

%-------------------------------------------------------------
\subsection*{II. Calcul détaillé de la covariance du bruit fractionnaire}

On veut obtenir la covariance
\[
  \gamma(m) = \Cov(X_0,X_m),
  \qquad m\in\mathbb{Z}.
\]
En utilisant la définition des incréments,
\[
  X_0 = \bH(\Delta t) - \bH(0),
  \qquad
  X_m = \bH((m+1)\Delta t) - \bH(m\Delta t).
\]

Par bilinéarité de la covariance,
\begin{align}
  \gamma(m)
  &= \Cov\bigl(\bH(\Delta t)-\bH(0),
               \bH((m+1)\Delta t)-\bH(m\Delta t)\bigr) \\[0.3em]
  &= \Cov\bigl(\bH(\Delta t),\bH((m+1)\Delta t)\bigr)
   - \Cov\bigl(\bH(\Delta t),\bH(m\Delta t)\bigr) \nonumber\\
  &\quad - \Cov\bigl(\bH(0),\bH((m+1)\Delta t)\bigr)
   + \Cov\bigl(\bH(0),\bH(m\Delta t)\bigr). \label{eq:gamma-expand}
\end{align}

On utilise la formule de covariance \eqref{eq:cov-fBm}. On traite
chaque terme de \eqref{eq:gamma-expand}.

\paragraph{1. Premier terme.}
\begin{align}
  \Cov\bigl(\bH(\Delta t),\bH((m+1)\Delta t)\bigr)
  &= \frac{1}{2}\Bigl[(\Delta t)^{2H}
  + ((m+1)\Delta t)^{2H}
  - |(m+1)\Delta t - \Delta t|^{2H}\Bigr] \\
  &= \frac{1}{2}\Bigl[(\Delta t)^{2H}
  + ((m+1)\Delta t)^{2H}
  - (m\Delta t)^{2H}\Bigr].
\end{align}

\paragraph{2. Deuxième terme.}
\begin{align}
  \Cov\bigl(\bH(\Delta t),\bH(m\Delta t)\bigr)
  &= \frac{1}{2}\Bigl[(\Delta t)^{2H}
  + (m\Delta t)^{2H}
  - |m\Delta t - \Delta t|^{2H}\Bigr] \\
  &= \frac{1}{2}\Bigl[(\Delta t)^{2H}
  + (m\Delta t)^{2H}
  - ((m-1)\Delta t)^{2H}\Bigr].
\end{align}

\paragraph{3. Troisième terme.}
\begin{align}
  \Cov\bigl(\bH(0),\bH((m+1)\Delta t)\bigr)
  &= \frac{1}{2}\Bigl[0^{2H}
  + ((m+1)\Delta t)^{2H}
  - |(m+1)\Delta t|^{2H}\Bigr] \\
  &= 0,
\end{align}
car les deux derniers termes se compensent.

\paragraph{4. Quatrième terme.}
\begin{align}
  \Cov\bigl(\bH(0),\bH(m\Delta t)\bigr)
  &= \frac{1}{2}\Bigl[0^{2H}
  + (m\Delta t)^{2H}
  - |m\Delta t|^{2H}\Bigr] \\
  &= 0.
\end{align}

En reportant ces expressions dans \eqref{eq:gamma-expand},
\begin{align}
  \gamma(m)
  &= \frac{1}{2}\Bigl[(\Delta t)^{2H}
  + ((m+1)\Delta t)^{2H}
  - (m\Delta t)^{2H}\Bigr] \\
  &\quad - \frac{1}{2}\Bigl[(\Delta t)^{2H}
  + (m\Delta t)^{2H}
  - ((m-1)\Delta t)^{2H}\Bigr] \\[0.4em]
  &= \frac{1}{2}\Bigl[(\Delta t)^{2H}
                      + ((m+1)\Delta t)^{2H}
                      - (m\Delta t)^{2H} \\
  &\qquad\qquad
        - (\Delta t)^{2H}
        - (m\Delta t)^{2H}
        + ((m-1)\Delta t)^{2H}\Bigr] \\[0.4em]
  &= \frac{1}{2}\Bigl[((m+1)\Delta t)^{2H}
        - 2(m\Delta t)^{2H}
        + ((m-1)\Delta t)^{2H}\Bigr].
\end{align}

On factorise $(\Delta t)^{2H}$ :
\begin{equation}
  \gamma(m)
  = \frac{(\Delta t)^{2H}}{2}
    \Bigl(|m+1|^{2H} - 2|m|^{2H} + |m-1|^{2H}\Bigr),
  \qquad m\in\mathbb{Z}.
\end{equation}

Pour alléger les notations, on introduit la covariance \emph{normalisée}
(associée au pas unité) :
\begin{equation}
  \gamma_0(m)
  = \frac{1}{2}\Bigl(|m+1|^{2H} - 2|m|^{2H} + |m-1|^{2H}\Bigr),
\end{equation}
de sorte que
\begin{equation}
  \gamma(m) = (\Delta t)^{2H}\,\gamma_0(m).
\end{equation}

La suite $(\gamma_0(m))_{m\in\mathbb{Z}}$ est la covariance
théorique du bruit fractionnaire unit-temps.

%-------------------------------------------------------------
\subsection*{III. Construction du vecteur de covariance circulante}

La covariance des $N$ premiers incréments $X_0,\dots,X_{N-1}$
est la matrice Toeplitz
\begin{equation}
  \Gamma_N =
  \begin{pmatrix}
    \gamma_0(0)   & \gamma_0(1)   & \cdots & \gamma_0(N-1) \\
    \gamma_0(1)   & \gamma_0(0)   & \cdots & \gamma_0(N-2) \\
    \vdots        & \vdots        & \ddots & \vdots        \\
    \gamma_0(N-1) & \gamma_0(N-2) & \cdots & \gamma_0(0)
  \end{pmatrix}.
\end{equation}
La méthode de Davies--Harte consiste à \emph{encastrer} cette
matrice Toeplitz dans une matrice \emph{circulante} de taille
$M=2N$, ce qui permet d'utiliser la FFT.

On définit le vecteur de longueur $M=2N$ :
\begin{equation}
  c =
  \bigl(
  c_0,c_1,\dots,c_{M-1}
  \bigr)
  =
  \bigl(
    \gamma_0(0),\gamma_0(1),\dots,
    \gamma_0(N-1),\gamma_0(N),\gamma_0(N-1),\dots,\gamma_0(1)
  \bigr).
\end{equation}
C'est un vecteur \emph{pair} au sens où $c_j=c_{M-j}$ pour
$j=1,\dots,M-1$.

On construit ensuite la matrice circulante
\begin{equation}
  \mathsf{C} =
  \bigl(c_{(j-k)\bmod M}\bigr)_{0\le j,k\le M-1}
  \in\R^{M\times M}.
\end{equation}
Explicitement,
\[
  \mathsf{C} =
  \begin{pmatrix}
    c_0   & c_1   & c_2   & \cdots & c_{M-1} \\
    c_{M-1} & c_0 & c_1   & \cdots & c_{M-2} \\
    \vdots  & \vdots & \vdots & \ddots & \vdots \\
    c_1  & c_2  & c_3  & \cdots & c_0
  \end{pmatrix}.
\]
Le sous-bloc supérieur gauche de taille $N\times N$ est
exactement la matrice de covariance $\Gamma_N$ des $N$
premiers incréments, grâce à la définition de $c$.

%-------------------------------------------------------------
\subsection*{IV. Diagonalisation de la matrice circulante par FFT}

On pose $M=2N$ et on considère la matrice de Fourier discrète
normalisée (unitaire) :
\begin{equation}
  F_{jk}
  = \frac{1}{\sqrt{M}}\,
    e^{-2i\pi jk / M},
  \qquad j,k=0,\dots,M-1.
\end{equation}
On a alors $F^{-1} = F^{\ast}$ (conjuguée-transposée).

Propriété fondamentale : toute matrice circulante $\mathsf{C}$ se
diagonalise dans la base de Fourier :
\begin{equation}
  \mathsf{C}
  = F^{\ast}\,\Delta\,F,
\end{equation}
où $\Delta$ est diagonale. Les valeurs propres
$(\lambda_0,\dots,\lambda_{M-1})$ sont données par la transformée
de Fourier du premier rang de $\mathsf{C}$, c'est-à-dire
du vecteur $c$ :
\begin{equation}
  \lambda_k
  = \sum_{j=0}^{M-1}
      c_j\,e^{-2i\pi jk / M},
  \qquad k=0,\dots,M-1.
\end{equation}
En notation compacte,
\begin{equation}
  \bm{\lambda}
  =
  (\lambda_0,\dots,\lambda_{M-1})^\top
  = \mathrm{FFT}(c),
\end{equation}
et
\begin{equation}
  \Delta = \mathrm{diag}(\lambda_0,\dots,\lambda_{M-1}).
\end{equation}

Puisque $c$ est réel et pair, tous les $\lambda_k$ sont réels.
Pour que la méthode soit valide, il faut en plus
$\lambda_k\ge 0$ pour tout $k$ (ce qui est garanti pour la fBm
dans un certain domaine de paramètres; en pratique on vérifie
numériquement que les valeurs propres sont non négatives).

On obtient donc la factorisation
\begin{equation}
  \mathsf{C} = F^{\ast}\,\Delta\,F.
\end{equation}

%-------------------------------------------------------------
\subsection*{V. Construction d'un vecteur gaussien de covariance $\mathsf{C}$}

On cherche maintenant un vecteur gaussien complexe
\[
  Y = (Y_0,\dots,Y_{M-1})^\top \in\C^M
\]
tel que
\begin{equation}
  \E[Y] = 0,
  \qquad
  \E\bigl[ Y\,Y^{\ast} \bigr] = \mathsf{C}.
\end{equation}

\subsubsection*{V.1. Idée générale}

Soit $Z\in\C^M$ un vecteur gaussien complexe dit \emph{standard},
c'est-à-dire
\begin{equation}
  \E[Z] = 0,
  \qquad
  \E[Z\,Z^{\ast}] = I_M
\end{equation}
($I_M$ matrice identité). Posons
\begin{equation}
  Y = F^{\ast}\,\Delta^{1/2}\,Z,
\end{equation}
où $\Delta^{1/2}=\mathrm{diag}(\sqrt{\lambda_0},\dots,
\sqrt{\lambda_{M-1}})$.

Alors
\begin{align}
  \E[Y]
  &= F^{\ast}\,\Delta^{1/2}\,\E[Z]
   = 0,
\end{align}
et
\begin{align}
  \E\bigl[Y\,Y^{\ast}\bigr]
  &= \E\bigl[
      F^{\ast}\Delta^{1/2}Z
      \, (F^{\ast}\Delta^{1/2}Z)^{\ast}
    \bigr] \\[0.3em]
  &= F^{\ast}\Delta^{1/2}
     \E\bigl[
       Z\,Z^{\ast}
     \bigr]
     \Delta^{1/2}F \\[0.3em]
  &= F^{\ast}\Delta^{1/2}
     I_M
     \Delta^{1/2}F \\[0.3em]
  &= F^{\ast}\Delta F
   = \mathsf{C}.
\end{align}
Ainsi, tout vecteur de la forme $Y=F^{\ast}\Delta^{1/2}Z$ a bien la
covariance souhaitée.

\subsubsection*{V.2. Construction explicite de $Z$ et réalité du signal}

Pour que le vecteur final que l'on utilisera soit \emph{réel}, on impose
à $Z$ une symétrie conjuguée particulière. On pose $M=2N$.

\begin{itemize}
  \item On génère $Z_0$ et $Z_N$ réels :
  \[
    Z_0 \sim \N(0,1),\qquad
    Z_N \sim \N(0,1).
  \]
  \item Pour $k=1,\dots,N-1$, on génère des gaussiennes réelles
        indépendantes
        \[
          U_k,V_k \sim \N(0,1),
        \]
        indépendantes entre elles et de $Z_0,Z_N$, puis on pose
        \[
          Z_k = \frac{1}{\sqrt{2}}(U_k + i V_k),
          \qquad
          Z_{M-k} = \overline{Z_k}
          = \frac{1}{\sqrt{2}}(U_k - i V_k).
        \]
\end{itemize}

Par construction, $Z$ est gaussien complexe centré, et un calcul
direct montre que
\[
  \E[Z] = 0,
  \qquad
  \E\bigl[Z\,Z^{\ast}\bigr] = I_M.
\]
En outre, comme $\Delta$ est réelle et diagonale, la combinaison
$\Delta^{1/2}Z$ vérifie la même symétrie conjuguée :
\[
  (\Delta^{1/2}Z)_{M-k}
  = \overline{(\Delta^{1/2}Z)_k},
  \qquad k=1,\dots,N-1,
\]
et $(\Delta^{1/2}Z)_0$, $(\Delta^{1/2}Z)_N$ sont réels.

Cette symétrie conjuguée, injectée dans l’IFFT (multiplication par
$F^{\ast}$), garantit que les composantes de $Y$ sont réelles en
arithmétique exacte:
\begin{equation}
  Y_j \in\R,\qquad j=0,\dots,M-1.
\end{equation}
En pratique numérique, de petits résidus imaginaires dus aux erreurs
d’arrondi peuvent subsister; on prend alors la partie réelle
$\Re(Y_j)$.

On résume donc :
\begin{equation}
  \label{eq:def-Y}
  Y = F^{\ast}\bigl(\Delta^{1/2} Z\bigr)
  = \mathrm{IFFT}\bigl(\sqrt{\bm{\lambda}}\;\odot Z\bigr),
\end{equation}
où $\odot$ désigne le produit terme à terme, et
$\sqrt{\bm{\lambda}}=(\sqrt{\lambda_0},\dots,\sqrt{\lambda_{M-1}})$.

%-------------------------------------------------------------
\subsection*{VI. Extraction du fractional Gaussian noise}

Le vecteur $Y$ est de taille $M=2N$. La théorie de Davies--Harte montre
que les $N$ \emph{premières} composantes de $Y$ forment un vecteur
gaussien réel de covariance $\Gamma_N$ (cf. encastrement Toeplitz
dans la circulante $\mathsf{C}$).

On définit donc le bruit fractionnaire \emph{normalisé} (pas unité)
par :
\begin{equation}
  \label{eq:def-fGn-unit}
  \tilde{X}_k = Y_k,
  \qquad k=0,\dots,N-1.
\end{equation}
Ces variables vérifient
\begin{equation}
  \Cov(\tilde{X}_j,\tilde{X}_{j+m})
  = \gamma_0(m),
  \qquad 0\le j, j+m \le N-1,
\end{equation}
c’est-à-dire qu’elles ont la covariance correspondant à un pas de
temps unité.

Pour retrouver la covariance au pas $\Delta t$, on applique le
facteur d’échelle $(\Delta t)^H$ :
\begin{equation}
  X_k = (\Delta t)^H \tilde{X}_k,
  \qquad k=0,\dots,N-1.
\end{equation}
On a alors
\begin{align}
  \Cov(X_j,X_{j+m})
  &= (\Delta t)^{2H}\Cov(\tilde{X}_j,\tilde{X}_{j+m}) \\
  &= (\Delta t)^{2H}\gamma_0(m)
   = \gamma(m),
\end{align}
ce qui est exactement la covariance théorique des incréments de fBm
au pas $\Delta t$.

En résumé, les $X_k$ construits ainsi forment un fractional
Gaussian noise au pas $\Delta t$.

%-------------------------------------------------------------
\subsection*{VII. Reconstruction détaillée du fBm discret}

À partir du bruit fractionnaire $(X_0,\dots,X_{N-1})$, on reconstruit
la trajectoire discrète de la fBm par cumul des incréments :
\begin{equation}
  \bH(t_0) = 0,
  \qquad
  \bH(t_k) = \sum_{j=0}^{k-1} X_j,
  \quad k=1,\dots,N.
\end{equation}
En insistant sur l’échelle en temps :
\begin{equation}
  \bH(t_k)
  = \sum_{j=0}^{k-1} X_j
  = (\Delta t)^H \sum_{j=0}^{k-1} \tilde{X}_j,
  \qquad t_k = k\Delta t.
\end{equation}
On obtient ainsi une trajectoire
$(\bH(t_0),\dots,\bH(t_N))$ dont les incréments ont la bonne
covariance et donc la bonne régularité (Hölder de paramètre $H$
quasi-sûrement).

%-------------------------------------------------------------
\subsection*{VIII. Résumé algorithmique complet (Davies--Harte)}

\begin{enumerate}
  \item \textbf{Discrétisation en temps.}
        Fixer $T>0$, $N\in\mathbb{N}^\ast$,
        $\Delta t = T/N$, $t_k = k\Delta t$.
  \item \textbf{Covariance théorique du bruit fractionnaire.}
        Calculer pour $m=0,\dots,N$ :
        \[
          \gamma_0(m)
          = \frac{1}{2}
            \Bigl(|m+1|^{2H} - 2|m|^{2H} + |m-1|^{2H}\Bigr).
        \]
  \item \textbf{Vecteur circulant.}
        Construire
        \[
          c = \bigl(\gamma_0(0),\gamma_0(1),\dots,
                    \gamma_0(N-1),\gamma_0(N),
                    \gamma_0(N-1),\dots,\gamma_0(1)\bigr)
          \in\R^{2N}.
        \]
  \item \textbf{Valeurs propres par FFT.}
        Calculer
        \[
          \bm{\lambda} = \mathrm{FFT}(c).
        \]
        Vérifier numériquement que $\lambda_k\ge 0$ pour tous $k$,
        puis poser $\Delta^{1/2}=\mathrm{diag}(\sqrt{\lambda_0},
        \dots,\sqrt{\lambda_{2N-1}})$.
  \item \textbf{Vecteur gaussien complexe symétrique.}
        Construire $Z\in\C^{2N}$ :
        \begin{itemize}
          \item $Z_0,Z_N\sim\N(0,1)$ réels,
          \item pour $k=1,\dots,N-1$,
                $Z_k = (U_k+iV_k)/\sqrt{2}$,
                $Z_{2N-k} = \overline{Z_k}$,
                avec $U_k,V_k\sim\N(0,1)$ indépendants.
        \end{itemize}
  \item \textbf{IFFT et bruit fractionnaire.}
        Calculer
        \[
          Y = \mathrm{IFFT}\bigl(\sqrt{\bm{\lambda}}\;\odot Z\bigr)
          = F^{\ast}\bigl(\Delta^{1/2}Z\bigr)
          \in\R^{2N}.
        \]
        Poser le bruit fractionnaire normalisé
        \[
          \tilde{X}_k = Y_k,\qquad k=0,\dots,N-1,
        \]
        puis le bruit fractionnaire au pas $\Delta t$ :
        \[
          X_k = (\Delta t)^H \tilde{X}_k,\qquad k=0,\dots,N-1.
        \]
  \item \textbf{Reconstruction de la fBm.}
        Définir
        \[
          \bH(t_0) = 0,
          \qquad
          \bH(t_k) = \sum_{j=0}^{k-1} X_j,
          \quad k=1,\dots,N.
        \]
\end{enumerate}

On obtient ainsi une trajectoire discrète simulée du mouvement
brownien fractionnaire sur $[0,T]$ par la méthode de
Davies--Harte.

%-------------------------------------------------------------
\subsection*{IX. Lien détaillé avec le modèle rough Bergomi}

On relie maintenant cette construction de fBm à la dynamique du
modèle \emph{rough Bergomi} (rBergomi).

%-------------------------------------------------------------
\subsubsection*{IX.1. Définition continue du modèle rough Bergomi}

Sous la mesure risque-neutre, le modèle rough Bergomi est défini
sur un espace probabilisé portant deux mouvements browniens
indépendants $W$ et $W^\perp$. On fixe un paramètre de corrélation
$\rho\in[-1,1]$ et on définit
\begin{equation}
  Z_t = \rho\,W_t + \sqrt{1-\rho^2}\,W_t^\perp.
\end{equation}
Le prix $S_t$ d’un actif sans dividende suit
\begin{equation}
  dS_t = S_t \sqrt{V_t}\,dZ_t,
\end{equation}
où $V_t$ est le processus de variance instantanée.

Le facteur de volatilité est donné par
\begin{equation}
  \label{eq:def-WH-Volterra}
  W_t^H = \sqrt{2H}\int_0^t (t-s)^{H-\frac12}\,dW_s,
  \qquad H\in(0,\tfrac12),
\end{equation}
et la variance est
\begin{equation}
  \label{eq:def-V-rBergomi}
  V_t = \xi_0(t)\,
        \exp\!\left(
          \eta W_t^H - \frac12 \eta^2 t^{2H}
        \right).
\end{equation}
Ici :
\begin{itemize}
  \item $\xi_0(t)$ est la courbe de variance \emph{forward} initiale
        (déterministe) ;
  \item $\eta>0$ est le paramètre de ``vol-of-vol'' ;
  \item $W^H$ est un processus gaussien de type fBm défini par
        l’intégrale de Volterra \eqref{eq:def-WH-Volterra}.
\end{itemize}

%-------------------------------------------------------------
\subsubsection*{IX.2. Vérification que $W^H$ est un mouvement brownien fractionnaire}

On vérifie que $W^H$ défini par \eqref{eq:def-WH-Volterra} a bien
la même covariance qu’un fBm de paramètre $H$.

\paragraph{Espérance.}
Par linéarité de l’intégrale d’Itô,
\begin{equation}
  \E[W_t^H]
  = \sqrt{2H}\,\E\Bigl[\int_0^t (t-s)^{H-\frac12}\,dW_s\Bigr]
  = 0.
\end{equation}

\paragraph{Variance.}
Pour $t\ge 0$,
\begin{align}
  \Var(W_t^H)
  &= 2H\,\E\left[\left(
        \int_0^t (t-s)^{H-\frac12} dW_s
      \right)^2\right] \\
  &= 2H \int_0^t (t-s)^{2H-1}\,ds
     \quad\text{(isométrie de l'intégrale stochastique)}  \\[0.3em]
  &= 2H \int_0^t u^{2H-1}\,du
     \quad (u = t-s) \\[0.3em]
  &= 2H \cdot \frac{t^{2H}}{2H}
   = t^{2H}.
\end{align}
On retrouve donc la variance d’un fBm de paramètre $H$.

\paragraph{Covariance.}
Pour $s,t\ge 0$,
\begin{align}
  \Cov(W_t^H,W_s^H)
  &= 2H\,\E\left[
      \int_0^t (t-u)^{H-\frac12} dW_u
      \int_0^s (s-v)^{H-\frac12} dW_v
    \right] \\[0.3em]
  &= 2H \int_0^{t\wedge s}
          (t-u)^{H-\frac12}(s-u)^{H-\frac12}\,du
     \quad\text{(isométrie)}.
\end{align}
Un calcul classique (changement de variable $u=s-x$, puis utilisation
de la fonction Bêta) donne
\begin{equation}
  2H \int_0^{t\wedge s}
          (t-u)^{H-\frac12}(s-u)^{H-\frac12}\,du
  = \frac{1}{2}\Bigl(t^{2H}+s^{2H}-|t-s|^{2H}\Bigr)
  = R_H(t,s).
\end{equation}
Ainsi,
\begin{equation}
  \Cov(W_t^H,W_s^H) = R_H(t,s),
\end{equation}
exactement la covariance d’un mouvement brownien fractionnaire de
paramètre $H$. Par unicité de la loi gaussienne à partir de la
covariance, on a
\begin{equation}
  (W_{t_1}^H,\dots,W_{t_n}^H)
  \stackrel{d}{=}
  (\bH(t_1),\dots,\bH(t_n))
\end{equation}
pour tout $n$ et tout choix de temps $t_1,\dots,t_n$. Le processus
$W^H$ de rough Bergomi est donc \emph{en loi} un fBm standard de
paramètre $H$.

Conséquence importante : sur une grille
$t_k=k\Delta t$, on peut simuler directement $(\bH(t_k))$ par
Davies--Harte et l’utiliser comme échantillonnage de
$(W_{t_k}^H)$, puisque les lois finies-dimensionnelles coïncident.

%-------------------------------------------------------------
\subsubsection*{IX.3. Log-volatilité et espérance conditionnelle}

La log-volatilité définie par \eqref{eq:def-V-rBergomi} s’écrit
\begin{equation}
  \log V_t
  = \log \xi_0(t)
    + \eta W_t^H - \frac12 \eta^2 t^{2H}.
\end{equation}
On a déjà vu que $W_t^H\sim\N(0,t^{2H})$. Donc
\begin{equation}
  \eta W_t^H \sim \N\bigl(0,\eta^2 t^{2H}\bigr).
\end{equation}
Le moment exponentiel de $W_t^H$ est donné par
\begin{equation}
  \E\bigl[e^{\eta W_t^H}\bigr]
  = \exp\!\left(\frac12 \eta^2 \Var(W_t^H)\right)
  = \exp\!\left(\frac12 \eta^2 t^{2H}\right).
\end{equation}
On calcule alors l'espérance de $V_t$ :
\begin{align}
  \E[V_t]
  &= \E\left[
        \xi_0(t)
        \exp\!\left(
          \eta W_t^H - \frac12 \eta^2 t^{2H}
        \right)
      \right] \\
  &= \xi_0(t)\,
     e^{-\frac12\eta^2 t^{2H}}\,
     \E\bigl[e^{\eta W_t^H}\bigr] \\
  &= \xi_0(t)\,
     e^{-\frac12\eta^2 t^{2H}}\,
     e^{\frac12\eta^2 t^{2H}} \\
  &= \xi_0(t).
\end{align}
Par construction, $\xi_0(t)$ est donc bien la \emph{variance
forward} initiale :
\begin{equation}
  \xi_0(t) = \E[V_t],\qquad t\ge 0.
\end{equation}

%-------------------------------------------------------------
\subsubsection*{IX.4. Discrétisation de rough Bergomi à partir de la fBm simulée}

On se donne maintenant la grille $t_k=k\Delta t$,
$k=0,\dots,N$, et une trajectoire simulée de fBm
\[
  \bigl(\bH(t_0),\dots,\bH(t_N)\bigr)
\]
obtenue par la méthode de Davies--Harte décrite aux sections précédentes.

\paragraph{1. Identification $W^H\equiv\bH$.}
Comme démontré ci-dessus,
\[
  (W_{t_k}^H)_{k=0,\dots,N}
  \stackrel{d}{=}
  (\bH(t_k))_{k=0,\dots,N}.
\]
Pour la simulation, on pose donc simplement
\begin{equation}
  W_{t_k}^H := \bH(t_k),\qquad k=0,\dots,N.
\end{equation}

\paragraph{2. Discrétisation de la variance.}
On obtient la variance rough Bergomi sur la grille en posant
\begin{equation}
  V_k := V_{t_k}
  = \xi_0(t_k)\,
    \exp\!\left(
      \eta \bH(t_k) - \frac12 \eta^2 t_k^{2H}
    \right),
  \qquad k=0,\dots,N.
\end{equation}
Cette suite $(V_k)$ est log-normale, et la relation
$\E[V_k]=\xi_0(t_k)$ reste vérifiée.

\paragraph{3. Discrétisation de la dynamique du prix.}
L’équation
\[
  dS_t = S_t \sqrt{V_t}\,dZ_t
\]
est un SDE de type Black–Scholes avec volatilité stochastique
$\sqrt{V_t}$. Par Itô,
\begin{equation}
  d\log S_t
  = \frac{1}{S_t}dS_t - \frac{1}{2S_t^2}(dS_t)^2
  = \sqrt{V_t}\,dZ_t - \frac12 V_t\,dt.
\end{equation}
En supposant $V_t$ constant sur $[t_k,t_{k+1})$ et égal à $V_k$,
on obtient le schéma de type Euler exponentiel :
\begin{equation}
  \log S_{t_{k+1}}
  = \log S_{t_k}
    - \frac12 V_k \Delta t
    + \sqrt{V_k}\,\Delta Z_k,
\end{equation}
où
\[
  \Delta Z_k = Z_{t_{k+1}} - Z_{t_k} \sim \N(0,\Delta t).
\]
En exponentiant,
\begin{equation}
  S_{t_{k+1}}
  = S_{t_k}\,
    \exp\!\left(
      -\frac12 V_k \Delta t
      + \sqrt{V_k}\,\Delta Z_k
    \right).
\end{equation}

\paragraph{4. Corrélation avec le facteur rough.}
Dans la définition continue, $Z$ est construit à partir de $W$ via
\[
  Z_t = \rho W_t + \sqrt{1-\rho^2}\,W_t^\perp,
\]
et $W^H$ est une intégrale stochastique de $W$ (voir
\eqref{eq:def-WH-Volterra}). Ceci implique une structure de
corrélation non triviale entre $W^H$ et $Z$.

Si l’on se contente de simuler la fBm $(\bH(t_k))$ \emph{indépendamment}
d’un Brownien discrétisé $(Z_{t_k})$, alors le prix $S$ est
indépendant du facteur rough : ce n’est plus le modèle rBergomi mais
un modèle à volatilité rough décorrélée. Pour avoir le rBergomi
correct, il faut construire conjointement un vecteur gaussien
\[
  \bigl(W_{t_1}^H,\dots,W_{t_N}^H,
        Z_{t_1},\dots,Z_{t_N}\bigr)
\]
avec la bonne matrice de covariance. Celle-ci est de la forme
\[
  \Sigma =
  \begin{pmatrix}
    \Sigma^{HH} & \Sigma^{HZ} \\
    (\Sigma^{HZ})^\top & \Sigma^{ZZ}
  \end{pmatrix},
\]
où :
\begin{itemize}
  \item $\Sigma^{HH}_{ij}
         = \Cov(W_{t_i}^H,W_{t_j}^H)
         = R_H(t_i,t_j)$ ;
  \item $\Sigma^{ZZ}_{ij}
         = \Cov(Z_{t_i},Z_{t_j})
         = \min(t_i,t_j)$ (Brownien standard) ;
  \item $\Sigma^{HZ}_{ij}
         = \Cov(W_{t_i}^H,Z_{t_j})$.
\end{itemize}

Pour ce dernier terme, en utilisant \eqref{eq:def-WH-Volterra} et
$Z_t = \rho W_t +\sqrt{1-\rho^2}W_t^\perp$, on obtient pour
$t_i\ge t_j$ :
\begin{align}
  \Cov(W_{t_i}^H,Z_{t_j})
  &= \rho\,\Cov\!\left(
        \sqrt{2H}\int_0^{t_i}(t_i-s)^{H-\frac12}dW_s,\;
        W_{t_j}
      \right) \\
  &= \rho\sqrt{2H}\int_0^{t_j}(t_i-s)^{H-\frac12}ds \\[0.3em]
  &= \rho\sqrt{2H}
     \int_{t_i-t_j}^{t_i} u^{H-\frac12}du
     \quad (u=t_i-s) \\[0.3em]
  &= \rho\sqrt{2H}\,
     \frac{t_i^{H+\frac12}-(t_i-t_j)^{H+\frac12}}{H+\frac12}.
\end{align}
Pour $t_i<t_j$, on échange les rôles de $t_i$ et $t_j$.

En pratique, pour une implémentation exacte de rBergomi, on doit
simuler le vecteur gaussien de dimension $2N$ avec cette covariance
$\Sigma$ (par exemple via une factorisation de type Cholesky ou une
méthode FFT multivariée). La méthode Davies--Harte décrite plus haut
donne efficacement le bloc $\Sigma^{HH}$ ; la construction complète
est plus lourde mais repose sur le même principe : tout est
gaussien, donc entièrement déterminé par la covariance.

%-------------------------------------------------------------
\subsubsection*{IX.5. Résumé : utilisation pratique de la fBm simulée pour rBergomi}

Pour relier concrètement la simulation de fBm par Davies--Harte au
modèle rough Bergomi sur une grille $t_k=k\Delta t$ :

\begin{enumerate}
  \item Simuler une trajectoire de fBm $(\bH(t_k))_{k=0,\dots,N}$ via
        Davies--Harte.
  \item L’identifier en loi au facteur rough
        $W_{t_k}^H$, i.e. poser $W_{t_k}^H=\bH(t_k)$.
  \item Construire la variance
        \[
          V_k = \xi_0(t_k)
                \exp\!\left(
                  \eta \bH(t_k) - \frac12 \eta^2 t_k^{2H}
                \right).
        \]
  \item Simuler des incréments browniens corrélés $\Delta Z_k$ de sorte
        que la corrélation entre $W^H$ et $Z$ soit bien celle imposée
        par le modèle (via la covariance $\Sigma^{HZ}$ ci-dessus) ;
        dans une implémentation approximative plus simple, on peut
        prendre $\Delta Z_k$ simplement corrélés avec les incréments
        de $W$ ou considérer l’approximation de corrélation locale.
  \item Mettre à jour le prix par
        \[
          S_{t_{k+1}}
          = S_{t_k}
            \exp\!\left(
              -\frac12 V_k \Delta t
              + \sqrt{V_k}\,\Delta Z_k
            \right).
        \]
\end{enumerate}

La dérivation montre explicitement où intervient la fBm dans
rough Bergomi et comment la trajectoire obtenue par Davies--Harte
se branche dans la construction de la variance et du prix.



\fi
\subsection*{Annexe D - Mod\`ele SABR (expansion asymptotique)}
\addcontentsline{toc}{section}{Annexe D - Mod\`ele SABR}



\subsection*{Annexe D — Mod\`ele SABR (expansion asymptotique)}

%%%%%%%%%%%%%%%%%%%%%%%%%%%%%%%%%%%%%%%%%%%%%%%%%%%%%%%%%%%%%%
\subsection*{D.1. Dynamique}

\[
\begin{cases}
dF_t = \alpha_t F_t^{\beta} \, dW_t^{(1)}, \\[4pt]
d\alpha_t = \nu \alpha_t \, dW_t^{(2)}, \\[4pt]
dW_t^{(1)}dW_t^{(2)} = \rho dt.
\end{cases}
\]

%%%%%%%%%%%%%%%%%%%%%%%%%%%%%%%%%%%%%%%%%%%%%%%%%%%%%%%%%%%%%%
\subsection*{D.2. Formule asymptotique de la vol implicite}

La volatilité implicite obtenue par expansion en $\nu$ (Hagan et al.) est :
\[
\sigma_{\mathrm{SABR}}(K)
=
\frac{\alpha_0}{(F_0 K)^{\frac{1-\beta}{2}}}
\left[
1 + \frac{(1-\beta)^2}{24}(\ln(F_0/K))^2 
+ \frac{(1-\beta)^4}{1920}(\ln(F_0/K))^4
\right]
\Bigg/
z,
\]

avec
\[
z = \frac{\nu}{\alpha_0} (F_0 K)^{\frac{1-\beta}{2}} \ln\frac{F_0}{K},
\qquad
\chi(z) = \ln\!\left( \frac{\sqrt{1-2\rho z + z^2} + z - \rho}{1-\rho} \right).
\]

La forme finale couramment utilisée est :
\[
\sigma_{\mathrm{SABR}}(K)
=
\frac{\alpha_0}{(F_0 K)^{\frac{1-\beta}{2}}}
\left[
1 
+
\left(
\frac{(1-\beta)^2}{24}\ln^2\frac{F_0}{K}
+
\frac{(1-\beta)^4}{1920}\ln^4\frac{F_0}{K}
\right)
\right]
\frac{z}{\chi(z)}.
\]

%%%%%%%%%%%%%%%%%%%%%%%%%%%%%%%%%%%%%%%%%%%%%%%%%%%%%%%%%%%%%%
\subsection*{D.3. Utilisation en calibration}

La structure $(\beta,\nu,\rho,\alpha_0)$ permet de calibrer finement le smile :

\[
\theta^\*
=
\arg\min_{\theta}
\sum_{K\in\mathcal{K}}
\left(
\sigma_{\mathrm{SABR}}(K;\theta)
-
\sigma_{\mathrm{impl}}^{\mathrm{mkt}}(K)
\right)^2.
\]

%%%%%%%%%%%%%%%%%%%%%%%%%%%%%%%%%%%%%%%%%%%%%%%%%%%%%%%%%%%%%%
\subsection*{D.4. Remarques}

\begin{itemize}
  \item L’expansion est très précise pour $\nu$ modéré et $|F_0-K|$ raisonnable.
  \item Des variantes existent (SABR normal, shifted-SABR).
  \item La dérivation complète nécessite des développements asymptotiques 
        détaillés (reportés en annexe technique).
\end{itemize}

%%%%%%%%%%%%%%%%%%%%%%%%%%%%%%%%%%%%%%%%%%%%%%%%%%%%%%%%%%%%%%



\subsection*{Annexe E - Mod\`ele de Rough Heston}
\addcontentsline{toc}{section}{Annexe E - Rough Heston}



\subsection*{Annexe E — Mod\`ele de Rough Heston (dynamique, caract\'eristique, simulation)}

%%%%%%%%%%%%%%%%%%%%%%%%%%%%%%%%%%%%%%%%%%%%%%%%%%%%%%%%%%%%%%
\subsection*{E.1. Motivation}

Le mod\`ele de Heston classique g\'enère une variance $v_t$ markovienne de type
CIR.  
Cependant, les observations empiriques montrent un comportement plus ``rugbueux''
de la variance, avec une r\'egularité H\"older strictement inférieure à $1/2$ et
une d\'ecorrélation lente.

Le \emph{Rough Heston} introduit cette rugosit\'e en remplaçant la dynamique CIR
par une SDE de Volterra reposant sur un noyau de type $t^{H-\frac12}$ avec
$H<\frac12$.

%%%%%%%%%%%%%%%%%%%%%%%%%%%%%%%%%%%%%%%%%%%%%%%%%%%%%%%%%%%%%%
\subsection*{E.2. Dynamique du Rough Heston}

On considère le modèle suivant (sous la mesure risque--neutre) :
\[
\begin{cases}
dS_t = S_t \sqrt{v_t}\, dW_t^{(1)} + r S_t\, dt, \\[6pt]
v_t = v_0 + \frac{1}{\Gamma(H+\tfrac12)} 
      \int_0^t (t-s)^{H-\frac12}
      \bigl( \kappa(\theta - v_s)\, ds + \xi \sqrt{v_s}\, dW_s^{(2)} \bigr),
\end{cases}
\]
avec covariance :
\[
dW_t^{(1)} dW_t^{(2)} = \rho \, dt.
\]

Le noyau :
\[
K(t) = \frac{ t^{H-\frac12} }{\Gamma(H+\frac12)}, \qquad H<\frac12,
\]
introduit une m\'emoire longue et une r\'egularit\'e faible.

%%%%%%%%%%%%%%%%%%%%%%%%%%%%%%%%%%%%%%%%%%%%%%%%%%%%%%%%%%%%%%
\subsection*{E.3. Formulation Volterra}

La variance satisfait la Volterra SDE affine :
\[
v_t = v_0 + \int_0^t K(t-s)\, \kappa(\theta - v_s)\, ds
      + \int_0^t K(t-s)\, \xi \sqrt{v_s}\, dW_s^{(2)}.
\]

La structure affine demeure, ce qui permet des formules fermées pour la fonction
caractéristique du log--prix.

%%%%%%%%%%%%%%%%%%%%%%%%%%%%%%%%%%%%%%%%%%%%%%%%%%%%%%%%%%%%%%
\subsection*{E.4. Fonction caract\'eristique du log--prix}

On pose $X_t = \ln S_t$.  
On cherche la fonction caractéristique :
\[
\phi(u,t) 
=
\mathbb{E}\bigl[ e^{iu X_t} \mid X_0,v_0 \bigr].
\]

Comme le modèle est affine Volterra, on a :
\[
\phi(u,t)
=
\exp\!\bigl( iu X_0 + C(u,t) + D(u,t)v_0 \bigr).
\]

Les fonctions $C(u,t)$ et $D(u,t)$ satisfont des \emph{équations de Riccati fractionnaires}
\[
D(u,t) = \int_0^t K(t-s)
\biggl( 
\frac12 (u^2+iu)
- (\kappa - \rho\xi iu) D(u,s)
- \frac12\xi^2 D(u,s)^2
\biggr)\, ds,
\]
\[
C(u,t) = iu r t + \kappa\theta \int_0^t D(u,s)\, ds.
\]

Il s'agit de convolutions Volterra nécessitant des méthodes numériques dédiées 
(FFT accélérée, quadrature rapide, schémas hybrides).

%%%%%%%%%%%%%%%%%%%%%%%%%%%%%%%%%%%%%%%%%%%%%%%%%%%%%%%%%%%%%%
\subsection*{E.5. Prix du call europ\'een}

Le prix du call est obtenu via la représentation de Carr–Madan / Fourier :
\[
C_0 = S_0 P_1 - K e^{-rT} P_2,
\]
où
\[
P_j 
= 
\frac12 + \frac1\pi 
\int_0^\infty
\Re\left(
\frac{ e^{-iu\ln K}\, \phi_j(u,T) }{iu}
\right) du,
\qquad j=1,2.
\]

La différence avec Heston classique provient du calcul de $\phi_j(u,T)$, qui 
utilise les Riccati fractionnaires ci-dessus.

%%%%%%%%%%%%%%%%%%%%%%%%%%%%%%%%%%%%%%%%%%%%%%%%%%%%%%%%%%%%%%
\subsection*{E.6. Simulation : schéma hybride (hybrid scheme)}

La simulation directe du rough Heston est coûteuse car la convolution Volterra
est de complexité naïve $\mathcal{O}(n^2)$.  
Bayer, Friz et Gatheral (2016) proposent une approche \emph{hybride} qui combine
une d\'ecomposition du noyau (partie singuli\`ere vs r\'eguli\`ere) et une acc\'el\'eration des convolutions,
afin de r\'eduire le co\^ut de calcul sur une grille. Le point essentiel (au niveau th\'eorique) est la gestion
de la non-markovianit\'e : l'\'etat d\'epend d'un historique, ce qui impose des sch\'emas sp\'ecialis\'es pour obtenir
des trajectoires utilisables en pricing/greeks.

%%%%%%%%%%%%%%%%%%%%%%%%%%%%%%%%%%%%%%%%%%%%%%%%%%%%%%%%%%%%%%
\subsection*{E.7. Calibration}

La calibration consiste à résoudre :
\[
\theta^\*
=
\arg\min_{\theta}
\sum_{(K,T)}
\bigl(
\sigma_{\mathrm{mod}}(K,T;\theta)
-
\sigma_{\mathrm{mkt}}(K,T)
\bigr)^2,
\]
où $\theta = (H,\eta,\kappa,\theta,\xi,\rho,v_0)$.

Le paramètre clé est le Hurst $H$ :

\begin{itemize}
  \item $H<0.5$ est la signature ``rough'' : il accentue la g\'eom\'etrie court terme du smile (motivation empirique).
  \item La structure Volterra augmente l'expressivit\'e par rapport \`a Heston classique, au prix d'une calibration
        et d'une validation plus d\'elicates (instabilit\'e, sur-ajustement possible).
  \item Sans donn\'ees fournies ici, une comparaison ``Rough Heston vs Heston'' doit \^etre pos\'ee comme un protocole :
        calibration in-sample sur une date, puis test out-of-sample sur d'autres dates/maturit\'es, avec contr\^oles
        de stabilit\'e des param\`etres et de violation d'arbitrage statique.
\end{itemize}

%%%%%%%%%%%%%%%%%%%%%%%%%%%%%%%%%%%%%%%%%%%%%%%%%%%%%%%%%%%%%%
\subsection*{E.8. Synthèse}

Le Rough Heston combine :
\begin{itemize}
  \item structure affine (donc transformées de Fourier disponibles),
  \item mémoire longue via Volterra,
  \item réalisme empirique, notamment sur le skew court terme.
\end{itemize}

Il constitue un cadre avanc\'e pour le pricing d'options equity et la mod\'elisation de surfaces de volatilit\'e.

Plus g\'en\'eralement, c'est un cadre avanc\'e (structure affine + m\'emoire) dont la pertinence pratique doit \^etre
\'evalu\'ee par des diagnostics de calibration et de couverture, en tenant compte des co\^uts et du risque de mod\`ele.

%%%%%%%%%%%%%%%%%%%%%%%%%%%%%%%%%%%%%%%%%%%%%%%%%%%%%%%%%%%%%%
\section*{Annexe F — M\'ethodes num\'eriques (FFT, Crank--Nicolson)}
\addcontentsline{toc}{section}{Annexe F — M\'ethodes num\'eriques}

\subsection*{F.1. M\'ethode de Crank--Nicolson pour la PDE de Black--Scholes}

On consid\`ere la PDE de Black--Scholes :
\[
\frac{\partial V}{\partial t} + \frac12 \sigma^2 S^2 \frac{\partial^2 V}{\partial S^2}
+ r S \frac{\partial V}{\partial S} - r V = 0.
\]

En changeant de variable $x = \ln S$ et $\tau = T-t$, on obtient une PDE de type
chaleur pour $\tilde V(\tau,x)$, que l’on discrétise sur une grille $(\tau_n, x_j)$.

Le schéma de Crank--Nicolson consiste à moyenniser les schémas implicite et explicite :
\[
\frac{V^{n+1}_j - V^n_j}{\Delta \tau}
=
\frac12 \bigl( \mathcal{L}V^{n+1}_j + \mathcal{L}V^{n}_j \bigr),
\]
o\`u $\mathcal{L}$ est l’opérateur spatial issu de la PDE.

En pratique, cela donne un syst\`eme tridiagonal à chaque pas de temps :
\[
A V^{n+1} = B V^n + b^n,
\]
que l’on r\'esout via un algorithme de Thomas (co\^ut en $\mathcal{O}(M)$ pour 
$M$ points en espace).

\subsection*{F.2. FFT pour le pricing d’options (Carr--Madan)}

L’id\'ee est d’exprimer le prix d’un call en termes de transformée de Fourier
de la fonction caract\'eristique du log--prix.

Soit
\[
\phi_T(u) = \mathbb{E}\bigl[ e^{iu \log S_T} \bigr].
\]

On introduit un call dampt\'e :
\[
c_\alpha(k) = e^{\alpha k} C(k),
\qquad
k = \ln K,
\]
avec $\alpha > 0$.

La transformée de Fourier de $c_\alpha(k)$ s’exprime alors en fonction de 
$\phi_T(u)$ :
\[
\hat c_\alpha(u)
=
\int_{-\infty}^{+\infty} e^{i u k} c_\alpha(k)\, dk
=
\frac{e^{-rT} \phi_T(u - i(\alpha+1))}{\alpha^2 + \alpha - u^2 + iu(2\alpha +1)}.
\]

Les prix pour un ensemble de strikes équidistants $k_m$ sont obtenus par FFT :
\[
c_\alpha(k_m)
\approx 
\frac{e^{-i u_0 k_m}}{2\pi}
\sum_{n=0}^{N-1} e^{i u_n k_m} \hat c_\alpha(u_n) \Delta u,
\]
o\`u $(u_n)$ est une grille uniforme de fréquences.

Enfin, on récupère le prix du call :
\[
C(k_m) = e^{-\alpha k_m} c_\alpha(k_m).
\]

Cette m\'ethode est particulièrement efficace pour les mod\`eles à fonction 
caract\'eristique explicite (Heston, Merton, Variance Gamma, etc.).

\subsection*{F.3. Résumé}

\begin{itemize}
  \item Le schéma de Crank--Nicolson fournit une solution stable et seconde-ordre
        en temps pour les PDE de pricing.
  \item Les méthodes FFT de type Carr--Madan permettent d’obtenir de nombreux prix
        d’options en un seul appel de Fourier, particulièrement bien adaptées aux
        mod\`eles de type affine.
\end{itemize}

\section*{Annexe G --- Projet annexe : pipeline de pr\'ediction tabulaire structur\'e par eras (cas Numerai)}
\addcontentsline{toc}{section}{Annexe G --- Projet annexe : pipeline Numerai}

\subsection*{Objectif et positionnement}
Cette annexe pr\'esente, au niveau conceptuel (sans code), un \emph{pipeline} de pr\'ediction tabulaire structur\'e par p\'eriodes (``eras'') tel
qu'on en rencontre dans certains tournois de pr\'evision (cas Numerai). L'int\'er\^et est m\'ethodologique : (i) g\'erer une non-stationnarit\'e
explicite, (ii) \'eviter le \emph{leakage} temporel, et (iii) privil\'egier la stabilit\'e du signal plut\^ot qu'une performance moyenne illusoire.

\subsection*{Donn\'ees : entra\^{i}nement vs pr\'ediction}
On dispose conceptuellement de deux ensembles :
\begin{itemize}
  \item un jeu d'\textbf{entra\^{i}nement} contenant des variables explicatives $x$ et une cible $y$ ;
  \item un jeu \textbf{``live''} contenant les variables $x$ \`a pr\'edire (sans $y$).
\end{itemize}
Les observations sont regroup\'ees en eras $e=1,\dots,E$, permettant une \'evaluation et une validation \emph{par blocs temporels}.

\subsection*{Cha\^{i}ne end-to-end (vue macro)}
Le pipeline peut \^etre r\'esum\'e par les \'etapes suivantes :
\begin{enumerate}[leftmargin=*, itemsep=3pt]
  \item \textbf{Contr\^oles de coh\'erence} : variables disponibles, valeurs manquantes, d\'erive grossi\`ere de distribution.
  \item \textbf{R\'eduction de dimension} : filtre de valeurs manquantes puis s\'election (par score) des variables les plus informatives.
  \item \textbf{D\'ecoupage temporel} : validation par blocs d'eras (rolling/expanding) et r\'eservation d'un holdout final.
  \item \textbf{Entra\^{i}nement} : ajustement d'un mod\`ele (ex.\ GBDT) sur chaque split, avec r\'egularisation et arr\^et anticip\'e.
  \item \textbf{Pr\'edictions hors-fold (OOF)} : construction de pr\'edictions sur les parties non vues pour estimer la performance hors \'echantillon.
  \item \textbf{\'Evaluation} : calcul d'une m\'etrique par era (ex.\ corr\'elation), puis agr\'egation (moyenne, dispersion, ratio type Sharpe).
  \item \textbf{Ensembles} : agr\'egation de plusieurs mod\`eles (moyenne) et, si justifi\'e, combinaison via meta-mod\`ele (stacking) sous validation stricte.
  \item \textbf{Post-traitement} : transformation monotone (ranking) pour stabiliser l'\'echelle et limiter l'impact d'outliers.
  \item \textbf{Production} : g\'en\'eration des pr\'edictions sur le jeu live dans le format standard attendu par la plateforme.
\end{enumerate}

\subsection*{Choix m\'ethodologiques : justification et limites}
\begin{itemize}
  \item \textbf{S\'election univari\'ee (score)} : efficace pour ma\^{i}triser le bruit et la dimension, mais ignore interactions ; elle doit \^etre compatible avec la validation temporelle (sinon leakage).
  \item \textbf{M\'etrique par era} : la stabilit\'e inter-eras est centrale en non-stationnarit\'e ; une moyenne seule peut masquer des inversions de signal.
  \item \textbf{Ensembles} : r\'eduisent la variance et l'instabilit\'e, au prix d'un co\^ut de calcul accru.
  \item \textbf{Ranking} : stabilise la distribution des scores ; il ne cr\'ee pas de signal et doit \^etre valid\'e (effet possible sur corr\'elation lin\'eaire).
\end{itemize}

\subsection*{Lien avec le m\'emoire}
Ce projet annexe illustre les principes d\'evelopp\'es au chapitre sur l'apprentissage supervis\'e sous non-stationnarit\'e
(\emph{validation temporelle, leakage, m\'etriques par p\'eriode, ensembles et post-traitements monotones}). Dans le contexte de
\emph{PaperTradingApp}, ces principes cadrent la validation de briques apprenantes (surrogates, r\'egimes, diagnostics) et compl\`etent les exigences
de tra\c{c}abilit\'e et de gouvernance des hypoth\`eses discut\'ees dans la partie m\'ethodologie.

\end{document}
